% આ ભાગનો પૂર્ણ સંપાદનીય પાઠ છે. શૈલી, ફોન્ટ, ચિત્રો અને પ્રકરણસંદર્ભ માટે સંપૂર્ણ સ્રોત ZIP તથા reader/full-reader.tex જુઓ. આ એકલી સ્વતંત્ર કમ્પાઇલ થતી મુખ્ય ફાઇલ નથી.

% BEGIN OLP-0252
\olpart{tur}{ટ્યુરિંગ મશીનો}
\phantomsection\label{guunit:OLP-0252}


\begingroup
\makeatletter\def\ol@sectioncs{section}\makeatother

% BEGIN OLP-0253
\olchapter{tur}{mac}{ટ્યુરિંગ મશીનની સંગણનાઓ}
\olchapterid{tur}{mac}
\phantomsection\label{guunit:OLP-0253}


\begingroup
\makeatletter\def\ol@sectioncs{section}\makeatother

% BEGIN OLP-0254
\olfileid{tur}{mac}{int}
\olsection{પરિચય}
\phantomsection\label{guunit:OLP-0254}


ઉદાહરણ તરીકે $\Nat$થી $\Nat$માં જતું કોઈ વિધેય \emph{સંગણનીય}
હોવાનો અર્થ શું? શરૂઆતના જવાબોમાંનો એક, અને સૌથી જાણીતો જવાબ,
એ છે કે કોઈ વિધેય ટ્યુરિંગ મશીન વડે ગણી શકાય તો તે સંગણનીય છે.
એલન ટ્યુરિંગે 1936માં આ કલ્પના રજૂ કરી. ટ્યુરિંગ મશીન
\emph{સંગણનના નિદર્શ}નું એક ઉદાહરણ છે---એ ``સંગણન પ્રક્રિયા''નો
વિચાર ગણિતીય રીતે ચોક્કસ વ્યાખ્યાયિત કરવાની રીત છે. આ વિચારનો
ચોક્કસ અર્થ શું છે તે ચર્ચાનો વિષય છે, પરંતુ ટ્યુરિંગ મશીનો
સંગણન પ્રક્રિયા નિર્દિષ્ટ કરવાની એક રીત છે એ બાબતે વ્યાપક સહમતિ
છે. ``ટ્યુરિંગ મશીન'' શબ્દથી ચાલતા ભાગો ધરાવતા ભૌતિક મશીનની
છબી ઊભી થાય છે, પરંતુ કડક અર્થમાં ટ્યુરિંગ મશીન શુદ્ધ ગણિતીય
રચના છે; તેથી તે સંગણન પ્રક્રિયાના વિચારનું આદર્શીકરણ કરે છે.
ઉદાહરણ તરીકે, આપણે ટ્યુરિંગ મશીનને લાગતા સમય કે સ્મૃતિની
જરૂરિયાત પર કોઈ મર્યાદા મૂકતા નથી: સંગણના માટે બ્રહ્માંડના
પરમાણુઓની સંખ્યા કરતાં વધુ સંગ્રહસ્થાન અથવા વધુ પગલાં જોઈએ
તોપણ ટ્યુરિંગ મશીન તે ગણતરી કરી શકે છે.

\begin{explain}
ટ્યુરિંગ મશીનને કોઈ વિશિષ્ટ પ્રકારના કલ્પિત ઉપકરણના પ્રોગ્રામ
તરીકે સમજવું કદાચ સૌથી યોગ્ય છે. આ ઉપકરણમાં એક \emph{ટેપ} અને
એક \emph{વાંચન-લેખન હેડ} હોય છે. ટ્યુરિંગ મશીનની આપણી આવૃત્તિમાં
ટેપ એક દિશામાં (જમણી તરફ) અનંત હોય છે અને \emph{ખાનાં}માં
વહેંચાયેલી હોય છે. દરેક ખાનામાં મર્યાદિત \emph{વર્ણમાળા}માંથી
એક પ્રતીક હોઈ શકે છે. આવી વર્ણમાળામાં સાન્ત સંખ્યાનાં ગમે તેટલાં જુદાં પ્રતીકો
હોઈ શકે, પરંતુ આપણે મુખ્યત્વે ત્રણથી કામ ચલાવીશું:
$\TMendtape$, $\TMblank$ અને $\TMstroke$. ઉપકરણ શરૂ થાય ત્યારે
ટેપ ખાલી હોય છે (એટલે દરેક ખાનામાં $\TMblank$ હોય છે), સિવાય કે
સૌથી ડાબા ખાનામાં $\TMendtape$ અને સાન્ત સંખ્યામાં અન્ય ખાનાંમાં
\emph{ઇનપુટ} હોય છે. કોઈ પણ સમયે ઉપકરણ સાન્ત સંખ્યાની
\emph{અવસ્થાઓ}માંથી એક અવસ્થામાં હોય છે. શરૂઆતમાં હેડ સૌથી
ડાબું ખાનું વાંચે છે અને ઉપકરણ નિર્દિષ્ટ \emph{પ્રારંભિક અવસ્થા}માં
હોય છે. ઉપકરણના ચાલવાના દરેક પગલે, હાલમાં હેડ જે ખાનું વાંચે
છે તેની સામગ્રી અને ઉપકરણની અવસ્થા મળીને, ટ્યુરિંગ મશીનના
પ્રોગ્રામ મુજબ આગળ શું થશે તે નક્કી કરે છે. ટ્યુરિંગ મશીનનો
પ્રોગ્રામ એક આંશિક વિધેય વડે અપાય છે: તે અવસ્થા~$q$ અને
પ્રતીક~$\sigma$ ઇનપુટ તરીકે લે છે અને ત્રિક
$\tuple{q', \sigma', D}$ આપે છે. ઉપકરણ અવસ્થા~$q$માં હોય અને
પ્રતીક~$\sigma$ વાંચે ત્યારે તે હાલના ખાનામાંનું પ્રતીક
$\sigma'$થી બદલે છે; $D$ અનુક્રમે $\TMleft$, $\TMright$ કે
$\TMstay$ હોય તે મુજબ હેડ ડાબે, જમણે કે એ જ સ્થાને જાય છે;
અને ઉપકરણ અવસ્થા~$q'$માં પ્રવેશે છે.

ઉદાહરણ તરીકે, \olref{fig:tm}માં દર્શાવેલી સ્થિતિ જુઓ.
\begin{figure}
  \olasset{\olpath/assets/diagrams/turing-machine.tikz}
  \caption{પોતાનો પ્રોગ્રામ ચલાવતું ટ્યુરિંગ મશીન.}
  \ollabel{fig:tm}
\end{figure}
ટ્યુરિંગ મશીનની ટેપનો દેખાતો ભાગ સૌથી ડાબા ખાનામાં ટેપના
છેડાનું પ્રતીક~$\TMendtape$ ધરાવે છે. ત્યાર પછી ત્રણ $1$,
એક $0$ અને બીજા ચાર $1$ છે. હેડ ડાબેથી ત્રીજું ખાનું વાંચે
છે, જેમાં $1$ છે, અને મશીન અવસ્થા~$q_1$માં છે---અપણે કહીએ
કે ``મશીન અવસ્થા~$q_1$માં $1$ વાંચે છે.'' જો ટ્યુરિંગ મશીનનો
પ્રોગ્રામ ઇનપુટ~$\tuple{q_1, 1}$ માટે ત્રિક
$\tuple{q_2, 0, \TMstay}$ આપે, તો મશીન હવે ત્રીજા ખાનાનો
$1$ બદલીને $0$ લખશે, વાંચન-લેખન હેડને એ જ સ્થાને રાખશે અને
અવસ્થા~$q_2$માં જશે. પછી જો પ્રોગ્રામ ઇનપુટ~$\tuple{q_2, 0}$
માટે $\tuple{q_3, 0, \TMright}$ આપે, તો મશીન હવે તે $0$ પર
ફરી $0$ લખશે (એટલે હેડની નીચેની ટેપની સામગ્રી હકીકતમાં
બદલાશે નહીં), એક ખાનું જમણે ખસશે અને અવસ્થા~$q_3$માં
પ્રવેશશે. આ રીતે આગળ ચાલે છે.

મશીન કોઈ અવસ્થા~$q_n$ અને પ્રતીક~$\sigma$ની એવી જોડી પાસે
પહોંચે કે જેના માટે $\tuple{q_n, \sigma}$ને અનુરૂપ કોઈ સૂચના
ન હોય---એટલે ઇનપુટ $\tuple{q_n,\sigma}$ પર સંક્રમણ વિધેય
અવ્યાખ્યાયિત હોય---
ત્યારે આપણે કહીએ છીએ કે મશીન \emph{થંભે} છે. બીજા શબ્દોમાં,
મશીન પાસે કરવા માટે કોઈ સૂચના રહેતી નથી અને તે બિંદુએ તેની
ક્રિયા બંધ થાય છે. થંભવું ક્યારેક વિશિષ્ટ થંભવાની
અવસ્થા~$h$થી દર્શાવાય છે. આ બાબત આગળ વધુ વિગતે બતાવીશું.
\end{explain}

\begin{digress}
ટ્યુરિંગના ``સંગણનીય સંખ્યાઓ અંગે'' (``On computable numbers'') લેખની
વિશેષતા એ છે કે તેમાં
માત્ર ઔપચારિક વ્યાખ્યા જ નથી, પરંતુ એ વ્યાખ્યા સંગણનીયતાની
સહજ કલ્પનાને પકડે છે એવી દલીલ પણ છે. વ્યાખ્યામાંથી એ સ્પષ્ટ
થવું જોઈએ કે ટ્યુરિંગ મશીન વડે સંગણનીય દરેક વિધેય સહજ અર્થમાં
સંગણનીય છે. ટ્યુરિંગે ઊલટી દિશા માટે---એટલે જેને આપણે
સ્વાભાવિક રીતે સંગણનીય ગણીએ એવું દરેક વિધેય આવા મશીન વડે
ગણી શકાય છે તે માટે---ત્રણ પ્રકારની દલીલો આપી. તેમના પોતાના
શબ્દોમાં તે આ પ્રમાણે છે:
\begin{enumerate}
\item સહજ સમજને સીધી અપીલ.
\item બે વ્યાખ્યાઓ સમકક્ષ છે તેની સાબિતી (નવી વ્યાખ્યા સહજ
  સમજને વધુ ગમે એવી હોય ત્યારે).
\item સંગણનીય સંખ્યાઓના મોટા વર્ગોનાં ઉદાહરણો આપવાં.
\end{enumerate}
આપણો હેતુ સંગણનીયતાને ``સૈદ્ધાંતિક રીતે'' વ્યાખ્યાયિત કરવાનો
છે, એટલે સમય અને સ્થાનની વ્યવહારુ મર્યાદાઓ ધ્યાનમાં લીધા
વિના. અલબત્ત, સંગણનીયતાની વ્યાપકતમ વ્યાખ્યા હાથવગી થયા
પછી મર્યાદિત સંસાધનો સાથેના સંગણનનો પણ અભ્યાસ કરી શકાય;
તે ``સંગણકીય જટિલતા'' તરીકે ઓળખાતા વિષયનું કેન્દ્ર છે.
\end{digress}

\begin{history}
એલન ટ્યુરિંગે 1936માં ટ્યુરિંગ મશીનોની શોધ કરી. તે સમયે તેમનો
રસ પ્રથમ ક્રમના તર્કશાસ્ત્રની નિર્ણેયતામાં હતો, છતાં તેમના
લેખને કમ્પ્યુટરની રચનાના પાયા અંગેનો નિર્ણાયક લેખ કહેવામાં
આવ્યો છે. તેમાં ટ્યુરિંગ સંગણનીય વાસ્તવિક સંખ્યાઓ પર ધ્યાન
આપે છે, એટલે એવી વાસ્તવિક સંખ્યાઓ કે જેના દશાંશ વિસ્તરણો
ગણી શકાય; પરંતુ તેઓ નોંધે છે કે પ્રાકૃતિક સંખ્યાઓ પરનાં
સંગણનીય વિધેયો વગેરે માટે પણ આ વિચારોને અનુકૂળ બનાવવાનું
અઘરું નથી. નોંધો કે પ્રથમ કાર્યરત સામાન્ય હેતુનો કમ્પ્યુટર
1941માં બનાવાયો તેના પૂરા પાંચ વર્ષ પહેલાંની આ વાત છે
(જર્મન કોનરાડ ઝૂઝેએ પોતાના માતાપિતાના ઘરના બેઠકખંડમાં તે
બનાવ્યો હતો); બ્લેચલી પાર્કમાં ટ્યુરિંગ અને તેમના સાથીઓએ
સંકેત ઉકેલતું કોલોસસ બનાવ્યું (1943) તેના સાત વર્ષ પહેલાં;
અમેરિકન ENIAC (1945)ના નવ વર્ષ પહેલાં; મૅન્ચેસ્ટરમાં પ્રથમ
બ્રિટિશ સામાન્ય હેતુનો કમ્પ્યુટર---Manchester Small-Scale
Experimental Machine---બન્યો (1948) તેના બાર વર્ષ પહેલાં;
અને અમેરિકનોએ BINAC (1949)નું પ્રથમ પરીક્ષણ કર્યું તેના તેર
વર્ષ પહેલાં. Manchester SSEM પ્રથમ સંગ્રહિત-પ્રોગ્રામવાળો
કમ્પ્યુટર હતો---અગાઉનાં મશીનોને દરેક નવા કામ માટે હાથેથી
ફરી તાર જોડવા પડતા હતા.
\end{history}
% END OLP-0254
\endgroup


\begingroup
\makeatletter\def\ol@sectioncs{section}\makeatother

% BEGIN OLP-0255
\olfileid{tur}{mac}{rep}
\olsection{ટ્યુરિંગ મશીનોનું નિરૂપણ}
\phantomsection\label{guunit:OLP-0255}


\begin{explain}
ટ્યુરિંગ મશીનોને \emph{અવસ્થા-આલેખો} વડે ચિત્રરૂપે દર્શાવી
શકાય છે. આવા આલેખોમાં તીરોથી જોડાયેલાં અવસ્થાનાં ખાનાં હોય
છે. દરેક અવસ્થાનું ખાનું મશીનની એક અવસ્થા દર્શાવે છે. દરેક
તીર તે અવસ્થામાંથી અમલમાં મૂકી શકાય એવી એક સૂચના દર્શાવે છે;
તે સૂચનાની વિગતો સંબંધિત તીરની ઉપર અથવા નીચે લખેલી હોય
છે. નીચેના મશીનમાં માત્ર બે આંતરિક અવસ્થાઓ, $q_0$ અને
$q_1$, તથા એક સૂચના છે:
\[
\begin{tikzpicture}[->,>=stealth',shorten >=1pt,auto,node distance=2.8cm,
                    semithick]
  \tikzstyle{every state}=[fill=none,draw=black,text=black]

  \node[initial,state]   (A)              {$q_0$};
  \node[state]   (B) [right of=A] {$q_1$};

  \path (A) edge  node {\TMtrans{\TMblank}{\TMstroke}{\TMright}} (B);
\end{tikzpicture}
\]
યાદ કરો કે ટ્યુરિંગ મશીન પાસે વાંચન-લેખન હેડ અને તેના પર
ઇનપુટ લખેલી ટેપ હોય છે. આ સૂચનાનો અર્થ છે: \emph{અવસ્થા
$q_0$માં $\TMblank$ વાંચતા હો, તો $\TMstroke$ લખો, જમણે ખસો
અને અવસ્થા $q_1$માં જાઓ}. આ વાત સંક્રમણ વિધેયમાં
$\tuple{q_0, \TMblank}$ને $\tuple{q_1, \TMstroke,
\TMright}$ સાથે જોડવા સમકક્ષ છે.
\end{explain}

\begin{ex}
\emph{સમ મશીન}: ટેપ પર $\TMstroke$ની સંખ્યા સમ હોય ત્યારે
અને ત્યારે જ નીચેનું ટ્યુરિંગ મશીન થંભે છે (એવી ધારણા હેઠળ કે
ટેપ પરના બધા $\TMstroke$ પ્રથમ $\TMblank$ પહેલાં આવે છે).
\[
\begin{tikzpicture}[->,>=stealth',shorten >=1pt,auto,node distance=2.8cm,
                    semithick]
  \tikzstyle{every state}=[fill=none,draw=black,text=black]

  \node[initial,state]         (A)              {$q_0$};
  \node[state]         (B) [right of=A] {$q_1$};

  \path (A) edge [bend left] node {\TMtrans{\TMstroke}{\TMstroke}{\TMright}} (B)
        (B) edge [loop above] node {\TMtrans{\TMblank}{\TMblank}{\TMright}} (B)
            edge [bend left] node {\TMtrans{\TMstroke}{\TMstroke}{\TMright}} (A);
\end{tikzpicture}
\]

આ અવસ્થા-આલેખને નીચેનું સંક્રમણ વિધેય અનુરૂપ છે:
\begin{align*}
\delta(q_0, \TMstroke) & = \tuple{q_1, \TMstroke, \TMright},\\
\delta(q_1, \TMstroke) & = \tuple{q_0, \TMstroke, \TMright},\\
\delta(q_1, \TMblank)  & = \tuple{q_1, \TMblank, \TMright}
\end{align*}
\end{ex}

\begin{explain}
ઉપરનું મશીન ઇનપુટમાં સ્ટ્રોકની સંખ્યા સમ હોય ત્યારે
જ થંભે છે. નહીંતર મશીન (સૈદ્ધાંતિક રીતે) અનંત સમય સુધી
ચાલુ રહે છે. કોઈ પણ મશીન અને ઇનપુટ માટે આઉટપુટ નક્કી
કરવા મશીનની \emph{સંરૂપણો}નો ક્રમ અનુસરવો શક્ય છે.
સંરૂપણની ઔપચારિક વ્યાખ્યા આપણે આગળ આપીશું. હાલ પૂરતું,
મશીન ચાલે ત્યારે જુદા જુદા સમયે તેની સ્થિતિ દર્શાવતા આલેખોની
શ્રેણી તરીકે સંરૂપણોને સહજ રીતે સમજી શકાય. સંરૂપણમાં ટેપની
સામગ્રી, મશીનની અવસ્થા અને વાંચન-લેખન હેડનું સ્થાન દેખાય છે.

ચાર $\TMstroke$નું ઇનપુટ આપીએ ત્યારે સમ મશીનનાં સંરૂપણો
ક્રમશઃ જોઈએ. આ કિસ્સામાં મશીન થંભશે એવી આપણી અપેક્ષા
છે. પછી તેને ત્રણ $\TMstroke$નું ઇનપુટ આપીશું; ત્યારે તે
હંમેશાં ચાલતું રહેશે.

મશીન અવસ્થા~$q_0$માં શરૂ થાય છે અને સૌથી ડાબો
$\TMstroke$ વાંચે છે. મશીનની આ પ્રારંભિક સ્થિતિ આ રીતે
દર્શાવી શકાય:
\[
\TMendtape \TMstroke_0 \TMstroke \TMstroke \TMstroke \TMblank \ldots
\]
ઉપરનું સંરૂપણ સમજવું સહેલું છે. મશીન પ્રથમ અવસ્થામાં
સૌથી ડાબો $\TMstroke$ વાંચતાં શરૂ થાય છે. પ્રથમ
$\TMstroke$ને અવસ્થાના નામનો નિમ્નસૂચક લગાવીને આ દર્શાવ્યું
છે. આ સમયે લાગુ પડતી સૂચના
$\delta(q_0, \TMstroke) = \tuple{q_1, \TMstroke, \TMright}$
છે. તેથી મશીન ટેપ પર જમણે ખસે છે અને અવસ્થા~$q_1$માં
જાય છે.
\[
\TMendtape \TMstroke \TMstroke_1 \TMstroke \TMstroke \TMblank \ldots
\]
હવે મશીન અવસ્થા~$q_1$માં $\TMstroke$ વાંચતું હોવાથી,
આપણે સૂચના $\delta(q_1, \TMstroke) = \tuple{q_0,
  \TMstroke, \TMright}$ને ``અનુસરવી'' પડશે. પરિણામે આ
સંરૂપણ મળે છે:
\[
\TMendtape \TMstroke \TMstroke \TMstroke_0 \TMstroke \TMblank \ldots
\]
મશીન આગળ ચાલે છે ત્યારે આ જ ક્રમમાં નિયમો ફરી લાગુ પડે
છે અને નીચેના બે સંરૂપણો મળે છે:
\[
\TMendtape \TMstroke \TMstroke \TMstroke \TMstroke_1 \TMblank \ldots
\]
\[
\TMendtape \TMstroke \TMstroke \TMstroke \TMstroke \TMblank_0 \ldots
\]
હવે મશીન અવસ્થા~$q_0$માં $\TMblank$ વાંચે છે. સંક્રમણ
આલેખ પરથી સહેલાઈથી દેખાય છે કે હવે અમલમાં મૂકવા માટે કોઈ
સૂચના નથી, એટલે મશીન થંભી ગયું છે. તેનો અર્થ છે કે
ઇનપુટ સ્વીકારાયું છે.

હવે ધારો કે મશીનને ત્રણ $\TMstroke$નું ઇનપુટ આપીએ.
ટેપ પરના ઇનપુટમાં નાનો તફાવત હોવા છતાં એ જ સૂચનાઓ લાગુ
પડતી હોવાથી શરૂઆતનાં કેટલાંક સંરૂપણો સરખાં હોય છે:
\[
\TMendtape \TMstroke_0 \TMstroke \TMstroke \TMblank \ldots
\]
\[
\TMendtape \TMstroke \TMstroke_1 \TMstroke \TMblank \ldots
\]
\[
\TMendtape \TMstroke \TMstroke \TMstroke_0 \TMblank \ldots
\]
\[
\TMendtape \TMstroke \TMstroke \TMstroke \TMblank_1 \ldots
\]
મશીન હવે બધા $\TMstroke$ વટાવી ચૂક્યું છે અને
અવસ્થા~$q_1$માં $\TMblank$ વાંચે છે. આલેખમાં બતાવ્યા
પ્રમાણે, અહીં $\delta(q_1, \TMblank) =\tuple{q_1,
\TMblank, \TMright}$ સ્વરૂપની સૂચના છે. ટેપ જમણી તરફ
અનંત સુધી $\TMblank$થી ભરેલી હોવાથી, મશીન
અવસ્થા~$q_1$માં રહીને વધુ ને વધુ જમણે ખસતું આ સૂચના
\emph{હંમેશાં} ચલાવશે. મશીન ક્યારેય થંભશે નહીં અને ઇનપુટ
સ્વીકારશે નહીં.
\end{explain}

\begin{explain}
દરેક મશીન થંભતું નથી, એ નોંધવું અગત્યનું છે. થંભવાનો અર્થ
અમલમાં મૂકવા માટેની સૂચનાઓ ખૂટી જવી એવો હોય, તો દરેક
અવસ્થામાં દરેક પ્રતીક માટે બહાર જતું તીર છે તેની ખાતરી
કરીને ક્યારેય ન થંભે એવું મશીન બનાવી શકાય. અવસ્થા~$q_0$માં
$\TMblank$ વાંચવા માટેની સૂચના ઉમેરીને સમ મશીનને અનંત સમય
સુધી ચાલે એવું બનાવી શકાય.
\end{explain}

\begin{ex}
\[
\begin{tikzpicture}[->,>=stealth',shorten >=1pt,auto,node distance=2.8cm,
                    semithick]
  \tikzstyle{every state}=[fill=none,draw=black,text=black]

  \node[initial,state] (A)              {$q_0$};
  \node[state]         (B) [right of=A] {$q_1$};

  \path
  (A) edge [bend left] node {\TMtrans{\TMstroke}{\TMstroke}{\TMright}} (B)
      edge [loop above] node {\TMtrans{\TMblank}{\TMblank}{\TMright}} (A)
  (B) edge [loop above] node {\TMtrans{\TMblank}{\TMblank}{\TMright}} (B)
      edge [bend left] node {\TMtrans{\TMstroke}{\TMstroke}{\TMright}} (A);
\end{tikzpicture}
\]
\end{ex}

\begin{explain}
ટ્યુરિંગ મશીનો દર્શાવવાની બીજી રીત \emph{મશીન-કોષ્ટકો}
છે. કોષ્ટકની $x$-અક્ષ પર ટેપની વર્ણમાળા અને $y$-અક્ષ પર
મશીનની અવસ્થાઓનો ગણ દર્શાવાય છે. દરેક અવસ્થા અને પ્રતીકના
છેદે સૂચનાનો બાકીનો ભાગ---નવી અવસ્થા, નવું પ્રતીક અને ખસવાની
દિશા---લખાય છે. કઈ અવસ્થામાં અને કયું પ્રતીક વાંચતાં મશીન
થંભે છે તે કોષ્ટકથી સહેલાઈથી જાણી શકાય છે. કોષ્ટકમાં કોઈ
ખાનું ખાલી હોય, તો ત્યાં મશીન થંભી શકે છે. અવસ્થા-આલેખો
અને સૂચનાના ગણમાં મશીન ક્યાં થંભે છે તે હંમેશાં તરત દેખાતું
નથી; મશીન-કોષ્ટકમાં ખાલી ખાનાં જોઈને આ બિંદુઓ તરત ઓળખી
શકાય છે.
\end{explain}

\begin{ex}
સમ મશીનનું મશીન-કોષ્ટક આ છે:
\[
\centering
\begin{tabular}{lllll}
\cline{2-4}
\multicolumn{1}{l|}{}      & \multicolumn{1}{c|}{$\TMblank$}
& \multicolumn{1}{c|}{$\TMstroke$}     & \multicolumn{1}{c|}{$\TMendtape$}   &  \\ \cline{1-4}
\multicolumn{1}{|l|}{$q_0$} & \multicolumn{1}{l|}{}
& \multicolumn{1}{l|}{\TMtrans{\TMstroke}{q_1}{\TMright}} &
\multicolumn{1}{l|}{\phantom{\TMtrans{\TMstroke}{q_1}{\TMright}}} &  \\ \cline{1-4}
\multicolumn{1}{|l|}{$q_1$} & \multicolumn{1}{l|}{\TMtrans{\TMblank}{q_1}{\TMright}}
& \multicolumn{1}{l|}{\TMtrans{\TMstroke}{q_0}{\TMright}} &
\multicolumn{1}{l|}{\phantom{\TMtrans{\TMstroke}{q_1}{\TMright}}} &  \\ \cline{1-4}
\end{tabular}
\]
જેમ જોઈ શકાય છે, મશીન અવસ્થા~$q_0$માં $\TMblank$
વાંચે ત્યારે થંભે છે.
\end{ex}

\begin{explain}
અત્યાર સુધી આપણે ઇનપુટ વાંચીને સ્વીકારતાં મશીનો જ જોયાં
છે. પરંતુ ટ્યુરિંગ મશીન વાંચી પણ શકે છે અને લખી પણ શકે
છે. આવા મશીનનું એક ઉદાહરણ (જોકે આવાં ઘણાં ઉદાહરણો છે)
\emph{બમણું કરનાર મશીન} છે. ટેપ પર $n$ $\TMstroke$નો
એક સળંગ સમૂહ આપીને શરૂ કરીએ, તો તે $2n$ $\TMstroke$નો
સમૂહ આઉટપુટ તરીકે આપે છે.
\end{explain}

\begin{ex}
  \ollabel{ex:doubler} બમણું કરનાર મશીન બનાવતાં પહેલાં સમસ્યા
  ઉકેલવાની \emph{રીત} નક્કી કરવી જરૂરી છે. મશીન (આપણે તેને
  જેમ ઘડ્યું છે તે મુજબ) કેટલા $\TMstroke$ વાંચ્યા છે તે યાદ
  રાખી શકતું નથી. તેથી ટેપ પરના બધા $\TMstroke$નો હિસાબ
  રાખવાની રીત શોધવી પડશે. એક રીત ઇનપુટને આઉટપુટથી
  એક $\TMblank$ વડે અલગ કરવાની છે. પછી મશીન ઇનપુટમાંથી
  પ્રથમ $\TMstroke$ ભૂંસી શકે, બાકીના ઇનપુટને વટાવી શકે,
  એક $\TMblank$ છોડી શકે અને બે નવા $\TMstroke$ લખી
  શકે. ત્યાર પછી મશીન પાછું જઈ ઇનપુટનો બીજો $\TMstroke$
  શોધશે અને તેને પણ બમણો કરશે. ઇનપુટના દરેક એક
  $\TMstroke$ માટે તે આઉટપુટમાં બે $\TMstroke$ લખશે.
  મશીન આગળ વધે તેમ ઇનપુટ ભૂંસવાથી કોઈ $\TMstroke$ છૂટી
  ન જાય કે બે વાર બમણો ન થાય તેની ખાતરી કરી શકાય. આખું
  ઇનપુટ ભૂંસાઈ જાય ત્યારે ટેપ પર $2n$ $\TMstroke$ બાકી
  હશે. મળેલા ટ્યુરિંગ મશીનનો અવસ્થા-આલેખ
  \olref{fig:doubler}માં દર્શાવ્યો છે.
  \begin{figure}
\[
\begin{tikzpicture}[->,>=stealth',shorten >=1pt,auto,node distance=2.8cm,
                    semithick]
  \tikzstyle{every state}=[fill=none,draw=black,text=black]

  \node[initial,state] (1)              {$q_0$};
  \node[state]         (2) [right of=1] {$q_1$};
  \node[state]         (3) [right of=2] {$q_2$};
  \node[state]         (4) [below of=3] {$q_3$};
  \node[state]         (5) [left of=4]  {$q_4$};
  \node[state]         (6) [left of=5]  {$q_5$};

  \path (1) edge node {\TMtrans{\TMstroke}{\TMblank}{\TMright}} (2)
    (2) edge [loop above] node {\TMtrans{\TMstroke}{\TMstroke}{\TMright}} (2)
      edge node {\TMtrans{\TMblank}{\TMblank}{\TMright}} (3)
    (3) edge [loop above] node {\TMtrans{\TMstroke}{\TMstroke}{\TMright}} (3)
        edge node {\TMtrans{\TMblank}{\TMstroke}{\TMright}} (4)
    (4) edge [loop below] node {\TMtrans{\TMblank}{\TMstroke}{\TMleft}} (4)
        edge node {\TMtrans{\TMstroke}{\TMstroke}{\TMleft}} (5)
    (5) edge [loop below]  node {\TMtrans{\TMstroke}{\TMstroke}{\TMleft}} (5)
        edge              node {\TMtrans{\TMblank}{\TMblank}{\TMleft}} (6)
    (6) edge [loop below] node {\TMtrans{\TMstroke}{\TMstroke}{\TMleft}} (6)
        edge              node {\TMtrans{\TMblank}{\TMblank}{\TMright}} (1);
\end{tikzpicture}
\]
\caption{બમણું કરનાર મશીન}
\ollabel{fig:doubler}
\end{figure}
\end{ex}

\begin{prob}
કોઈ પણ એક ઇનપુટ પસંદ કરો અને \olref[tur][mac][rep]{ex:doubler}માંના
બમણું કરનાર મશીનનાં સંરૂપણોનો ક્રમ અનુસરો.
\end{prob}

\begin{prob}
$\{\TMendtape,\TMblank, A, B\}$ વર્ણમાળા ધરાવતું એવું
ટ્યુરિંગ મશીન રચો કે જેમાં $A$ અને $B$ની સંખ્યા સરખી
હોય \emph{અને} બધા $A$, બધા $B$ પહેલાં આવતા હોય,
એવી $A$ અને $B$ની કોઈ પણ શબ્દમાળા તે સ્વીકારે,
એટલે તેના પર થંભે; અને $A$ની સંખ્યા $B$ની સંખ્યા
જેટલી ન હોય અથવા બધા $A$ બધા $B$ પહેલાં ન આવતા
હોય એવી શબ્દમાળા તે નકારે, એટલે તેના પર ન થંભે.
(દાખલા તરીકે, મશીને $AABB$ અને $AAABBB$ સ્વીકારવાં,
પણ $AAB$ અને $AABBAABB$ બંને નકારવાં જોઈએ.)
\end{prob}

\begin{prob}
$\{\TMendtape,\TMblank, A, B\}$ વર્ણમાળા ધરાવતું એવું
ટ્યુરિંગ મશીન રચો કે જે $A$ અને $B$ની કોઈ પણ
શબ્દમાળા~$\alpha$ ઇનપુટ તરીકે લે અને તેની નકલ કરીને
$\alpha\alpha$ સ્વરૂપનું આઉટપુટ આપે. (દાખલા તરીકે,
ઇનપુટ $ABBA$નું આઉટપુટ $ABBAABBA$ થવું જોઈએ.)
\end{prob}

\begin{prob}
\emph{વર્ણક્રમમાં?:} $\{\TMendtape,\TMblank, A, B\}$ વર્ણમાળા
ધરાવતું એવું ટ્યુરિંગ મશીન રચો કે જેને $A$ અને $B$નો
સાન્ત અનુક્રમ ઇનપુટ તરીકે આપતાં બધા $A$, બધા $B$ની
ડાબે આવે છે કે નહીં તે તપાસે. મશીને ઇનપુટ શબ્દમાળા ટેપ
પર જ રાખવી જોઈએ અને શબ્દમાળા ``વર્ણક્રમમાં'' હોય તો
થંભવું, નહીંતર હંમેશાં ચક્રમાં ચાલવું જોઈએ.
\end{prob}

\begin{prob}
\emph{વર્ણક્રમમાં ગોઠવનાર:} $\{\TMendtape,\TMblank, A, B\}$
વર્ણમાળા ધરાવતું એવું ટ્યુરિંગ મશીન રચો કે જે $A$ અને
$B$ના સાન્ત અનુક્રમને ઇનપુટ તરીકે લઈ તેમને ફરી ગોઠવે,
જેથી બધા $A$, બધા~$B$ની ડાબે હોય. (દાખલા તરીકે,
$BABAA$ અનુક્રમ $AAABB$ બનવો જોઈએ અને $ABBABB$
અનુક્રમ $AABBBB$ બનવો જોઈએ.)
\end{prob}
% END OLP-0255
\endgroup


\begingroup
\makeatletter\def\ol@sectioncs{section}\makeatother

% BEGIN OLP-0256
\olfileid{tur}{mac}{tur}
\olsection{ટ્યુરિંગ મશીનો}
\phantomsection\label{guunit:OLP-0256}


\begin{explain}
ટ્યુરિંગ મશીનની ઔપચારિક વ્યાખ્યા અમૂર્ત લાગે છે, પણ ખરેખર
સરળ છે: ટ્યુરિંગ મશીન કેવી રીતે કામ કરે છે તે નિર્દિષ્ટ કરવા
જરૂરી બધી માહિતી તે એક જ ગણિતીય રચનામાં સમાવે છે. તેમાં
(1) મશીન કઈ અવસ્થાઓમાં હોઈ શકે, (2) ટેપ પર કયાં પ્રતીકોને
મંજૂરી છે, (3) મશીન કઈ અવસ્થામાં શરૂ થવું જોઈએ અને
(4) મશીનનો સૂચનાઓનો ગણ શું છે, એ બધું આવે છે.
\end{explain}

\begin{defn}[ટ્યુરિંગ મશીન]
\emph{ટ્યુરિંગ મશીન} $M$ એ નીચેની વસ્તુઓનું બનેલું ચતુષ્ટય
$\langle Q, \Sigma, q_0, \delta\rangle$ છે:
\begin{enumerate}
\item \emph{અવસ્થાઓ}નો સાન્ત ગણ~$Q$,
\item $\TMendtape$ અને $\TMblank$ ધરાવતી સાન્ત \emph{વર્ણમાળા}
  $\Sigma$,
\item \emph{પ્રારંભિક અવસ્થા}~$q_0 \in Q$,
\item સાન્ત \emph{સૂચનાઓનો ગણ}~$\delta\colon Q \times \Sigma
  \pto Q \times \Sigma \times \{\TMleft, \TMright, \TMstay\}$.
\end{enumerate}
આંશિક વિધેય~$\delta$ને $M$નું \emph{સંક્રમણ વિધેય}
પણ કહે છે.
\end{defn}

\begin{explain}
આપણે ધારીએ છીએ કે ટેપ માત્ર એક દિશામાં અનંત છે.
તેથી ટેપના ડાબા છેડાને ચિહ્નિત કરવા માટે વિશિષ્ટ
પ્રતીક~$\TMendtape$ રાખવું ઉપયોગી છે. તેનાથી ટ્યુરિંગ
મશીનના પ્રોગ્રામને ટેપની બહાર નીકળી જવાનો ``ભય'' ક્યારે
છે તે ઓળખવાનું સરળ બને છે. આપણે એમ ધારી શકીએ કે આ
પ્રતીક ક્યારેય બદલવામાં આવતું નથી, એટલે
$\delta(q,\TMendtape)$ વ્યાખ્યાયિત હોય તો
$\delta(q,\TMendtape) = \tuple{q', \TMendtape, x}$.
કેટલાંક પાઠ્યપુસ્તકો આવી ધારણા કરે છે; આપણે કરતા નથી.
તમે તમારું ટ્યુરિંગ મશીન ઘડતી વખતે તે ક્યારેય
$\TMendtape$ને બદલે નહીં તેની કાળજી રાખી શકો છો.
ઉપરાંત, કેટલીક પરિસ્થિતિમાં આ પ્રતીક બદલવાની છૂટ
આપવાથી ઉપયોગી લવચીકતા મળે છે.
\end{explain}

\begin{ex}
\emph{સમ મશીન:} ઔપચારિક રીતે સમ મશીન એ ચતુષ્ટય
$\tuple{Q, \Sigma, q_0, \delta}$ છે, જ્યાં
\begin{align*}
Q & = \{ q_0, q_1 \} \\
\Sigma & = \{ \TMendtape, \TMblank, \TMstroke \}, \\
\delta(q_0, \TMstroke) & = \tuple{q_1, \TMstroke, \TMright},\\
\delta(q_1, \TMstroke) & = \tuple{q_0, \TMstroke, \TMright},\\
\delta(q_1, \TMblank)  & = \tuple{q_1, \TMblank, \TMright}.
\end{align*}
\end{ex}
% END OLP-0256
\endgroup


\begingroup
\makeatletter\def\ol@sectioncs{section}\makeatother

% BEGIN OLP-0257
\olfileid{cmp}{tur}{con}
\olsection{સંરૂપણો અને સંગણનાઓ}
\phantomsection\label{guunit:OLP-0257}


\begin{explain}
અગાઉ આપણે સમ મશીનનાં સંરૂપણોનો ક્રમ અનુસર્યો હતો તે યાદ
કરો. ટેપ, વાંચન-લેખન હેડ અને ટ્યુરિંગ મશીનના પ્રોગ્રામથી
બનેલું કલ્પિત ઉપકરણ ખરેખર ટ્યુરિંગ મશીનની સંગણના શું છે
તેની સહજ ચિત્રાત્મક સમજ આપવાની રીત છે. ઔપચારિક રીતે,
આપેલા ઇનપુટ પર ટ્યુરિંગ મશીનની સંગણનાને
\emph{સંરૂપણો}ના અનુક્રમ તરીકે વ્યાખ્યાયિત કરી શકાય છે.
સંરૂપણ પોતે પ્રતીકોનો અનુક્રમ (સંગણનાના કોઈ સમયે ટેપ
પર શું લખેલું છે તે દર્શાવતો), વાંચન-લેખન હેડનું સ્થાન
જણાવતી સંખ્યા અને એક અવસ્થા છે. આ બધાની મદદથી આપણે
આપેલા ઇનપુટ પર ટ્યુરિંગ મશીન~$M$ શું ગણે છે તે
વ્યાખ્યાયિત કરી શકીએ છીએ.
\end{explain}

\begin{defn}[સંરૂપણ]
ટ્યુરિંગ મશીન $M = \tuple{Q, \Sigma, q_0,
\delta}$નું \emph{સંરૂપણ} એ $\tuple{C, m, q}$ ત્રિક છે, જ્યાં
\begin{enumerate}
\item $C \in \Sigma^*$ એ $\Sigma$નાં પ્રતીકોનો સાન્ત અનુક્રમ છે,
\item $m \in \Nat$ એ $< \len{C}$ એવી સંખ્યા છે, અને
\item $q \in Q$ છે.
\end{enumerate}
સહજ રીતે, અનુક્રમ~$C$ ટેપની સામગ્રી છે (સૌથી ડાબા
ખાનાથી માંડીને છેલ્લાં ખાલી ન હોય અથવા અગાઉ મુલાકાત
લીધેલાં ખાનાં સુધીનાં પ્રતીકો); $m$ વાંચન-લેખન હેડ જે
ખાનું વાંચે છે તેનો ક્રમાંક છે (સૌથી ડાબા ખાનાનો ક્રમાંક
$0$ રાખીને), અને $q$ મશીનની વર્તમાન અવસ્થા છે.
\end{defn}

\begin{explain}
ટ્યુરિંગ મશીનનું સંભવિત ઇનપુટ પ્રતીકોનો અનુક્રમ હોય છે,
સામાન્ય રીતે કોઈ સ્વરૂપમાં એક સંખ્યાનું સંકેતીકરણ કરતો
અનુક્રમ. ટ્યુરિંગ મશીનનું પ્રારંભિક સંરૂપણ એ સંરૂપણ છે જેમાં
મશીનને તે ઇનપુટ પર કામ કરવા શરૂ કરીએ છીએ: ટેપ પર પહેલાં
ટેપના છેડાનું ચિહ્ન હોય છે અને તેની તરત જમણેનાં ખાનાંમાં
ઇનપુટ લખેલું હોય છે; વાંચન-લેખન હેડ ઇનપુટનું સૌથી ડાબું
ખાનું (એટલે છેડાના ચિહ્નની જમણેનું ખાનું) વાંચે છે; અને
ઉપકરણ નિર્દિષ્ટ પ્રારંભિક અવસ્થા~$q_0$માં હોય છે.
\end{explain}

\begin{defn}[પ્રારંભિક સંરૂપણ]
ઇનપુટ $I \in \Sigma^*$ માટે $M$નું \emph{પ્રારંભિક સંરૂપણ}
આ છે:
\[
\tuple{\TMendtape \frown I, 1, q_0}.
\]
\end{defn}

\sourcecorrection{OLTM-001}{સ્થિર અંગ્રેજી મૂળનો પરિચય
શરૂઆતમાં હેડ ટેપનું સૌથી ડાબું ખાનું વાંચે છે એમ કહે છે;
અહીંની ઔપચારિક વ્યાખ્યા અને $m=1$નું સૂત્ર હેડને
ઇનપુટના પ્રથમ ખાનામાં, છેડાના ચિહ્નની જમણે, મૂકે છે.
આ વિભાગમાં પ્રારંભિક સંરૂપણ માટે ઔપચારિક સૂત્ર અનુસર્યું છે.}

\begin{explain}
ચિહ્ન~$\frown$ \emph{અનુક્રમોનું જોડાણ} દર્શાવે છે---ઇનપુટ
શબ્દમાળા ડાબા છેડાના ચિહ્નની તરત જમણે શરૂ થાય છે.
\end{explain}

\sourcecorrection{OLTM-002}{સ્થિર અંગ્રેજી મૂળ આ વાક્યમાં
ઇનપુટ છેડાના ચિહ્નની તરત ``ડાબે'' શરૂ થાય છે એમ કહે છે.
પરંતુ બાજુનું $\TMendtape \frown I$ સૂત્ર અને અગાઉનું
ટેપનું વર્ણન ઇનપુટને તેની જમણે મૂકે છે; ગુજરાતી ગદ્યમાં
તે મુજબ દિશા સ્પષ્ટ કરી છે.}

\begin{defn}
અપણે કહીએ છીએ કે (મશીન~$M$ મુજબ) સંરૂપણ
$\tuple{C, m, q}$ \emph{એક પગલામાં સંરૂપણ
$\tuple{C', m', q'}$ આપે છે}, ત્યારે અને ત્યારે જ જ્યારે
\begin{enumerate}
\item $C$નું $m$-મું પ્રતીક $\sigma$ છે,
\item $M$નો સૂચનાઓનો ગણ $\delta(q, \sigma) =
  \tuple{q', \sigma', D}$ નિર્દિષ્ટ કરે છે,
\item $C'$નું $m$-મું પ્રતીક $\sigma'$ છે, અને
\item
\begin{enumerate}
\item $D = L$ હોય, અને એ કિસ્સામાં $m>0$ હોય તો
  $m' = m - 1$, નહીંતર $m'=0$; અથવા
\item $D = R$ હોય અને $m' = m + 1$ હોય, અથવા
\item $D = N$ હોય અને $m' = m$ હોય,
\end{enumerate}
\item $m' = \len{C}$ હોય તો $\len{C'} = \len{C} + 1$ છે
  અને $C'$નું $m'$-મું પ્રતીક~$\TMblank$ છે. નહીંતર
  $\len{C'}=\len{C}$ છે.
\item $i < \len{C}$ અને $i \neq m$ એવા દરેક $i$ માટે
  $C'(i) = C(i)$ છે,
\end{enumerate}
\end{defn}

\begin{defn}\ollabel{defn:run-output}
\emph{ઇનપુટ~$I$ પર $M$નું ચલન} એ $M$નાં સંરૂપણોનો
અનુક્રમ $C_i$ છે, જેમાં $C_0$ એ ઇનપુટ~$I$ માટે
$M$નું પ્રારંભિક સંરૂપણ છે અને દરેક $C_i$ એક પગલામાં
$C_{i+1}$ આપે છે.

જો $C_k = \tuple{C, m, q}$ હોય, $C$નું $m$-મું
પ્રતીક~$\sigma$ હોય અને $\delta(q, \sigma)$ અવ્યાખ્યાયિત
હોય, તો આપણે કહીએ છીએ કે $M$ \emph{ઇનપુટ $I$ પર
$k$ પગલાં પછી થંભે છે}. એ કિસ્સામાં, ઇનપુટ~$I$ માટે
$M$નું \emph{આઉટપુટ} $O$ છે, જ્યાં $O$ એ
$\TMblank$માં પૂરું ન થતું એવું પ્રતીકોનું શબ્દમાળા છે કે
કોઈ~$j \in \Nat$ માટે
$C = \TMendtape \concat O \concat \TMblank^j$.
($\TMblank^j$ એ $j$ $\TMblank$નો અનુક્રમ છે.)
\end{defn}

\begin{explain}
આ વ્યાખ્યા મુજબ $M$નું આઉટપુટ~$O$ હંમેશાં
$\TMblank$ સિવાયના પ્રતીકમાં પૂરું થાય છે; અથવા
સમય~$k$ પર આખી ટેપ $\TMblank$થી ભરેલી હોય
(સૌથી ડાબા $\TMendtape$ સિવાય), તો $O$ ખાલી
શબ્દમાળા છે.
\end{explain}
% END OLP-0257
\endgroup


\begingroup
\makeatletter\def\ol@sectioncs{section}\makeatother

% BEGIN OLP-0258
\olfileid{tur}{mac}{una}
\olsection{સંખ્યાઓનું એક-પ્રતીકી નિરૂપણ}
\phantomsection\label{guunit:OLP-0258}


\begin{explain}
ટ્યુરિંગ મશીનો પોતાની ટેપ પર લખાયેલા પ્રતીકોના અનુક્રમો
પર કામ કરે છે. ટ્યુરિંગ મશીન કઈ વર્ણમાળા વાપરે છે તેના
આધારે આ અનુક્રમો જુદા જુદા ઇનપુટ અને આઉટપુટ દર્શાવી
શકે છે. ખાસ રસ, અલબત્ત, \emph{અંકગણિતીય} વિધેયો,
એટલે પ્રાકૃતિક સંખ્યાઓનાં વિધેયો ગણતા ટ્યુરિંગ મશીનોમાં
છે. ધન પૂર્ણાંકો દર્શાવવાની સરળ રીત તેમને એક જ
પ્રતીક~$\TMstroke$ના અનુક્રમ તરીકે સંકેતિત કરવાની છે.
$n \in \Nat$ હોય, તો $n = 0$ માટે $\TMstroke^n$
ખાલી અનુક્રમ અને અન્યથા બરાબર $n$ $\TMstroke$ ધરાવતો
અનુક્રમ છે.
\end{explain}

\begin{defn}[સંગણના]
ટ્યુરિંગ મશીન~$M$ વિધેય $f\colon \Nat^k \to \Nat$
\emph{ગણે છે} ત્યારે અને ત્યારે જ જ્યારે $M$ ઇનપુટ
\[
\TMstroke^{n_1} \TMblank \TMstroke^{n_2} \TMblank \dots \TMblank \TMstroke^{n_k}
\]
પર થંભે અને આઉટપુટ $\TMstroke^{f(n_1, \dots, n_k)}$
આપે.
\end{defn}

\begin{prob}
  ટ્યુરિંગ મશીન~$M$ વિધેય $f\colon \Nat^k \to \Nat^m$
  ક્યારે ગણે છે તેની વ્યાખ્યા આપો.
\end{prob}

\begin{ex}\ollabel{ex:adder}
\emph{સરવાળો:}
ચાલો $f(n,m) = n + m$ વિધેય ગણતું મશીન બનાવીએ.
તે માટે ટેપ પર $n$ અને~$m$ લંબાઈના $\TMstroke$ના
બે સમૂહો લઈને શરૂ થતું અને $n+m$~$\TMstroke$ના એક
સમૂહ સાથે થંભતું મશીન જોઈએ. ઇનપુટના બે $\TMstroke$
સમૂહોને એક~$\TMblank$ અલગ કરે છે. તેથી એક રીત
$\TMblank$ ધરાવતા ખાનામાં સ્ટ્રોક લખી અને છેલ્લો
$\TMstroke$ ભૂંસી નાખવાની છે.
\begin{figure}\[
\begin{tikzpicture}[->,>=stealth',shorten >=1pt,auto,node distance=2.8cm,
                    semithick]
  \tikzstyle{every state}=[fill=none,draw=black,text=black]

  \node[initial,state]         (A)              {$q_0$};
  \node[state]         (B) [right of=A] {$q_1$};
  \node[state]         (C) [right of=B] {$q_2$};

  \path (A) edge node {\TMtrans{\TMblank}{\TMstroke}{\TMstay}} (B)
           edge [loop above] node {\TMtrans{\TMstroke}{\TMstroke}{\TMright}} (B)
        (B) edge [loop above] node {\TMtrans{\TMstroke}{\TMstroke}{\TMright}} (B)
            edge node {\TMtrans{\TMblank}{\TMblank}{\TMleft}} (C)
        (C) edge [loop above] node {\TMtrans{\TMstroke}{\TMblank}{\TMstay}} (C);
\end{tikzpicture}
\]\caption{$f(x,y) = x+y$ ગણતું મશીન}
\ollabel{fig:adder}
\end{figure}
\end{ex}

\begin{prob}
ઇનપુટ~$\tuple{3,2}$ માટે
\olref[tur][mac][una]{ex:adder}માંના મશીનનાં સંરૂપણોનો
ક્રમ અનુસરો. મશીન $0+0$ ગણે ત્યારે શું થાય છે?
\end{prob}

\begin{explain}
\olref[rep]{ex:doubler}માં આપણે એવું ટ્યુરિંગ મશીન
જોયું હતું જે ઇનપુટ તરીકે $\TMstroke$નો અનુક્રમ લે છે
અને ટેપ પર તેનાથી બમણા $\TMstroke$ના અનુક્રમ સાથે
થંભે છે---બમણું કરનાર મશીન. પરંતુ બમણા કરેલા
$\TMstroke$ના સમૂહની ડાબે આઉટપુટમાં $\TMblank$ રહે
છે. તેથી, તમે કદાચ ધાર્યું હોય તેમ, તે મશીન ખરેખર
$f(x) = 2x$ વિધેય ગણતું નથી. આ ખામી દૂર કરવાની
બે રીતો આપણે વર્ણવીશું.
\end{explain}

\begin{ex}
\olref{fig:doubler-disc}માંનું મશીન $f(x) = 2x$ વિધેય
ગણે છે. \olref[rep]{ex:doubler}માંનું મશીન ઇનપુટ
ભૂંસીને તેના દરેક $\TMstroke$ માટે છેક જમણે બે
$\TMstroke$ લખે છે; તેના બદલે આ મશીન ઇનપુટના
દરેક $\TMstroke$ માટે જમણે એક જ $\TMstroke$ ઉમેરે
છે. ઇનપુટ ક્યાં પૂરું થાય છે તેનો હિસાબ તેને રાખવો
પડે છે. તેથી તે ઇનપુટ અને ઉમેરેલા સ્ટ્રોક વચ્ચે એક
$\TMblank$ રાખે છે, જેને છેવટે $\TMstroke$થી ભરે
છે. આપણે ઇનપુટમાં ક્યાં છીએ તે પણ ``યાદ'' રાખવું
પડે છે, એટલે ઇનપુટના સમૂહમાંના એક $\TMstroke$ને
કામચલાઉ રીતે $\TMblank$થી બદલીએ છીએ.
  \begin{figure}
\[
\begin{tikzpicture}[->,>=stealth',shorten >=1pt,auto,node distance=2.8cm,
                    semithick]
  \tikzstyle{every state}=[fill=none,draw=black,text=black]

  \node[initial,state] (0)              {$q_0$};
  \node[state]         (1) [above of=0] {$q_1$};
  \node[state]         (2) [above of=1] {$q_2$};
  \node[state]         (3) [right of=2] {$q_3$};
  \node[state]         (4) [below of=3] {$q_4$};
  \node[state]         (5) [below of=4] {$q_5$};
  \node[state]         (6) [above right of=4] {$q_6$};
  \node[state]         (7) [below of=6] {$q_7$};
  \node[state]         (8) [below of=7] {$q_8$};

  \path (0) edge node {\TMtrans{\TMstroke}{\TMblank}{\TMright}} (1)
    (1) edge [loop left] node {\TMtrans{\TMstroke}{\TMstroke}{\TMright}} (1)
      edge node {\TMtrans{\TMblank}{\TMblank}{\TMright}} (2)
    (2) edge [loop above] node {\TMtrans{\TMstroke}{\TMstroke}{\TMright}} (2)
        edge node {\TMtrans{\TMblank}{\TMstroke}{\TMleft}} (3)
    (3) edge [loop above] node {\TMtrans{\TMstroke}{\TMstroke}{\TMleft}} (3)
        edge node[left] {\TMtrans{\TMblank}{\TMblank}{\TMleft}} (4)
    (4) edge        node[left] {\TMtrans{\TMstroke}{\TMstroke}{\TMleft}} (5)
    (5) edge [loop below] node {\TMtrans{\TMstroke}{\TMstroke}{\TMleft}} (5)
        edge              node {\TMtrans{\TMblank}{\TMstroke}{\TMright}} (0)
    (4) edge      node[sloped] {\TMtrans{\TMblank}{\TMstroke}{\TMright}} (6)
    (6) edge              node {\TMtrans{\TMblank}{\TMstroke}{\TMright}} (7)
    (7) edge [loop right] node {\TMtrans{\TMstroke}{\TMstroke}{\TMright}} (7)
        edge node {\TMtrans{\TMblank}{\TMblank}{\TMleft}} (8)
    (8) edge [loop right] node {\TMtrans{\TMstroke}{\TMblank}{\TMstay}} (8);
    \end{tikzpicture}
\]
\caption{$f(x) = 2x$ ગણતું મશીન}
\ollabel{fig:doubler-disc}
\end{figure}
\end{ex}

\begin{ex}\ollabel{ex:mover}
$f(x) = 2x$ ગણવાની બીજી રીત મૂળ બમણું કરનાર મશીન
રાખીને છેવટે બમણા સ્ટ્રોકના સમૂહને ટેપની છેક ડાબે
લઈ જતી અવસ્થાઓ અને સૂચનાઓ ઉમેરવાની છે.
\olref{fig:mover}માંનું મશીન આ છેલ્લો ભાગ જ કરે છે:
$\TMblank$ના સમૂહ પછી $\TMstroke$નો સમૂહ હોય
એવી ટેપ પર (અને હેડ $\TMblank$ના સમૂહમાં ગમે ત્યાં
હોય) તેને શરૂ કરીએ, તો તે એક પછી એક $\TMstroke$
ભૂંસીને ટેપની શરૂઆતમાં લખે છે. કામ પૂરું થયું છે તે
ઓળખી શકાય એ માટે તે પહેલાં $\TMstroke$ના સમૂહનો
છેડો $\TMendtape$ પ્રતીક વડે ચિહ્નિત કરે છે; છેવટે
આ પ્રતીક ભૂંસાઈ જાય છે. આપણે અવસ્થાઓને~$q_6$થી
ક્રમ આપવાનું શરૂ કર્યું છે, જેથી તેને બમણું કરનાર મશીનમાં
ઉમેરી શકાય. તમારે વધારાની સૂચના
$\delta(q_0, \TMblank) = \tuple{q_6,\TMblank,\TMstay}$,
એટલે~$q_0$થી~$q_6$ સુધીનું
$\TMtrans{\TMblank}{\TMblank}{\TMstay}$ ચિહ્નવાળું તીર,
ઉમેરવાનું છે. (એક સૂક્ષ્મ સમસ્યા છે: મળેલું મશીન
ઇનપુટ~$x=0$ માટે કામ કરતું નથી. તેને આપણે અભ્યાસપ્રશ્ન
તરીકે રાખીએ છીએ.)
\begin{figure}
  \[
  \begin{tikzpicture}[->,>=stealth',shorten >=1pt,auto,node distance=2.8cm,
                      semithick]
    \tikzstyle{every state}=[fill=none,draw=black,text=black]
    \node[initial,state] (6)              {$q_6$};
    \node[state]         (7) [right of=6] {$q_7$};
    \node[state]         (8) [right of=7] {$q_8$};
    \node[state]         (9) [below of=8] {$q_9$};
    \node[state]         (10) [left of=9] {$q_{10}$};
    \node[state]         (11) [left of=10]  {$q_{11}$};
    \node[state]         (12) [below of=11]  {$q_{12}$};
    \node[state]         (13) [right of=12] {$q_{13}$};
    \node[state]         (14) [right of=13] {$q_{14}$};
    \path
    (6)  edge node {\TMtrans{\TMblank}{\TMblank}{\TMright}} (7)
    (7)  edge [loop above] node {\TMtrans{\TMblank}{\TMblank}{\TMright}} (7)
         edge node {\TMtrans{\TMstroke}{\TMstroke}{\TMright}} (8)
    (8)  edge [loop above] node {\TMtrans{\TMstroke}{\TMstroke}{\TMright}} (8)
         edge node {\TMtrans{\TMblank}{\TMendtape}{\TMleft}} (9)
    (9)  edge [loop right] node {\TMtrans{\TMstroke}{\TMstroke}{\TMleft}} (9)
         edge node {\TMtrans{\TMblank}{\TMblank}{\TMright}} (10)
    (10) edge node[above] {\TMtrans{\TMstroke}{\TMblank}{\TMleft}} (11)
    (11) edge [loop above] node {\TMtrans{\TMblank}{\TMblank}{\TMleft}} (11)
         edge node[left] {\begin{tabular}{@{}l@{}}
          \TMtrans{\TMendtape}{\TMendtape}{\TMright}\\
          \TMtrans{\TMstroke}{\TMstroke}{\TMright}
         \end{tabular}} (12)
    (12) edge node[] {\TMtrans{\TMblank}{\TMstroke}{\TMright}} (13)
    (13) edge [loop below] node {\TMtrans{\TMblank}{\TMblank}{\TMright}} (13)
         edge node[sloped] {\TMtrans{\TMstroke}{\TMblank}{\TMleft}} (11)
    (13) edge node {\TMtrans{\TMendtape}{\TMblank}{\TMstay}} (14);
  \end{tikzpicture}
  \]
  \caption{$\TMstroke$ના સમૂહને ડાબે ખસેડવું}
  \ollabel{fig:mover}
\end{figure}
\end{ex}

\begin{prob}
\olref[tur][mac][una]{ex:mover}માં આપણે
\olref[rep]{ex:doubler}માંના બમણું કરનાર મશીન અને
\olref[tur][mac][una]{fig:mover}માંના ખસેડનાર મશીનને
જોડીને બનેલા મશીનનું વર્ણન કર્યું. આ સંયુક્ત મશીનને
ઇનપુટ~$x=0$, એટલે ખાલી ટેપ, પર શરૂ કરીએ તો શું
થાય છે? આ કિસ્સામાં મશીન આઉટપુટ~$2x=0$ સાથે
થંભે એ માટે તેને કેવી રીતે સુધારશો? (એક અવસ્થા અને
એક સંક્રમણ ઉમેરીને આ કરી શકવું જોઈએ.)
\end{prob}

\sourcecorrection{OLTM-003}{સ્થિર અંગ્રેજી મૂળના આ પ્રશ્નમાં
બમણું કરનાર મશીન માટે $f(x)=2x$ આપતી
\olref[tur][mac][una]{fig:doubler-disc} આકૃતિનો સંદર્ભ છે.
પરંતુ તરત અગાઉનું વર્ણન મૂળ
\olref[rep]{ex:doubler} મશીનને
\olref[tur][mac][una]{fig:mover} સાથે જોડે છે; પ્રશ્નમાં
તે જ મૂળ મશીનનો સંદર્ભ આપ્યો છે.}

\begin{prob}
\emph{બાદબાકી:} $n$ અને~$m$ લંબાઈની બે ખાલી ન હોય
એવી સ્ટ્રોકની શબ્દમાળાઓ ઇનપુટ તરીકે મળે, જ્યાં
$n > m$ હોય, ત્યારે $f(n,m) = n - m$ વિધેય ગણતું
ટ્યુરિંગ મશીન રચો.
\end{prob}

\begin{prob}
\emph{સમાનતા:} નીચેનું વિધેય ગણતું ટ્યુરિંગ મશીન રચો:
\[
\fn{equality}(n,m) =
\begin{cases}
  \text{1} & \text{જો~$n = m$ હોય} \\
  \text{0} & \text{જો~$n \neq m$ હોય}
\end{cases}
\]
અહીં~$n$ અને~$m \in \PosInt$ છે.
\end{prob}

\begin{prob}
$x$ અને $y$ ધન પૂર્ણાંકોને ટેપ પર $\TMblank$થી
અલગ કરેલી $\TMstroke$ની શબ્દમાળાઓ વડે દર્શાવ્યા હોય,
ત્યારે $\min(x,y)$ વિધેય ગણતું ટ્યુરિંગ મશીન રચો.
મશીનની વર્ણમાળામાં વધારાનાં પ્રતીકો વાપરી શકો છો.

વિધેય $\min$ તેના આર્ગ્યુમેન્ટોમાંથી સૌથી નાનું મૂલ્ય
પસંદ કરે છે: તેથી $\min(3,5)=3$, $\min(20,16)=16$
અને $\min(4,4)=4$ વગેરે.
\end{prob}

\begin{defn}
  ટ્યુરિંગ મશીન~$M$ આંશિક વિધેય $f\colon \Nat^k
  \pto \Nat$ ગણે છે ત્યારે અને ત્યારે જ જ્યારે
  \begin{enumerate}
    \item $f(n_1, \dots, n_k) = m$ હોય તો $M$ ઇનપુટ
    $\TMstroke^{n_1}\concat\TMblank\concat
    \dots \concat\TMblank\concat\TMstroke^{n_k}$ પર
    આઉટપુટ $\TMstroke^{m}$ સાથે થંભે છે.
    \item $f(n_1, \dots, n_k)$ અવ્યાખ્યાયિત હોય તો
    $M$ થંભતું જ નથી, અથવા $\TMstroke$ના એક જ
    સમૂહથી ન બનેલા આઉટપુટ સાથે થંભે છે.
  \end{enumerate}
\end{defn}
% END OLP-0258
\endgroup


\begingroup
\makeatletter\def\ol@sectioncs{section}\makeatother

% BEGIN OLP-0259
\olfileid{tur}{mac}{hal}
\olsection{થંભવાની અવસ્થાઓ}
\phantomsection\label{guunit:OLP-0259}


\begin{explain}
આપણે મશીનને એવી રીતે વ્યાખ્યાયિત કર્યું છે કે અમલમાં મૂકવા
કોઈ સૂચના ન હોય ત્યારે જ તે થંભે. જોકે ટ્યુરિંગ મશીનનાં
સામાન્ય નિરૂપણોમાં એક ખાસ \emph{થંભવાની અવસ્થા}~$h$
હોય છે, જેમાં $h \in Q$.

થંભવાની અવસ્થા પાછળનો વિચાર સરળ છે: મશીનનું કામ પૂરું
થાય (તે ઇનપુટ સ્વીકારવા તૈયાર થાય અથવા આઉટપુટ
લખવાનું પૂરું કરે), ત્યારે તે અવસ્થા~$h$માં જઈને થંભે.
કેટલાંક મશીનોમાં થંભવાની બે અવસ્થાઓ હોય છે: એક
ઇનપુટ સ્વીકારવા માટે અને બીજી ઇનપુટ નકારવા માટે.
\end{explain}

\begin{ex}\emph{થંભવાની અવસ્થાઓ}.
આ વિચાર સ્પષ્ટ કરવા પહેલાં સમ મશીનમાં ફેરફાર કરીએ.
ઇનપુટ સમ હોય ત્યારે મશીન અવસ્થા~$q_0$માં થંભે તે
બદલે, તેને થંભવાની અવસ્થામાં મોકલતી સૂચના ઉમેરી
શકીએ છીએ.
\[
\begin{tikzpicture}[->,>=stealth',shorten >=1pt,auto,node distance=2.8cm,
                    semithick]
  \tikzstyle{every state}=[fill=none,draw=black,text=black]

  \node[initial, state]         (A)                     {$q_0$};
  \node[state]         (B) [right of=A] {$q_1$};
  \node[state]         (C) [below of=A] {$h$};

  \path (A) edge [bend left] node {\TMtrans{\TMstroke}{\TMstroke}{R}} (B)
      edge node {\TMtrans{\TMblank}{\TMblank}{N}} (C)
        (B) edge [loop above] node {\TMtrans{\TMblank}{\TMblank}{R}} (B)
            edge [bend left] node {\TMtrans{\TMstroke}{\TMstroke}{R}} (A);
\end{tikzpicture}
\]

હવે ઉદાહરણને વધુ વિસ્તૃત કરીએ. ઇનપુટ વિષમ છે એવું
મશીન નક્કી કરે ત્યારે તે ક્યારેય થંભતું નથી. આ
ચક્રમાં ચાલવાની સૂચનાને નકારવાની અવસ્થા~$r$માં જવાની
સૂચનાથી બદલીને મશીનમાં \emph{નકારવાની} અવસ્થા
ઉમેરી શકીએ છીએ.
\[
\begin{tikzpicture}[->,>=stealth',shorten >=1pt,auto,node distance=2.8cm,
                    semithick]
  \tikzstyle{every state}=[fill=none,draw=black,text=black]

  \node[initial,state]         (A)                     {$q_0$};
  \node[state]         (B) [right of=A] {$q_1$};
  \node[state]         (C) [below of=A] {$h$};
  \node[state]         (D) [below of=B] {$r$};

  \path (A) edge [bend left] node {\TMtrans{\TMstroke}{\TMstroke}{R}} (B)
      edge node {\TMtrans{\TMblank}{\TMblank}{N}} (C)
        (B) edge node {\TMtrans{\TMblank}{\TMblank}{N}} (D)
            edge [bend left] node {\TMtrans{\TMstroke}{\TMstroke}{R}} (A);
\end{tikzpicture}
\]
\end{ex}

\begin{explain}
આવા કિસ્સામાં ખાસ થંભવાની અવસ્થા ઉમેરવી લાભકારક છે:
મશીન કયા ઇનપુટ સ્વીકારે કે નકારે છે તે સ્પષ્ટ થાય છે.
પરંતુ ખાસ થંભવાની અવસ્થા ઉમેરવાથી મશીનની સંગણનશક્તિ
વધતી નથી એ નોંધવું જરૂરી છે. એવી જ રીતે, થંભવાનો
ઓછો ઔપચારિક ખ્યાલ પણ પોતાના લાભ ધરાવે છે. આ
પ્રકરણમાં અત્યાર સુધી વપરાયેલી થંભવાની વ્યાખ્યા
\emph{અટકવાની સમસ્યા}ની સાબિતી સહજ અને રજૂ કરવામાં
સરળ બનાવે છે. તેથી આપણે આપણી મૂળ વ્યાખ્યા જ
ચાલુ રાખીએ છીએ.
\end{explain}
% END OLP-0259
\endgroup


\begingroup
\makeatletter\def\ol@sectioncs{section}\makeatother

% BEGIN OLP-0260
\olfileid{tur}{mac}{dis}
\olsection{શિસ્તબદ્ધ મશીનો}
\phantomsection\label{guunit:OLP-0260}


\begin{explain}
\olref[hal]{sec} વિભાગમાં આપણે એવી ટ્યુરિંગ મશીનો
જોઈ હતી જેમાં એક જ નિર્ધારિત થંભવાની અવસ્થા~$h$
હોય છે---આવાં મશીનો થંભે તો અવસ્થા~$h$માં જ થંભે
તેની ખાતરી હોય છે. આ રીતે, એક જ થંભવાની અવસ્થાવાળાં
મશીનો સામાન્ય ટ્યુરિંગ મશીનોને આપણે આપેલી છૂટ કરતાં
વધુ ``શિસ્તબદ્ધ'' છે. ટ્યુરિંગ મશીનોના વર્તન પર બીજી
શરતો પણ મૂકી શકીએ. દાખલા તરીકે, મશીન ખાનું~$0$માંનું
ટેપના છેડાનું ચિહ્ન ક્યારેય ભૂંસે નહીં અથવા ખાનું~$0$માંથી
ડાબે જવાનો પ્રયત્ન ન કરે એવી મનાઈ આપણે મૂકી નથી.
(આવા કિસ્સામાં હેડ ખાનું~$0$માં જ રહે છે એમ આપણી
વ્યાખ્યા કહે છે; બીજી વ્યાખ્યાઓમાં મશીન થંભે છે.)
એવી જ રીતે, ટ્યુરિંગ મશીન થંભે તો ખાનું~$1$ વાંચતાં
થંભે એવી ધારણા કરી શકાય તે ક્યારેક ઇચ્છનીય છે.
\end{explain}

\begin{defn}\ollabel{defn:disciplined}
ટ્યુરિંગ મશીન~$M$ \emph{શિસ્તબદ્ધ} છે ત્યારે અને
ત્યારે જ જ્યારે
\begin{enumerate}
    \item તેમાં એક જ નિર્ધારિત થંભવાની અવસ્થા~$h$ છે,
    \item તે થંભે તો ખાનું~$1$ વાંચતાં થંભે છે,
    \item તે ખાનું~$0$માંનું $\TMendtape$ પ્રતીક ક્યારેય
    ભૂંસતું નથી, અને
    \item તે ખાનું~$0$માંથી ડાબે જવાનો પ્રયત્ન ક્યારેય
    કરતું નથી.
\end{enumerate}
\end{defn}

\begin{explain}
આપણે અગાઉ ચર્ચ્યું છે કે કોઈ પણ ટ્યુરિંગ મશીનને એ જ
વર્તન ધરાવતું પરંતુ નિર્ધારિત થંભવાની અવસ્થાવાળું મશીન
બનાવી શકાય છે. આ માટે માત્ર નવી અવસ્થા~$h$ ઉમેરીને
મૂળ $\delta$ જે જોડી $\tuple{q,\sigma}$ પર અવ્યાખ્યાયિત
હોય તે દરેક માટે સૂચના
$\delta(q, \sigma) = \tuple{h, \sigma, N}$ ઉમેરવી પડે.
સાબિતી કંટાળાજનક હોવા છતાં એ સાચું છે કે કોઈ પણ
ટ્યુરિંગ મશીન~$M$ને એવા શિસ્તબદ્ધ ટ્યુરિંગ મશીન~$M'$
માં ફેરવી શકાય છે જે એ જ ઇનપુટ પર થંભે છે અને એ જ
આઉટપુટ આપે છે. દાખલા તરીકે, ટ્યુરિંગ મશીન થંભે અને
તે ખાનું~$1$ પર ન હોય, તો હેડ ટેપના છેડાનું ચિહ્ન
મળે ત્યાં સુધી ડાબે ખસે, પછી એક ખાનું જમણે ખસે અને
પછી થંભે એવી સૂચનાઓ ઉમેરી શકાય. બીજી શરતો કેવી
રીતે પૂરી કરી શકાય તે તમે વિચારો.
\end{explain}

\begin{ex}
    \olref{fig:adder-disc}માં આપણે
    \olref[una]{ex:adder}ના સરવાળાના મશીનને શિસ્તબદ્ધ
    મશીનમાં ફેરવીએ છીએ.
    \begin{figure}\[
\begin{tikzpicture}[->,>=stealth',shorten >=1pt,auto,node distance=2.8cm,
                    semithick]
  \tikzstyle{every state}=[fill=none,draw=black,text=black]

  \node[initial,state]         (A)              {$q_0$};
  \node[state]         (B) [right of=A] {$q_1$};
  \node[state]         (C) [below right of=B] {$q_2$};
  \node[state]         (D) [below left of=C] {$q_3$};
  \node[state]         (H) [left of=D] {$h$};

  \path (A) edge node {\TMtrans{\TMblank}{\TMstroke}{\TMstay}} (B)
           edge [loop below] node {\TMtrans{\TMstroke}{\TMstroke}{\TMright}} (B)
        (B) edge [loop below] node {\TMtrans{\TMstroke}{\TMstroke}{\TMright}} (B)
            edge node {\TMtrans{\TMblank}{\TMblank}{\TMleft}} (C)
        (C) edge node {\TMtrans{\TMstroke}{\TMblank}{\TMleft}} (D)
        (D) edge [loop above] node {\TMtrans{\TMstroke}{\TMstroke}{\TMleft}} (D)
        edge node {\TMtrans{\TMendtape}{\TMendtape}{\TMright}} (H);
\end{tikzpicture}
\]\caption{શિસ્તબદ્ધ સરવાળાનું મશીન}
\ollabel{fig:adder-disc}
\end{figure}
\end{ex}

\begin{prop}\ollabel{prop:disciplined} દરેક ટ્યુરિંગ
મશીન~$M$ માટે એવું શિસ્તબદ્ધ ટ્યુરિંગ મશીન~$M'$
છે કે $M$ આઉટપુટ~$O$ સાથે થંભે તો તે પણ
આઉટપુટ~$O$ સાથે થંભે, અને $M$ ન થંભે તો તે
પણ ન થંભે. ખાસ કરીને, ટ્યુરિંગ મશીન વડે સંગણનીય
દરેક વિધેય $f\colon\Nat^n \to \Nat$ શિસ્તબદ્ધ
ટ્યુરિંગ મશીન વડે પણ સંગણનીય છે.
\end{prop}

\begin{prob}\label{tur:mac:dis:prob:disc-succ}
$f(x) = x+1$ ગણતું શિસ્તબદ્ધ મશીન આપો.
\end{prob}


\begin{prob}\label{tur:mac:dis:prob:copier}
ઇનપુટ $\TMstroke^n$ પર શરૂ કરતાં આઉટપુટ
$\TMstroke^n \concat \TMblank \concat \TMstroke^n$
આપે એવું શિસ્તબદ્ધ મશીન શોધો.
\end{prob}
% END OLP-0260
\endgroup


\begingroup
\makeatletter\def\ol@sectioncs{section}\makeatother

% BEGIN OLP-0261
\olfileid{tur}{mac}{cmb}
\olsection{ટ્યુરિંગ મશીનોનું સંયોજન}
\phantomsection\label{guunit:OLP-0261}


\begin{explain}
અત્યાર સુધી આપણે જોયેલાં ટ્યુરિંગ મશીનો સ્વરૂપે ઘણાં
સરળ હતાં. પરંતુ હકીકતમાં, કોઈ પણ આધુનિક પ્રોગ્રામિંગ
ભાષામાં ઉકેલી શકાય તેવી સમસ્યા ટ્યુરિંગ મશીનોથી પણ
ઉકેલી શકાય છે. વધુ જટિલ ટ્યુરિંગ મશીનો બનાવવા માટે
તેમને જોડવાનું શક્ય છે તે સમજવું જરૂરી છે. એમ હોય તો
પ્રક્રિયાને સરળ ભાગોમાં વહેંચીને વધુ જટિલ સમસ્યાઓ
ઉકેલતાં મશીનો બનાવી શકીએ. જટિલ સમસ્યાને તેના
ઘટક ભાગોમાં વહેંચવાની સ્વાભાવિક રીત મળે, તો અનેક
સરળ ટ્યુરિંગ મશીનો બનાવીને તેમને એક મશીનમાં જોડતાં
સમસ્યાને કેટલાંક તબક્કામાં ઉકેલી શકીએ. આગળના વિભાગમાં
અટકવાની સમસ્યા હાથ ધરતી વખતે આ વાત ખાસ અગત્યની છે.

ટ્યુરિંગ મશીનો $M = \tuple{Q, \Sigma, q_0, \delta}$
અને~$M' = \tuple{Q', \Sigma', q_0', \delta'}$ને કેવી રીતે
જોડીએ? હવે $M$ થંભે તે પછી ટેપનું સંરૂપણ મશીન~$M'$
ના ચલનના ઇનપુટ સંરૂપણ તરીકે વાપરીએ. આ કામ કરતું
એક ટ્યુરિંગ મશીન $M \frown M'$ મેળવવા માટે આવું કરો:
\begin{enumerate}
    \item $M'$ની બધી અવસ્થાઓ~$Q'$ના ક્રમાંક (અથવા નામ)
    ફરી ગોઠવો, જેથી $M$ અને~$M'$ની કોઈ અવસ્થા
    સમાન ન હોય ($Q \cap Q' = \emptyset$).
    \item $M \frown M'$ની અવસ્થાઓ $Q \cup Q'$ છે.
    \item ટેપની વર્ણમાળા $\Sigma \cup \Sigma'$ છે.
    \item પ્રારંભિક અવસ્થા~$q_0$ છે.
    \item સંક્રમણ વિધેય નીચે આપેલું વિધેય $\delta''$ છે:
    \[\delta''(q,\sigma) =
    \begin{cases}
      \delta(q,\sigma) & \text{જો $q \in Q$ અને $\delta(q,\sigma)$ વ્યાખ્યાયિત હોય}\\
      \delta'(q,\sigma) & \text{જો $q \in Q'$ હોય}\\
      \tuple{q_0', \sigma, \TMstay} & \text{જો $q \in Q$ હોય અને
      $\delta(q,\sigma)$ અવ્યાખ્યાયિત હોય}
    \end{cases}\]
\end{enumerate}
મળેલું મશીન અવસ્થા $q \in Q$માં હોય ત્યારે~$M$ની
સૂચનાઓ અને અવસ્થા~$q \in Q'$માં હોય ત્યારે~$M'$
ની સૂચનાઓ વાપરે છે. તે અવસ્થા $q \in Q$માં હોય
અને એવું પ્રતીક~$\sigma$ વાંચે કે જેના માટે $M$માં
સંક્રમણ નથી (એટલે $M$ થંભ્યું હોત), ત્યારે તે~$M'$
ની પ્રારંભિક અવસ્થામાં જાય છે (અને ટેપની સામગ્રી તથા
હેડનું સ્થાન યથાવત્ રાખે છે).
\sourcecorrection{OLTM-004}{સ્થિર અંગ્રેજી મૂળના પ્રથમ
કિસ્સામાં માત્ર $q\in Q$ લખ્યું છે, તેથી $\delta(q,\sigma)$
અવ્યાખ્યાયિત હોય ત્યારે પણ પહેલો કિસ્સો લાગુ પડે અને
ત્રીજો કિસ્સો ઢંકાઈ જાય. ગુજરાતી સૂત્રના પહેલા કિસ્સામાં
$\delta(q,\sigma)$ વ્યાખ્યાયિત હોવાની શરત સ્પષ્ટ કરી છે;
થંભેલા $M$માંથી $M'$માં જવાનું બાજુનું વર્ણન આ અર્થ
નક્કી કરે છે.}

નોંધો કે મશીન~$M$ શિસ્તબદ્ધ ન હોય, તો $M$ થંભે
ત્યારે ટેપનો હેડ ક્યાં હશે તે આપણે જાણતા નથી; તેથી
$M$ના થંભવાના સંરૂપણમાં હેડ ખાનું~$1$ વાંચતો હોય
એવું જરૂરી નથી. મશીનો જોડતી વખતે આ વાત ધ્યાનમાં
રાખવી અગત્યની છે.
\end{explain}

\begin{ex}
\emph{મશીનોનું સંયોજન:} એવું મશીન રચીશું કે તેને ઇનપુટ
તરીકે $n$ અને~$m$ લંબાઈના $\TMstroke$ના બે સમૂહો
આપીએ, તો તે ટેપ પર $2(m+n)$ $\TMstroke$ના એક
જ સમૂહ સાથે થંભે. તેને બનાવવા આપણે પહેલેથી જાણીએ
છીએ એવાં બે મશીનો---સરવાળાનું મશીન અને બમણું કરનાર
મશીન---જોડી શકીએ. પહેલાં સરવાળાના મશીનનો અવસ્થા-આલેખ
દોરીએ.
\[
\begin{tikzpicture}[->,>=stealth',shorten >=1pt,auto,node distance=2.8cm,
                    semithick]
  \tikzstyle{every state}=[fill=none,draw=black,text=black]

  \node[initial,state] (A)              {$q_0$};
  \node[state]         (B) [right of=A] {$q_1$};
  \node[state]         (C) [right of=B] {$q_2$};

  \path (A) edge node {\TMtrans{\TMblank}{\TMstroke}{\TMstay}} (B)
            edge [loop above] node {\TMtrans{\TMstroke}{\TMstroke}{\TMright}} (B)
        (B) edge [loop above] node {\TMtrans{\TMstroke}{\TMstroke}{\TMright}} (B)
            edge node {\TMtrans{\TMblank}{\TMblank}{\TMleft}} (C)
        (C) edge [loop above] node {\TMtrans{\TMstroke}{\TMblank}{\TMstay}} (C);
\end{tikzpicture}
\]
અવસ્થા~$q_2$માં થંભવાને બદલે, આઉટપુટ બમણું કરવા
મશીન આગળ ચાલે એવું જોઈએ. યાદ કરો કે બમણું કરનાર
મશીન ઇનપુટનો પ્રથમ સ્ટ્રોક ભૂંસીને અલગ આઉટપુટમાં
બે સ્ટ્રોક લખે છે. ટેપનો હેડ સરવાળાના મશીનના
આઉટપુટનો પ્રથમ સ્ટ્રોક વાંચે તેની ખાતરી કરતી સૂચના
ઉમેરીએ.
\[
\begin{tikzpicture}[->,>=stealth',shorten >=1pt,auto,node distance=2.8cm,
                    semithick]
  \tikzstyle{every state}=[fill=none,draw=black,text=black]

  \node[initial,state] (A)              {$q_0$};
  \node[state]         (B) [right of=A] {$q_1$};
  \node[state]         (C) [right of=B] {$q_2$};
  \node[state]         (D) [below left of=C] {$q_3$};
  \node[state]         (E) [below left of=D] {$q_4$};

  \path (A) edge node {\TMtrans{\TMblank}{\TMstroke}{\TMstay}} (B)
            edge [loop above] node {\TMtrans{\TMstroke}{\TMstroke}{\TMright}} (B)
        (B) edge [loop above] node {\TMtrans{\TMstroke}{\TMstroke}{\TMright}} (B)
            edge node {\TMtrans{\TMblank}{\TMblank}{\TMleft}} (C)
        (C) edge node {\TMtrans{\TMstroke}{\TMblank}{\TMleft}} (D)
        (D) edge [loop left] node {\TMtrans{\TMstroke}{\TMstroke}{\TMleft}} (D)
            edge node[left, xshift=-2mm] {\TMtrans{\TMendtape}{\TMendtape}{\TMright}} (E);
\end{tikzpicture}
\]
હવે ઇનપુટ બમણું કરવું સહેલું છે---બમણું કરનાર મશીનને
અવસ્થા~$q_4$ સાથે જોડવાનું જ બાકી છે. એ માટે
બમણું કરનાર મશીનની અવસ્થાઓનાં નામ બદલવાં પડે,
જેથી તે~$q_0$ને બદલે~$q_4$થી શરૂ થાય---આ રીતે
બે પ્રારંભિક અવસ્થાઓ ઊભી થતી નથી. આખરી આલેખ
\olref{fig:combined} જેવો હોવો જોઈએ.
\begin{figure}
\[
\begin{tikzpicture}[->,>=stealth',shorten >=1pt,auto,node distance=2.8cm,
                    semithick]
  \tikzstyle{every state}=[fill=none,draw=black,text=black]
  \node[initial,state] (A)              {$q_0$};
  \node[state]         (B) [right of=A] {$q_1$};
  \node[state]         (C) [right of=B] {$q_2$};
  \node[state]         (D) [below left of=C] {$q_3$};
  \node[state]         (E) [below left of=D] {$q_4$};
  \node[state]         (2) [right of=E] {$q_5$};
  \node[state]         (3) [right of=2] {$q_6$};
  \node[state]         (4) [below of=3] {$q_7$};
  \node[state]         (5) [left of=4]  {$q_8$};
  \node[state]         (6) [left of=5]  {$q_9$};

  \path (A) edge node {\TMtrans{\TMblank}{\TMstroke}{\TMstay}} (B)
            edge [loop above] node {\TMtrans{\TMstroke}{\TMstroke}{\TMright}} (B)
        (B) edge [loop above] node {\TMtrans{\TMstroke}{\TMstroke}{\TMright}} (B)
            edge node {\TMtrans{\TMblank}{\TMblank}{\TMleft}} (C)
        (C) edge node {\TMtrans{\TMstroke}{\TMblank}{\TMleft}} (D)
        (D) edge [loop left] node {\TMtrans{\TMstroke}{\TMstroke}{\TMleft}} (D)
            edge node[left, xshift=-2mm] {\TMtrans{\TMendtape}{\TMendtape}{\TMright}} (E)
    (E) edge node {\TMtrans{\TMstroke}{\TMblank}{\TMright}} (2)
    (2) edge [loop above] node {\TMtrans{\TMstroke}{\TMstroke}{\TMright}} (2)
      edge node {\TMtrans{\TMblank}{\TMblank}{\TMright}} (3)
    (3) edge [loop above] node {\TMtrans{\TMstroke}{\TMstroke}{\TMright}} (3)
        edge node {\TMtrans{\TMblank}{\TMstroke}{\TMright}} (4)
    (4) edge [loop below] node {\TMtrans{\TMblank}{\TMstroke}{\TMleft}} (4)
        edge node {\TMtrans{\TMstroke}{\TMstroke}{\TMleft}} (5)
    (5) edge [loop below]  node {\TMtrans{\TMstroke}{\TMstroke}{\TMleft}} (5)
        edge              node {\TMtrans{\TMblank}{\TMblank}{\TMleft}} (6)
    (6) edge [loop below] node {\TMtrans{\TMstroke}{\TMstroke}{\TMleft}} (6)
        edge              node {\TMtrans{\TMblank}{\TMblank}{\TMright}} (E);
\end{tikzpicture}
\]
\caption{સરવાળાનું અને બમણું કરનાર મશીન જોડવાં}
\ollabel{fig:combined}
\end{figure}
\end{ex}

\begin{prop}
    $M$ અને $M'$ શિસ્તબદ્ધ હોય અને અનુક્રમે
    $f\colon \Nat^k \to \Nat$ અને
    $f'\colon \Nat \to \Nat$ વિધેયો ગણતા હોય, તો
    $M \frown M'$ પણ શિસ્તબદ્ધ છે અને
    $\comp{f}{f'}$ ગણે છે.
\end{prop}

\begin{proof}
    $M$ શિસ્તબદ્ધ હોવાથી તે
    આઉટપુટ~$f(n_1,\dots,n_k) = m$ સાથે થંભે ત્યારે
    હેડ ખાનું~$1$ વાંચે છે. હવે~$M'$ની પ્રારંભિક
    અવસ્થામાં જઈએ, તો મશીન આઉટપુટ $f'(m)$ સાથે,
    ફરી ખાનું~$1$ વાંચતાં, થંભશે.
    \olref[dis]{defn:disciplined}ની બીજી શરતો પણ
    પૂરી થાય છે.
\end{proof}

\begin{prob}
    \cref{tur:mac:dis:prob:disc-succ}નું મશીન~$M$
    લઈને અને $M \frown M$ રચીને $f(x) = x+2$
    ગણતું શિસ્તબદ્ધ ટ્યુરિંગ મશીન આપો.
\end{prob}
% END OLP-0261
\endgroup


\begingroup
\makeatletter\def\ol@sectioncs{section}\makeatother

% BEGIN OLP-0262
\olfileid{tur}{mac}{var}
\olsection{ટ્યુરિંગ મશીનોનાં વિવિધ સ્વરૂપો}
\phantomsection\label{guunit:OLP-0262}


હકીકતમાં ટ્યુરિંગ મશીનોની વ્યાખ્યા આપવાની ઘણી રીતો છે;
આપણી રીત માત્ર એક છે. કેટલીક બાબતમાં આપણી વ્યાખ્યા
બીજી વ્યાખ્યાઓ કરતાં વધુ છૂટ આપે છે. આપણે મનસ્વી
સાન્ત વર્ણમાળાઓને મંજૂરી આપીએ છીએ; વધુ મર્યાદિત
વ્યાખ્યા ટેપનાં માત્ર બે પ્રતીકો, $\TMstroke$ અને~$\TMblank$,
મંજૂર કરી શકે. આપણે મશીનને ટેપ પર પ્રતીક લખતાં જ
ખસવાની છૂટ આપીએ છીએ; બીજી વ્યાખ્યાઓમાં લખવું
અથવા ખસવું, બેમાંથી એક જ થઈ શકે. આપણે ટેપના હેડને
ખસેડ્યા વગર લખવાની શક્યતા મંજૂર કરીએ છીએ; બીજી
વ્યાખ્યાઓ $\TMstay$ ``સૂચના'' છોડે છે. બીજી બાબતમાં
આપણી વ્યાખ્યા વધુ મર્યાદિત છે. આપણે ટેપ માત્ર એક
દિશામાં અનંત છે એમ ધાર્યું છે; બીજી વ્યાખ્યાઓ તેને ડાબે
અને જમણે બંને તરફ અનંત રાખે છે. હકીકતમાં જુદી જુદી
ગમે તેટલી ટેપ, અથવા ખાનાંની અનંત જાળી પણ મંજૂર કરી
શકાય. આપણે ટ્યુરિંગ મશીનનો સૂચનાઓનો ગણ સંક્રમણ
વિધેયથી દર્શાવીએ છીએ; બીજી વ્યાખ્યાઓ સંક્રમણ સંબંધ
વાપરે છે, જેમાં કોઈ પણ આપેલી સ્થિતિમાં મશીન પાસે
એકથી વધુ શક્ય સૂચનાઓ હોય છે.

વ્યાખ્યામાં કરેલી આ છેલ્લી છૂટ ખાસ રસપ્રદ છે. આપણી
વ્યાખ્યામાં મશીન અવસ્થા~$q$માં પ્રતીક~$\sigma$ વાંચે,
ત્યારે $\delta(q, \sigma)$ નવું પ્રતીક, નવી અવસ્થા અને
ટેપના હેડનું નવું સ્થાન નક્કી કરે છે. પરંતુ સૂચનાઓનો
ગણ વર્તમાન અવસ્થા-પ્રતીક જોડીઓ $\tuple{q,
  \sigma}$ અને નવી અવસ્થા-પ્રતીક-દિશા ત્રિકો
$\tuple{q', \sigma',
  D}$ વચ્ચેનો સંબંધ હોય, તો ટ્યુરિંગ મશીનની ક્રિયા
એક જ રીતે નક્કી થાય એવું જરૂરી નથી---સૂચનાઓના
સંબંધમાં $\tuple{q,
  \sigma, q', \sigma', D}$ અને
$\tuple{q, \sigma, q'', \sigma'', D'}$ બંને હોઈ શકે.
આ કિસ્સામાં આપણી પાસે \emph{અનિર્ધારિત} ટ્યુરિંગ મશીન
છે. સંગણકીય જટિલતાના સિદ્ધાંતમાં આવાં મશીનો અગત્યની
ભૂમિકા ભજવે છે.

ટ્યુરિંગ મશીન ક્યારે થંભે તેની પણ જુદી રીતો છે: આપણે
સંક્રમણ વિધેય અવ્યાખ્યાયિત હોય ત્યારે તે થંભે એમ કહીએ
છીએ; બીજી વ્યાખ્યાઓમાં મશીન ખાસ નિર્ધારિત થંભવાની
અવસ્થામાં હોય તે જરૂરી છે. આવી અવસ્થા જરૂરી રાખવાથી
ટ્યુરિંગ મશીનો શું ગણી શકે તેમાં ફેર પડતો નથી તે આપણે
\olref[hal]{sec}માં સમજાવ્યું છે. આપણા ટ્યુરિંગ મશીનોની
ટેપ માત્ર એક દિશામાં અનંત હોવાથી ક્યારેક મશીન સૂચના
યોગ્ય રીતે ચલાવી શકતું નથી: સૌથી ડાબું ખાનું વાંચતું
હોય અને ડાબે ખસવાનું કહેવામાં આવે ત્યારે. આપણી
વ્યાખ્યા મુજબ ``નીચે પડી જવા''ને બદલે તે એ જ સ્થાને
રહે છે; પરંતુ આવું બને ત્યારે મશીન થંભે એવી વ્યાખ્યા
પણ આપી શકીએ. આ વ્યાખ્યા પણ સમકક્ષ છે: ખાનું~$0$માંથી
ડાબે ખસવાનો પ્રયત્ન કરતાં થંભતું ટ્યુરિંગ મશીનનું વર્તન
અનુકરવા, દરેક સંક્રમણ $\delta(q,\TMendtape) =
\tuple{q',\sigma,\TMleft}$ દૂર કરી શકીએ---પછી
મશીન~$\TMendtape$ પર ડાબે ખસવાનો પ્રયત્ન કરવાને
બદલે થંભે છે.\footnote{આ રીત \emph{પૂરેપૂરી} કામ
કરતી નથી, કારણ કે ખાનું~$0$ સિવાયનાં ખાનાંમાં પણ
$\TMendtape$ લખતાં અને વાંચતાં આપણને કોઈ રોકતું નથી
(જુઓ \olref[una]{ex:mover}). આવાં કામ માટે બીજું
પ્રતીક~$\TMendtape'$ ઉમેરીને આ મુશ્કેલી દૂર કરી શકીએ.}

સંખ્યાઓ દર્શાવવાની પણ જુદી રીતો છે (અને તેથી ટ્યુરિંગ
મશીન ગણતું ઇનપુટ-આઉટપુટ વિધેય પણ બદલાય છે): આપણે
એક-પ્રતીકી નિરૂપણ વાપરીએ છીએ, પરંતુ દ્વિ-આધારી
નિરૂપણ પણ વાપરી શકાય. તે માટે $\TMblank$ અને~$\TMendtape$
ઉપરાંત બે પ્રતીકો જોઈએ.

હવે એક રસપ્રદ હકીકત: કયાં વિધેયો ટ્યુરિંગ-સંગણનીય છે
તેના પર આમાંથી કોઈ ફેરફારનો પ્રભાવ પડતો નથી.
\emph{એક વ્યાખ્યા અનુસાર કોઈ વિધેય ટ્યુરિંગ-સંગણનીય
હોય, તો બધી વ્યાખ્યાઓ અનુસાર તે ટ્યુરિંગ-સંગણનીય છે.}

આ વાતની ચકાસણીની વિગતોમાં નહીં જઈએ. માત્ર એક
ઉદાહરણ જોઈએ: ટેપ બંને દિશામાં અનંત હોય અથવા
અનેક ટેપ હોય એવી છૂટ આપવાથી સંગણનશક્તિ વધતી
નથી. લગભગ આવું કારણ છે: એક જ, એક દિશામાં
અનંત ટેપવાળું ટ્યુરિંગ મશીન અનેક ટેપનું અથવા બે
દિશામાં અનંત ટેપનું અનુકરણ કરી શકે છે. દાખલા તરીકે,
વધારાની અવસ્થાઓ અને સૂચનાઓ વાપરીને અનેક ટેપવાળા
અથવા બે દિશામાં અનંત ટેપવાળા મશીનના પ્રોગ્રામને
એક જ, એક દિશામાં અનંત ટેપવાળા મશીનના પ્રોગ્રામમાં
``રૂપાંતરિત'' કરી શકીએ. રૂપાંતરિત મશીન સમ ક્રમાંકનાં
ખાનાં ટેપ~$1$નાં ખાનાં (અથવા બે દિશામાં અનંત ટેપનાં
``ધન'' ખાનાં) માટે અને વિષમ ક્રમાંકનાં ખાનાં ટેપ~$2$
નાં ખાનાં (અથવા ``ઋણ'' ખાનાં) માટે વાપરી શકે છે.
% END OLP-0262
\endgroup


\begingroup
\makeatletter\def\ol@sectioncs{section}\makeatother

% BEGIN OLP-0263
\olfileid{tur}{mac}{ctt}
\olsection{ચર્ચ--ટ્યુરિંગ માન્યતા}
\phantomsection\label{guunit:OLP-0263}


ટ્યુરિંગ મશીનો અસરકારક પ્રક્રિયાના ખ્યાલનું ચોક્કસ
ગણિતીય સ્વરૂપ આપવાનાં છે. ટ્યુરિંગના મતે, જે કોઈ
અસરકારક પ્રક્રિયા અને ટ્યુરિંગ મશીન બંનેના ખ્યાલને
સમજે તેને સહજ રીતે લાગશે કે અસરકારક પ્રક્રિયા વડે
જે કંઈ કરી શકાય તે ટ્યુરિંગ મશીન વડે પણ કરી શકાય.
આ દાવાને એ હકીકત ટેકો આપે છે કે અસરકારક પ્રક્રિયાના
ખ્યાલનાં અન્ય તમામ પ્રસ્તાવિત ચોક્કસ સ્વરૂપો ટ્યુરિંગ
મશીનના ખ્યાલ સાથે વિસ્તારની દૃષ્ટિએ સમકક્ષ નીવડે
છે---એટલે તેઓ બરાબર એ જ વિધેયોનો ગણ ગણી શકે છે.
આ દાવો \emph{ચર્ચ--ટ્યુરિંગ માન્યતા} કહેવાય છે.

\begin{defn}[ચર્ચ--ટ્યુરિંગ માન્યતા]
\emph{ચર્ચ--ટ્યુરિંગ માન્યતા} કહે છે કે અસરકારક
પ્રક્રિયા વડે સંગણનીય હોય તે બધું ટ્યુરિંગ-સંગણનીય છે.
\end{defn}

ચર્ચ--ટ્યુરિંગ માન્યતાનો બે રીતે આધાર લેવાય છે. પ્રથમ
રીત આળસને બહાનું આપે છે. ધારો કે આપણને કંઈક ગણવાની
અસરકારક પ્રક્રિયાનું વર્ણન, માનો ``આભાસી કોડ''માં,
મળ્યું છે. તો આપણે ખરેખર ટ્યુરિંગ મશીન બનાવ્યું ન
હોય તોપણ એ જ વિધેય કોઈ ટ્યુરિંગ મશીન ગણે છે એવો
દાવો યોગ્ય ઠેરવવા ચર્ચ--ટ્યુરિંગ માન્યતા લઈ શકીએ.

ચર્ચ--ટ્યુરિંગ માન્યતાનો બીજો ઉપયોગ તત્ત્વચિંતનની
દૃષ્ટિએ વધુ રસપ્રદ છે. કેટલાંક વિધેયો ટ્યુરિંગ મશીન
વડે ગણી શકાતા નથી તે બતાવી શકાય છે. તેના પરથી,
ચર્ચ--ટ્યુરિંગ માન્યતા વાપરીને, કોઈ પણ પ્રકારની પ્રક્રિયા
વડે તેમને અસરકારક રીતે ગણી શકાતા નથી એવો નિષ્કર્ષ
કાઢી શકાય છે. કારણ કે એવી પ્રક્રિયા હોય, તો
ચર્ચ--ટ્યુરિંગ માન્યતા મુજબ તેને માટે ટ્યુરિંગ મશીન
હોવું જોઈએ. તેથી જો તેને ગણતું ટ્યુરિંગ મશીન નથી
એવું સાબિત કરી શકીએ, તો અસરકારક પ્રક્રિયા પણ
હોઈ શકે નહીં. ખાસ કરીને, કહેવાતી અટકવાની સમસ્યા
ટ્યુરિંગ મશીનોથી જ નહીં પરંતુ કોઈ પણ અસરકારક રીતે
ઉકેલી શકાતી નથી એવો દાવો કરવા ચર્ચ--ટ્યુરિંગ
માન્યતા વપરાય છે.
% END OLP-0263
\endgroup


\OLEndChapterHook
% END OLP-0253
\endgroup


\begingroup
\makeatletter\def\ol@sectioncs{section}\makeatother

% BEGIN OLP-0264
\olchapter{tur}{und}{અનિર્ણેયતા}
\olchapterid{tur}{und}
\phantomsection\label{guunit:OLP-0264}


\begingroup
\makeatletter\def\ol@sectioncs{section}\makeatother

% BEGIN OLP-0265
\olfileid{tur}{und}{int}
\olsection{પરિચય}
\phantomsection\label{guunit:OLP-0265}


પ્રથમ નજરે એવું સ્પષ્ટ લાગશે કે કેટલાક વિધેયો, અંકગણિતીય
વિધેયો પણ, અસંગણનીય છે. આવાં વિધેયો એટલાં બધાં છે
અને તેમનું વર્તન એટલું
જટિલ છે. દાખલા તરીકે, કિરણોત્સર્ગી કણોના ક્ષય પરથી, અથવા
બીજા અવ્યવસ્થિત કે યાદૃચ્છિક વર્તન પરથી વ્યાખ્યાયિત વિધેયો લો.
ધારો કે આપણે કોઈ આપેલી ક્ષણથી એક-સેકંડના સમયગાળાની ગણતરી
શરૂ કરીએ અને તે પ્રારંભિક ક્ષણ પછીના $n$મા એક-સેકંડના
સમયગાળામાં બ્રહ્માંડના ક્ષય પામતા કણોની સંખ્યા $f(n)$ તરીકે
વ્યાખ્યાયિત કરીએ. આ વિધેય કદાચ એવું લાગે કે તેને કદી
ગણવાની આશા રાખી શકાય નહીં.

પરંતુ આવા વિધેયો કેવી રીતે ગણવા તેની કલ્પના ન થઈ શકે એ એક
વાત છે; તેઓ અસંગણનીય છે તે ખરેખર સાબિત કરવું બીજી વાત છે.
હકીકતમાં, અત્યંત જટિલ લાગતાં વિધેયો પણ અમૂર્ત અર્થમાં
સંગણનીય હોઈ શકે. ઉદાહરણ તરીકે, ધારો કે બ્રહ્માંડ સમયની
દૃષ્ટિએ સાન્ત છે---કેટલાક બ્રહ્માંડવિજ્ઞાનના સિદ્ધાંતોની
આગાહી મુજબ, દૂરના ભવિષ્યમાં કોઈ દિવસે તે સંકોચાઈને એક જ
બિંદુ બની જશે. તો એ પ્રારંભિક ક્ષણ પછી $f(n)$ વ્યાખ્યાયિત
હોય એવા સેકંડોની સંખ્યા સાન્ત (જોકે અતિ વિશાળ) હશે.
માત્ર સાન્ત સંખ્યામાં ઇનપુટ માટે વ્યાખ્યાયિત કોઈ પણ વિધેય
સંગણનીય છે: આપણે તેના આઉટપુટ એક મોટા કોષ્ટકમાં લખી શકીએ,
અથવા તેમને અત્યંત મોટા ટ્યુરિંગ મશીનના અવસ્થા-આલેખમાં
સંકેતિત કરી શકીએ.

ઘણી વાર આપણને એવા વિધેયોના ખાસ કિસ્સાઓમાં રસ હોય છે
જેનાં મૂલ્યો હા/ના પ્રશ્નોના જવાબ આપે છે. ઉદાહરણ તરીકે,
``$n$ અવિભાજ્ય સંખ્યા છે?'' એ પ્રશ્નની સાથે નીચેનું વિધેય જોડાય છે:
\[
\fn{isprime}(n) = \begin{cases}
  1 & \text{જો $n$ અવિભાજ્ય હોય}\\
  0 & \text{અન્યથા.}
  \end{cases}
\]
સંબંધિત $1/0$-મૂલ્યવાળું વિધેય અસરકારક રીતે સંગણનીય હોય,
તો આપણે કહીએ કે હા/ના પ્રશ્નનો \emph{અસરકારક રીતે નિર્ણય}
કરી શકાય છે.

કેટલાંક વિધેયો અસરકારક રીતે ગણી શકાતાં નથી, અથવા કેટલીક
સમસ્યાઓનો અસરકારક રીતે નિર્ણય કરી શકાતો નથી, તે ગણિતીય
રીતે સાબિત કરવા સંગણનનું કોઈ એક ચોક્કસ નિદર્શ નક્કી કરવું
અને તે કેટલાંક વિધેયો ગણી કે કેટલીક સમસ્યાઓનો નિર્ણય
કરી શકતું નથી તે બતાવવું આવશ્યક છે. દાખલા તરીકે, દરેક
વિધેય ટ્યુરિંગ મશીન વડે ગણી શકાતું નથી અને દરેક સમસ્યાનો
ટ્યુરિંગ મશીન વડે નિર્ણય કરી શકાતો નથી તે આપણે બતાવી
શકીએ. પછી ચર્ચ--ટ્યુરિંગ માન્યતાનો આધાર લઈને તારવી
શકીએ કે માત્ર ટ્યુરિંગ મશીનો જ દરેક વિધેય ગણવા માટે
અસમર્થ નથી; ટ્યુરિંગ મશીન વડે ન ગણી શકાય તેવા વિધેયો
કોઈ અસરકારક પ્રક્રિયા વડે પણ ગણી શકાતા નથી.

આવા નકારાત્મક પરિણામોની સાબિતીમાં મુખ્ય બાબત એ છે કે
ટ્યુરિંગ મશીનોને જ સંખ્યાઓ આપી શકાય છે. એ માટેનો સૌથી
સરળ રસ્તો તેમની પરિગણના કરવાનો છે: મશીનો અને તેમના
પ્રોગ્રામ લખવાની કોઈ ચોક્કસ રીત નક્કી કરી, પછી તેમને
ક્રમબદ્ધ રીતે યાદીમાં મૂકી શકાય. આમ કરી શકાય છે એ
જોયા પછી, ટ્યુરિંગ મશીન વડે અસંગણનીય વિધેયોનું અસ્તિત્વ
ગણાંકના સરળ વિચારોથી ફલિત થાય છે: $\Nat$થી~$\Nat$માં
જતાં વિધેયોનો ગણ (હકીકતમાં, $\Nat$થી~$\{0, 1\}$માં
જતાં વિધેયોનો ગણ પણ) અપરિગણનીય છે; પરંતુ બધાં
ટ્યુરિંગ મશીનોની પરિગણના કરી શકાય છે, તેથી ટ્યુરિંગ-સંગણનીય
વિધેયોનો ગણ માત્ર ગણનીય છે.

આપણે \emph{ચોક્કસ} વિધેયો અને સમસ્યાઓ પણ વ્યાખ્યાયિત
કરી શકીએ, અને અનુક્રમે તેમનું અસંગણનીય તથા અનિર્ણેય
હોવું સાબિત કરી શકીએ. આવી એક સમસ્યા કહેવાતી
\emph{અટકવાની સમસ્યા} છે. ટ્યુરિંગ મશીનની સૂચનાઓની
યાદી આપીને તેનું સાન્ત વર્ણન કરી શકાય છે. ટ્યુરિંગ
મશીનનું એવું વર્ણન, એટલે ટ્યુરિંગ મશીનનો પ્રોગ્રામ,
અલબત્ત બીજા ટ્યુરિંગ મશીનને ઇનપુટ તરીકે આપી શકાય.
તેથી આપણે એવા ટ્યુરિંગ મશીનો વિચારીએ જે બીજાં
ટ્યુરિંગ મશીનો વિશેના પ્રશ્નોનો નિર્ણય કરે. ખાસ
રસપ્રદ પ્રશ્ન આ છે: ``આપેલું ટ્યુરિંગ મશીન ઇનપુટ~$n$
પર શરૂ કરતાં છેવટે અટકશે?'' આ પ્રશ્નનો નિર્ણય
કરી શકતું ટ્યુરિંગ મશીન હોય તો સારું: તેને ગુણવત્તા
તપાસતું ટ્યુરિંગ મશીન માનો, જે બીજા મશીનો અનંત
ચક્રમાં ફસાઈ ન જાય તેની ખાતરી કરે. ટ્યુરિંગે સાબિત
કરેલી રસપ્રદ હકીકત એ છે કે આવું મશીન હોઈ શકે નહીં.
ટ્યુરિંગ મશીન~$M$ના વર્ણન અને કોઈ સંખ્યા~$n$નો
ઇનપુટ આપતાં હંમેશાં અટકે અને, $M$ ઇનપુટ~$n$ પર
અટકે કે ન અટકે તે મુજબ અનુક્રમે $1$ કે $0$ આપે,
એવું કોઈ એક ટ્યુરિંગ મશીન હોઈ શકે નહીં.

ચોક્કસ અનિર્ણેય સમસ્યાઓનાં ઉદાહરણો મળ્યા પછી બીજી
સમસ્યાઓ પણ અનિર્ણેય છે તે બતાવવા તેમનો ઉપયોગ
કરી શકીએ. દાખલા તરીકે, પ્રસિદ્ધ અનિર્ણેય પ્રશ્ન છે:
``પ્રથમ-ક્રમનું સૂત્ર~$!A$ પ્રામાણ્ય ધરાવે છે?''
ઇનપુટ તરીકે પ્રથમ-ક્રમનું સૂત્ર~$!A$ આપવામાં આવે
ત્યારે, $!A$ પ્રામાણ્ય ધરાવે છે કે નહીં તે મુજબ $1$ કે
$0$ આઉટપુટ સાથે અટકશે તેની ખાતરી હોય એવું કોઈ ટ્યુરિંગ
મશીન નથી. ઐતિહાસિક રીતે, આ સમસ્યાને અસરકારક રીતે
ઉકેલવાની પ્રક્રિયા શોધવાના પ્રશ્નને સાદી રીતે ``નિર્ણય
સમસ્યા'' કહેવાતો; તેથી આપણે કહીએ છીએ કે નિર્ણય સમસ્યા
ઉકેલી શકાતી નથી. ટ્યુરિંગ અને ચર્ચે આશરે એક જ સમયે સ્વતંત્ર
રીતે આ પરિણામ સાબિત કર્યું હતું, તેથી તેને ચર્ચ--ટ્યુરિંગ
પ્રમેય પણ કહે છે.
% END OLP-0265
\endgroup


\begingroup
\makeatletter\def\ol@sectioncs{section}\makeatother

% BEGIN OLP-0266
\olfileid{tur}{und}{enu}
\olsection{ટ્યુરિંગ મશીનોની પરિગણના}
\phantomsection\label{guunit:OLP-0266}


\begin{explain}
બધાં ટ્યુરિંગ મશીનોનો ગણ પરિગણનીય છે એ આપણે બતાવી શકીએ.
દરેક ટ્યુરિંગ મશીનનું સાન્ત વર્ણન આપી શકાય છે, તેમાંથી આ
વાત ફલિત થાય છે. અવસ્થાઓનો ગણ અને ટેપની વર્ણમાળા સાન્ત
ગણો છે. સંક્રમણ વિધેય એ $Q \times \Sigma$થી
$Q \times \Sigma \times \{\TMleft, \TMright, \TMstay\}$માં
જતું આંશિક વિધેય છે; તેથી જે સાન્ત સંખ્યામાં નિવેશ-જોડીઓ માટે તે
વ્યાખ્યાયિત છે, તેમનાં મૂલ્યોની યાદી આપીને તેને પણ
નિર્દિષ્ટ કરી શકાય છે.

અત્યાર સુધીની વાત સાચી છે, છતાં અહીં એક સૂક્ષ્મ ભેદ છે.
ટ્યુરિંગ મશીનની વ્યાખ્યાએ અવસ્થાઓના ગણ અને ટેપની
વર્ણમાળામાં કેવા ઘટકો હોઈ શકે તેના પર કોઈ
પ્રતિબંધ મૂક્યો નથી. તેથી, દાખલા તરીકે, દરેક વાસ્તવિક
સંખ્યા માટે તકનીકી અર્થમાં એવું ટ્યુરિંગ મશીન છે જે તે
સંખ્યાને અવસ્થા તરીકે વાપરે છે. જોકે મશીનનું \emph{વર્તન}
કયા પદાર્થો અવસ્થાઓ અને ચિહ્નો તરીકે કામ કરે છે તેના પર
આધારિત નથી. \olref{fig:variants}નાં બે ટ્યુરિંગ મશીનો જુઓ.
\begin{figure}
\begin{center}
\begin{tikzpicture}[->,>=stealth',shorten >=1pt,auto,node distance=2.8cm,
                    semithick]
  \tikzstyle{every state}=[fill=none,draw=black,text=black]

  \node[initial,state]         (A)              {$q_0$};
  \node[state]         (B) [right of=A] {$q_1$};

  \path (A) edge [bend left] node {\TMtrans{\TMstroke}{\TMstroke}{\TMright}} (B)
        (B) edge [loop above] node {\TMtrans{\TMblank}{\TMblank}{\TMright}} (B)
            edge [bend left] node {\TMtrans{\TMstroke}{\TMstroke}{\TMright}} (A);
\end{tikzpicture}\\
\begin{tikzpicture}[->,>=stealth',shorten >=1pt,auto,node distance=2.8cm,
  semithick]
\tikzstyle{every state}=[fill=none,draw=black,text=black]

\node[initial,state]         (A)              {$s$};
\node[state]         (B) [right of=A] {$h$};

\path (A) edge [bend left] node {\TMtrans{A}{A}{\TMright}} (B)
(B) edge [loop above] node {\TMtrans{\TMblank}{\TMblank}{\TMright}} (B)
edge [bend left] node {\TMtrans{A}{A}{\TMright}} (A);
\end{tikzpicture}
\end{center}
\caption{\emph{સમ} મશીનનાં બે રૂપ}
\ollabel{fig:variants}
\end{figure}
આ બે આલેખ બે મશીનોને અનુરૂપ છે: $M$ની ટેપ-વર્ણમાળા
$\Sigma = \{\TMendtape,\TMblank,\TMstroke\}$ અને અવસ્થાઓનો
ગણ $\{q_0,q_1\}$ છે, જ્યારે $M'$ની વર્ણમાળા $\Sigma' =
\{\TMendtape,\TMblank,A\}$ અને અવસ્થાઓનો ગણ $\{s,h\}$
છે. પરંતુ તેમની બાકીની સૂચનાઓ સમાન છે: $n$ સંખ્યાના
$\TMstroke$ના ક્રમ પર $M$ ત્યારે અને માત્ર ત્યારે અટકે છે
જ્યારે $n$ સમ હોય; અને $n$ સંખ્યાના $A$ના ક્રમ પર $M'$
પણ ત્યારે અને માત્ર ત્યારે અટકે છે જ્યારે $n$ સમ હોય.
આપણે માત્ર $\TMstroke$નું નામ~$A$, $q_0$નું નામ~$s$ અને
$q_1$નું નામ~$h$ રાખ્યાં છે. આ દાખલાને સામાન્ય બનાવી
શકાય: ચિહ્નો અને અવસ્થાઓનાં આવાં નામ બદલવાથી એકમાંથી
બીજું મળે તો ટ્યુરિંગ મશીનોને સમાન માની શકીએ. હકીકતમાં,
આપણે મશીનનાં ચિહ્નો અને અવસ્થાઓને સીધાં ધન પૂર્ણાંકો
જ માની શકીએ: $\sigma_0$ને બદલે~$1$, $\sigma_1$ને
બદલે~$2$ વગેરે; $\TMendtape$ એ~$1$, $\TMblank$ એ~$2$
વગેરે. આ રીતે \emph{સમ} મશીન \olref{fig:standard-even}માં
દર્શાવેલું મશીન બને છે.
\begin{figure}
\[\begin{tikzpicture}[->,>=stealth',shorten >=1pt,auto,node distance=2.8cm,
  semithick]
\tikzstyle{every state}=[fill=none,draw=black,text=black]

\node[initial,state]         (A)              {$1$};
\node[state]         (B) [right of=A] {$2$};

\path (A) edge [bend left] node {\TMtrans{3}{3}{\TMright}} (B)
(B) edge [loop above] node {\TMtrans{2}{2}{\TMright}} (B)
edge [bend left] node {\TMtrans{3}{3}{\TMright}} (A);
\end{tikzpicture}
\]
\caption{માનક \emph{સમ} મશીન}
\ollabel{fig:standard-even}
\end{figure}
જે ટ્યુરિંગ મશીનની અવસ્થાઓ અને ચિહ્નો ધન પૂર્ણાંકો છે
તેને \emph{માનક} મશીન કહી શકીએ અને હવે પછી માત્ર માનક
મશીનો જ વિચારીએ.\footnote{``માનક મશીન'' પરિભાષા
સ્વયં માનક નથી.}

આપણે ટ્યુરિંગ મશીનોનો ગણ પરિગણનીય છે તે બતાવવું
હતું; ઉપરની ચર્ચા મુજબ માનક ટ્યુરિંગ મશીનોનો ગણ
પરિગણનીય છે તે બતાવવું પૂરતું છે. ધારો કે માનક
ટ્યુરિંગ મશીન $M = \tuple{Q, \Sigma, q_0, \delta}$
આપેલું છે. ધન પૂર્ણાંકોની સાન્ત શ્રેણી વડે તેનું વર્ણન
કેવી રીતે કરીએ? પહેલાં અવસ્થાઓની સંખ્યા, પછી અવસ્થાઓ,
ચિહ્નોની સંખ્યા, ચિહ્નો અને પ્રારંભિક અવસ્થાની યાદી આપીએ.
($M$ માનક હોવાથી આ બધાં ધન પૂર્ણાંકો છે એ યાદ રાખો.)
હવે~$\delta$નું શું? શક્ય દલીલોનો ગણ, એટલે કે
$\tuple{q,\sigma}$ જોડો, સાન્ત છે, કારણ કે $Q$ અને
~$\Sigma$ સાન્ત છે. તેથી~$\delta$માં રહેલી માહિતી એ
એવી બધી $5$ ઘટકોવાળી ક્રમયુક્ત શ્રેણીઓ
$\tuple{q, \sigma, q', \sigma', d}$ની
સાન્ત યાદી છે, જેમના માટે
$\delta(q,\sigma) = \tuple{q', \sigma', D}$; અહીં $d$
દિશા~$D$નો સંકેત આપતો અંક છે (જેમ કે~$\TMleft$ માટે
$1$,~$\TMright$ માટે $2$ અને~$\TMstay$ માટે $3$).

આ રીતે દરેક માનક ટ્યુરિંગ મશીનને ધન પૂર્ણાંકોની સાન્ત
યાદી વડે, એટલે કે શ્રેણી $s_M \in (\PosInt)^*$ વડે,
વર્ણવી શકાય છે. દાખલા તરીકે, માનક \emph{સમ} મશીનનો
સંકેત આ શ્રેણી છે:
\[
2, \underbrace{1, 2}_Q, 3, \overbrace{1, 2, 3}^\Sigma, 1, \underbrace{1, 3, 2, 3, 2}_{\delta(1,3) = \tuple{2,3,R}},
\overbrace{2, 2, 2, 2, 2}^{\delta(2,2) = \tuple{2,2,R}},
\underbrace{2, 3, 1, 3, 2}_{\delta(2,3) = \tuple{1,3,R}}.
\]
\end{explain}

\begin{thm}
$\Nat$થી~$\Nat$માં જતાં કેટલાંક વિધેયો ટ્યુરિંગ-સંગણનીય
નથી.
\end{thm}

\begin{proof}
ધન પૂર્ણાંકોની સાન્ત શ્રેણીઓનો ગણ~$(\PosInt)^*$
પરિગણનીય છે એ આપણે જાણીએ છીએ
(\cref{sfr:siz:zigzag:prob:posint-star}). તેથી માનક
ટ્યુરિંગ મશીનોનાં વર્ણનોનો ગણ, $(\PosInt)^*$નો ઉપગણ
હોવાથી, પોતે પણ પરિગણનીય છે. $\Nat$થી~$\Nat$માં જતું
દરેક ટ્યુરિંગ-સંગણનીય વિધેય કોઈ ટ્યુરિંગ મશીન વડે
(હકીકતમાં, અનેક મશીનો વડે) ગણાય છે. તેની અવસ્થાઓ અને
ચિહ્નોનાં નામ ધન પૂર્ણાંકો વડે બદલીને (ખાસ કરીને
$\TMendtape$ને~$1$, $\TMblank$ને~$2$ અને
$\TMstroke$ને~$3$ રાખીને) જોઈ શકાય છે કે દરેક
ટ્યુરિંગ-સંગણનીય વિધેય કોઈ માનક ટ્યુરિંગ મશીન વડે
ગણાય છે. એટલે $\Nat$થી~$\Nat$માં જતાં બધાં
ટ્યુરિંગ-સંગણનીય વિધેયોનો ગણ પણ પરિગણનીય છે.

બીજી તરફ, $\Nat$થી~$\Nat$માં જતાં બધાં વિધેયોનો
ગણ પરિગણનીય નથી (\cref{sfr:siz:red:prob:nat-nat}).
જો દરેક વિધેય કોઈ ટ્યુરિંગ મશીન વડે સંગણનીય હોત,
તો તેમનું સંગણન કરતાં મશીનોનાં બધાં વર્ણનોની યાદી
આપીને બધાં વિધેયોની પરિગણના કરી શકાત. તેથી કેટલાંક
વિધેયો ટ્યુરિંગ-સંગણનીય નથી.
\end{proof}

\begin{prob}
  અવસ્થાઓ અને વર્ણમાળાનાં ચિહ્નોની સ્પષ્ટ યાદી આપવી ન
  પડે એ રીતે ટ્યુરિંગ મશીનોનું વર્ણન કેવી રીતે આપી શકાય
  તે વિચારો. ``માનક'' મશીનની તમારી પોતાની વ્યાખ્યા આપી
  શકો છો; પરંતુ તમારા નવા અર્થમાં દરેક ટ્યુરિંગ મશીન
  જે વિધેય ગણે છે તે કોઈ ``માનક'' મશીન પણ કેમ ગણી
  શકે તે વિશે કંઈક કહો.
  \sourcecorrection{OLUN-004}{સ્થિર અંગ્રેજી મૂળમાં
  મશીનની જ સંગણના કરવાની વાત છે; અહીં અર્થ તેના
  અવસ્થા અને ચિહ્નોનાં નામ બદલીને સમાન સંગણન
  કરતું માનક મશીન મેળવવાનો છે.}
\end{prob}
% END OLP-0266
\endgroup


\begingroup
\makeatletter\def\ol@sectioncs{section}\makeatother

% BEGIN OLP-0267
\olfileid{tur}{und}{uni}
\olsection{સર્વવ્યાપી ટ્યુરિંગ મશીનો}
\phantomsection\label{guunit:OLP-0267}


\olref[enu]{sec}માં આપણે ચર્ચ્યું કે દરેક ટ્યુરિંગ મશીનનું
વર્ણન પૂર્ણાંકોની સાન્ત શ્રેણી વડે કરી શકાય છે. આ શ્રેણી
મશીનની અવસ્થાઓ, વર્ણમાળા, પ્રારંભિક અવસ્થા અને સૂચનાઓને
સંકેતિત કરે છે. આ બધાં વર્ણનોનો ગણ પરિગણનીય છે એ
પણ આપણે નોંધ્યું. આવાં વર્ણનોનો ગણ ગણનીય હોવાથી
~$\Nat$થી આ વર્ણનોમાં જતું કોઈ વ્યાપ્ત વિધેય છે.
દાખલા તરીકે, કેન્ટરની આડીઅવળી રીત વડે આવું
વ્યાપ્ત વિધેય મેળવી શકાય. તે બધાં ટ્યુરિંગ
મશીનોનાં (વર્ણનોનાં) પરિગણનાનો રસ્તો આપે છે. આવી કોઈ
એક પરિગણના નક્કી કરીએ, તો હવે $1$મું, $2$જું, \dots,
$e$મું ટ્યુરિંગ મશીન કહેવાનો અર્થ બને છે. આ અંકોને
\emph{સૂચકો} કહે છે.

\begin{defn}
આપણી નક્કી કરેલી પરિગણનામાં $M$ એ $e$મું ટ્યુરિંગ
મશીન હોય, તો $e$ને~$M$નો \emph{સૂચક} કહીએ છીએ.
$e$મા ટ્યુરિંગ મશીન માટે $M_e$ લખીએ છીએ.
\end{defn}

એક મશીનને એક કરતાં વધુ સૂચકો હોઈ શકે. દાખલા તરીકે,
~$M$નાં બે વર્ણનોમાં તેની સૂચનાઓની યાદીનો ક્રમ જુદો
હોઈ શકે; આવા જુદા વર્ણનોને જુદા સૂચકો મળશે.

અગત્યની વાત એ છે કે ટ્યુરિંગ મશીનોનાં વર્ણનોની પરિગણના
એવી રીતે આપી શકાય છે કે મશીનના સૂચકમાંથી~$M$નું વર્ણન
અસરકારક રીતે ગણી શકાય, અને મશીન~$M$ના વર્ણનમાંથી
તેનો કોઈ સૂચક અસરકારક રીતે ગણી શકાય. પછી ચર્ચ--ટ્યુરિંગ
માન્યતા મુજબ એવું ટ્યુરિંગ મશીન મેળવી શકાય છે જે
સૂચક~$e$વાળા ટ્યુરિંગ મશીનનું વર્ણન પાછું મેળવે અને
તે વર્ણનને આઉટપુટ તરીકે પોતાની ટેપ પર લખે. આ વર્ણન
~$\TMstroke$ના ખંડોની શ્રેણી હશે (જે~$M_e$નું વર્ણન
કરતી શ્રેણીના ધન પૂર્ણાંકો દર્શાવે છે).

આથી હવે સ્વાભાવિક રીતે પ્રશ્ન થાય: ટ્યુરિંગ મશીનોના
સૂચકો પરનાં કયાં વિધેયો પોતે ટ્યુરિંગ મશીન વડે સંગણનીય
છે? સૂચકોનાં કયાં ગુણધર્મોનો ટ્યુરિંગ મશીન વડે નિર્ણય
કરી શકાય? દાખલો લો: સૂચક~$e$ને સૂચક~$e$વાળા ટ્યુરિંગ
મશીનની અવસ્થાઓની સંખ્યા સાથે જોડતું વિધેય ટ્યુરિંગ
મશીન વડે સંગણનીય છે. આવું મશીન શું કરશે તે જુઓ:
તેની ટેપ પર $e$ જેટલા~$\TMstroke$નો એક ખંડ હોય ત્યારે
તે પહેલાં $e$માંથી મશીનનું વર્ણન કાઢશે. હવે ટેપ પર
વર્ણન~$\TMstroke$ના ખંડોની શ્રેણી વડે દર્શાવાયેલું છે.
આ શ્રેણીનો પહેલો ઘટક અવસ્થાઓની સંખ્યા છે.
તેથી હવે પહેલા~$\TMstroke$ના ખંડ સિવાયનું બધું ભૂંસીને
મશીન અટકી જાય એટલું જ કરવાનું રહે છે.

નીચેનું પરિણામ નોંધપાત્ર છે:

\begin{thm}\ollabel{thm:universal-tm} એવું \emph{સર્વવ્યાપી
  ટ્યુરિંગ મશીન}~$U$ છે કે ઇનપુટ $\tuple{e,n}$ પર શરૂ કરતાં
  \begin{enumerate}
    \item $M_e$ ઇનપુટ~$n$ પર અટકે ત્યારે અને માત્ર ત્યારે જ તે અટકે છે, અને
    \item $M_e$ આઉટપુટ $m$ સાથે અટકે, તો~$U$ પણ તેમ જ અટકે છે.
  \end{enumerate}
  આમ, $U$ એ $f\colon \Nat \times \Nat \pto \Nat$ વિધેય
  ગણે છે: $M_e$ ઇનપુટ~$n$ પર શરૂ કરીને આઉટપુટ~$m$
  સાથે અટકે ત્યારે $f(e,n) = m$, અને અન્યથા તે અવ્યાખ્યાયિત છે.
\end{thm}

\begin{proof}
  ખરેખર~$U$નું સંપૂર્ણ બાંધકામ આપવું વ્યવહારમાં લગભગ અશક્ય
  છે, કારણ કે તે અત્યંત જટિલ મશીન છે. પરંતુ તે કેવી રીતે
  કામ કરે છે તેની રૂપરેખા આપી, પછી ચર્ચ--ટ્યુરિંગ માન્યતાનો
  આધાર લઈ શકીએ. શરૂઆતમાં~$U$ની ટેપ પર $e$ જેટલા
  $\TMstroke$નો ખંડ, ત્યાર પછી $n$ જેટલા~$\TMstroke$નો
  ખંડ હોય છે. તે પહેલાં સૂચક~$e$માંથી મશીનનું વર્ણન
  કાઢીને ઇનપુટ~$n$ની જમણી બાજુ લખે છે. આમ મશીન~$M_e$ની
  સૂચનાઓનું વર્ણન કરતી સંખ્યાઓની યાદી મળે છે (એટલે
  ~$\TMblank$થી છૂટા પડેલા~$\TMstroke$ના ખંડો).
  પછી~$U$,~$M_e$ના વર્ણનની જમણી બાજુ,~$M_e$ની પ્રારંભિક
  અવસ્થાનો અંક અને સંખ્યા~$1$ ટેપ પર લખે છે. (ફરીથી,
  આ અંકો~$\TMstroke$ના ખંડો તરીકે એકકી સ્વરૂપમાં દર્શાવાય
  છે.) ત્યાર પછી તે ઇનપુટ ($n$ જેટલા~$\TMstroke$નો ખંડ)
  જમણી બાજુ નકલ કરે છે---પરંતુ દરેક $\TMstroke$ને
  ત્રણ $\TMstroke$ના ખંડથી બદલે છે (યાદ રાખો,
  $\TMstroke$ ચિહ્નનો અંક~$3$ છે; $\TMendtape$નો અંક
  $1$ અને $\TMblank$નો અંક $2$ છે). આ ખંડોની શ્રેણીના
  ડાબા છેડે (~$U$ની ટેપ પર $\TMblank$ ચિહ્નોથી છૂટા
  પડેલા ખંડોમાં) તે એક~$\TMstroke$ લખે છે, જે
  ~$\TMendtape$નો સંકેત છે.

  હવે~$U$ની ટેપ પર આ બધું છે: સૂચક~$e$, સંખ્યા~$n$,
  પ્રારંભિક અવસ્થાનો સંકેત-અંક (``વર્તમાન અવસ્થા''),
  પ્રારંભિક હેડ-સ્થાનનો અંક~$1$ (``વર્તમાન હેડ-સ્થાન''),
  અને ``ટેપ''ની પ્રારંભિક સામગ્રી (~$\TMblank$થી છૂટા
  પડેલા~$\TMstroke$ના ખંડોની શ્રેણી, જે~$M_e$નાં
  ચિહ્નોના સંકેત-અંકો---``ચિહ્નો''---દર્શાવે છે).

  હવે નીચેના પગલાં લઈને તે ઇનપુટ~$n$ પર શરૂ કરાયેલું
  $M_e$ શું કરે તેની અનુકૃતિ કરે છે:
  \begin{enumerate}
    \item ``વર્તમાન હેડ-સ્થાન''નો અંક~$k$ શોધો
    (શરૂઆતમાં તે~$1$ છે),
    \item ``ટેપ''ના $k$મા ખંડ પાસે જઈને ત્યાં કયું
    ``ચિહ્ન'' છે તે જુઓ,
    \item\ollabel{find-inst}%
    વર્તમાન ``અવસ્થા'' અને ``ચિહ્ન''ને મળતી સૂચના શોધો,
    \item ``ટેપ''ના $k$મા ખંડ પર પાછા જઈ ત્યાંનું
    ``ચિહ્ન''~$M_e$ જે ચિહ્ન લખે તેના સંકેત-અંકથી બદલો,
    \item વર્તમાન ``અવસ્થા''ની નોંધ જ્યાં છે ત્યાં હેડ
    લઈ જઈ એ અંકને નવી અવસ્થાના અંકથી બદલો,
    \item ``ટેપ-સ્થાન''ની નોંધ પાસે જઈ એક~$\TMstroke$
    ભૂંસો અથવા એક~$\TMstroke$ ઉમેરો (સૂચના અનુક્રમે
    ડાબે કે જમણે જવાનું કહે ત્યારે).
    \item ફરી કરો.\footnote{અહીં કેટલીક સૂક્ષ્મ મુશ્કેલીઓ
    છોડીને ચાલીએ છીએ. દાખલા તરીકે, ``વર્તમાન હેડ-સ્થાન''નો
    હિસાબ રાખતું ગણક વધારતી વખતે~$U$ને વધારે જગ્યા
    જોઈતી થઈ શકે---એ કિસ્સામાં આખી ``ટેપ'' જમણી બાજુ
    ખસેડવી પડશે.}
  \end{enumerate}
  ઇનપુટ~$n$ પર શરૂ કરાયેલું $M_e$ કદી ન અટકે, તો~$U$
  પણ કદી નથી અટકતું, એટલે તેનું આઉટપુટ અવ્યાખ્યાયિત છે.

  પગલું~\olref{find-inst}માં~$M_e$ના વર્ણનમાં વર્તમાન
  ``અવસ્થા''/``ચિહ્ન'' જોડ માટે કોઈ સૂચના નથી એવું જણાય,
  તો~$M_e$ અટકતું. આ થાય તો~$U$ પોતાની ટેપનો
  ``ટેપ''ની ડાબી બાજુનો ભાગ ભૂંસે છે. તે ત્રણ
  ~$\TMstroke$ના દરેક ખંડ માટે (જે~$M_e$ની ટેપ પરના
  એક~$\TMstroke$ને દર્શાવે છે) પોતાની ટેપના ડાબા છેડે
  એક $\TMstroke$ લખે છે અને ``ટેપ''ને ક્રમે ભૂંસે છે.
  આ પૂરું થાય ત્યારે~$U$ની ટેપ પર લંબાઈ~$m$વાળો
  ~$\TMstroke$નો એક જ ખંડ હોય છે.

  ``ટેપ'' પર ત્રણ~$\TMstroke$ના ખંડ સિવાયનું કંઈ
  $U$ને મળે, તો તે તરત અટકી જાય છે. આ કિસ્સામાં
  ~$U$ની ટેપ પર~$\TMstroke$નો એક જ ખંડ ન હોવાથી
  તેનું આઉટપુટ પ્રાકૃતિક સંખ્યા નથી, એટલે આ કિસ્સામાં
  $f(e,n)$ અવ્યાખ્યાયિત છે.
\end{proof}
% END OLP-0267
\endgroup


\begingroup
\makeatletter\def\ol@sectioncs{section}\makeatother

% BEGIN OLP-0268
\olfileid{tur}{und}{hal}
\olsection{અટકવાની સમસ્યા}
\phantomsection\label{guunit:OLP-0268}


\begin{explain}
માનો કે આપણે ટ્યુરિંગ મશીનોનાં વર્ણનોની કોઈ પરિગણના
નક્કી કરી છે. તેથી દરેક ટ્યુરિંગ મશીનને એક
\emph{સૂચક} મળે છે: મશીનોનાં વર્ણનોની પરિગણના
$M_1$, $M_2$, $M_3$, \dots{}માં તેનું સ્થાન.

ટ્યુરિંગ-અસંગણનીય વિધેયો હોવાં જ જોઈએ એ આપણે
જાણીએ છીએ: ટ્યુરિંગ મશીનોનાં વર્ણનોનો ગણ---અને તેથી
ટ્યુરિંગ મશીનોનો ગણ---પરિગણનીય છે, પરંતુ
$\Nat$થી $\Nat$માં જતાં બધાં વિધેયોનો ગણ એવો નથી.
તેમ છતાં આપણે અસંગણનીય વિધેયોનાં ચોક્કસ ઉદાહરણો પણ
મેળવી શકીએ. આવું એક વિધેય અટકવાનું વિધેય છે.
\end{explain}

\begin{defn}[અટકવાનું વિધેય]
 \emph{અટકવાનું વિધેય}~$h$ આ રીતે વ્યાખ્યાયિત છે:
\[
h(e,n) =
\begin{cases}
  \text{0} & \text{જો મશીન~$M_e$ ઇનપુટ $n$ પર ન અટકે} \\
  \text{1} & \text{જો મશીન~$M_e$ ઇનપુટ $n$ પર અટકે}
\end{cases}
\]
\end{defn}

\begin{defn}[અટકવાની સમસ્યા]
\emph{અટકવાની સમસ્યા} એટલે કોઈ પણ $e$, $n$ માટે,
ટ્યુરિંગ મશીન~$M_e$~$n$ સ્ટ્રોકના
ઇનપુટથી શરૂ કરતાં અટકે છે કે નહીં તે નક્કી કરવાની
સમસ્યા.
\end{defn}

\begin{explain}
સંબંધિત વિધેય~$s$ ટ્યુરિંગ-સંગણનીય નથી એવું બતાવીને
આપણે~$h$ પણ ટ્યુરિંગ-સંગણનીય નથી એવું બતાવીએ છીએ.
આ સાબિતી એ હકીકતો પર આધાર રાખે છે કે ટ્યુરિંગ મશીન
વડે ગણી શકાય તે બધું શિસ્તબદ્ધ ટ્યુરિંગ મશીન વડે પણ
ગણી શકાય છે (\olref[mac][dis]{sec}), અને બે ટ્યુરિંગ
મશીનોને જોડીને એક મશીન બનાવી શકાય છે
(\olref[mac][cmb]{sec}).
\end{explain}

\begin{defn} વિધેય~$s$ આ રીતે વ્યાખ્યાયિત છે:
\[
s(e) =
\begin{cases}
  \text{0} & \text{જો મશીન~$M_e$ ઇનપુટ $e$ પર ન અટકે} \\
  \text{1} & \text{જો મશીન~$M_e$ ઇનપુટ $e$ પર અટકે}
\end{cases}
\]
\end{defn}

\begin{lem}
વિધેય~$s$ ટ્યુરિંગ-સંગણનીય નથી.
\end{lem}

\begin{proof}
વિરોધાભાસ માટે માનો કે વિધેય~$s$ ટ્યુરિંગ-સંગણનીય
છે. તો~$s$ ગણતું ટ્યુરિંગ મશીન~$S$ હોત. વ્યાપકતા
ગુમાવ્યા વિના ધારી શકીએ કે $S$ અટકે ત્યારે પ્રથમ
ખાનાને વાંચતું હોય (એટલે તે શિસ્તબદ્ધ હોય). આ
મશીનને બીજા મશીન~$J$ સાથે ``જોડી'' શકાય. તે ઇનપુટ
~$0$ પર શરૂ થાય તો અટકે છે (એટલે શરૂઆતની અવસ્થામાં,
ટેપના છેડાના ચિહ્નની જમણી બાજુના ખાનાને વાંચતી વખતે
$\TMblank$ વાંચે તો અટકે છે); અન્યથા તે કદી ન અટકીને
જમણી બાજુ ભટકતું રહે છે. $S$ને~$J$ સાથે જોડીને
બનેલું મશીન $S \concat J$ પણ ટ્યુરિંગ મશીન છે;
તેથી કોઈ~$e$ માટે તે $M_e$ છે (એટલે તે પરિગણનામાં
ક્યાંક આવે છે). હવે~$M_e$ને~$e$ જેટલા $\TMstroke$ના
ઇનપુટથી શરૂ કરો. બે શક્યતાઓ છે: $M_e$ અટકે છે
અથવા નથી અટકતું.
\begin{enumerate}
\item માનો કે~$M_e$,~$e$ જેટલા $\TMstroke$ના ઇનપુટ
  પર અટકે છે. તો $s(e) = 1$. તેથી~$S$ને ઇનપુટ~$e$
  પર શરૂ કરતાં તે ટેપ પર આઉટપુટ તરીકે એક જ
  $\TMstroke$ રાખીને અટકે છે. ત્યાર પછી~$J$ ટેપ પર
  $\TMstroke$ હોય ત્યારે શરૂ થાય છે. આ કિસ્સામાં~$J$
  અટકતું નથી. પરંતુ~$M_e$ એ $S \concat J$ મશીન જ
  છે, તેથી તેને $S$ અને પછી $J$ જે કરે તે જ કરવું
  પડે (એટલે આ કિસ્સામાં જમણી બાજુ ભટકીને કદી ન
  અટકવું પડે). તેથી~$e$ જેટલા $\TMstroke$ના ઇનપુટ
  પર~$M_e$ અટકી શકે નહીં.

\item હવે માનો કે~$M_e$,~$e$ જેટલા $\TMstroke$ના
  ઇનપુટ પર અટકતું નથી. તો $s(e) = 0$, અને~$S$ને
  ઇનપુટ~$e$ પર શરૂ કરતાં તે ખાલી ટેપ સાથે અટકે છે.
  ખાલી ટેપ પર શરૂ કરાયેલું~$J$ તરત અટકે છે. ફરી,
  $M_e$ એ $S$ અને પછી~$J$ જે કરે તે જ કરે છે;
  તેથી~$e$ જેટલા $\TMstroke$ના ઇનપુટ પર~$M_e$
  અટકવું જ જોઈએ.
\end{enumerate}
બંને કિસ્સામાં ધારણા સાથે વિરોધાભાસ મળે છે. તેથી
આવું ટ્યુરિંગ મશીન~$S$ હોઈ શકે નહીં: $s$ ટ્યુરિંગ-સંગણનીય
નથી.
\end{proof}

\begin{thm}[અટકવાની સમસ્યા ઉકેલી ન શકાય]
\ollabel{thm:halting-problem} અટકવાની સમસ્યા ઉકેલી
શકાતી નથી, એટલે વિધેય~$h$ ટ્યુરિંગ-સંગણનીય નથી.
\end{thm}

\begin{proof}
માનો કે~$h$ ટ્યુરિંગ-સંગણનીય છે અને તેને ટ્યુરિંગ
મશીન~$H$ ગણે છે. $H$ વડે~$s$ ગણતું ટ્યુરિંગ મશીન
બનાવી શકીએ: પહેલાં ઇનપુટની એક નકલ કરો (તેને
~$\TMblank$ ચિહ્નથી છૂટી રાખીને); પછી શરૂઆતમાં
પાછા જઈ~$H$ ચલાવો. પહેલું કામ કરતું મશીન સ્પષ્ટપણે
બનાવી શકાય છે (\cref{tur:mac:dis:prob:copier});
અને જો $H$ હોત, તો તેને આવા નકલ કરનારા મશીન સાથે
``જોડીને'' એવું નવું મશીન મેળવી શકાત જે~$M_e$
ઇનપુટ~$e$ પર અટકે છે કે નહીં તે નક્કી કરે, એટલે~$s$
ગણે. પરંતુ આવું કોઈ મશીન હોઈ શકે નહીં એ આપણે પહેલેથી
બતાવ્યું છે. તેથી~$h$ પણ ટ્યુરિંગ-સંગણનીય નથી.
\end{proof}

\begin{prob}
ત્રણ-અટક (3-Halt) સમસ્યા એટલે અન્યથા ખાલી ટેપ પર
ત્રણ $\TMstroke$ના ઇનપુટ માટે મનસ્વી પસંદ કરેલું
ટ્યુરિંગ મશીન અટકે છે કે નહીં તે નક્કી કરતી નિર્ણય-વિધિ
આપવાની સમસ્યા. ત્રણ-અટક સમસ્યા ઉકેલી શકાતી નથી
તે સાબિત કરો.
\end{prob}

\begin{prob}
બતાવો કે~$e \neq n$ હોય તેવા ટ્યુરિંગ મશીન અને ઇનપુટ
જોડ $M_e$ તથા~$n$ માટે અટકવાની સમસ્યા ઉકેલી શકાય,
તો $e = n$ હોય તેવા કિસ્સાઓ માટે પણ ઉકેલી શકાય છે.
\end{prob}

\begin{prob}
અમે સાબિત કર્યું કે ઇનપુટ તરીકે સંખ્યા~$e$ આપવામાં આવે,
જે બધાં ટ્યુરિંગ મશીનોની પરિગણનામાં મશીન~$M_e$ને
ઓળખાવે છે, ત્યારે અટકવાની સમસ્યા ઉકેલી શકાતી નથી.
જો ઇનપુટ તરીકે~\olref[tur][und][enu]{sec}માંનાં
ટ્યુરિંગ મશીનોનાં વર્ણનો સીધાં આપવાની છૂટ આપીએ તો?
શું સૂચકોને બદલે મશીનોનાં વર્ણનો ઇનપુટ તરીકે લેતું,
અટકવાની સમસ્યાનો નિર્ણય કરતું ટ્યુરિંગ મશીન હોઈ
શકે? કેમ અથવા કેમ નહીં તે સમજાવો.
\end{prob}

\begin{prob} બતાવો કે નીચે પ્રમાણે વ્યાખ્યાયિત
  \emph{આંશિક} વિધેય~$s'$ ટ્યુરિંગ-સંગણનીય \emph{છે}:
  \[
  s'(e) =
  \begin{cases}
    \text{1} & \text{જો મશીન~$M_e$ ઇનપુટ $e$ પર અટકે}\\
    \text{અવ્યાખ્યાયિત} & \text{જો મશીન~$M_e$ ઇનપુટ $e$ પર ન અટકે}
  \end{cases}
  \]
\end{prob}
% END OLP-0268
\endgroup


\begingroup
\makeatletter\def\ol@sectioncs{section}\makeatother

% BEGIN OLP-0269
\olfileid{tur}{und}{dec}
\olsection{નિર્ણય સમસ્યા}
\phantomsection\label{guunit:OLP-0269}


આપેલું વાક્ય પ્રામાણ્ય ધરાવે છે કે નહીં તે
નક્કી કરવાની અસરકારક રીત હોય ત્યારે અને માત્ર ત્યારે
જ આપણે પ્રથમ-ક્રમનું તર્કશાસ્ત્ર \emph{નિર્ણેય} છે એમ
કહીએ છીએ. હકીકતમાં, આવી કોઈ રીત નથી: પ્રથમ-ક્રમનાં
વાક્યોનું પ્રામાણ્ય નક્કી કરવાની સમસ્યા ઉકેલી શકાતી નથી.

આ મહત્ત્વનું નકારાત્મક પરિણામ સ્થાપિત કરવા આપણે
નિર્ણય સમસ્યા ટ્યુરિંગ મશીન વડે ઉકેલી શકાતી નથી તે
સાબિત કરીએ છીએ. એટલે, એવું કોઈ ટ્યુરિંગ મશીન નથી
જેની ટેપ પર પ્રથમ-ક્રમનું વાક્ય હોય ત્યારે
તે શરૂ કરીને છેવટે અટકે અને તે વાક્ય પ્રામાણ્ય ધરાવે
છે કે નહીં તે મુજબ $1$ અથવા~$0$ આઉટપુટ આપે.
ચર્ચ--ટ્યુરિંગ માન્યતા મુજબ દરેક સંગણનીય વિધેય
ટ્યુરિંગ-સંગણનીય છે. તેથી જો આ ``પ્રામાણ્ય વિધેય''
કોઈ પણ રીતે અસરકારક રીતે સંગણનીય હોત, તો તે
ટ્યુરિંગ-સંગણનીય હોત. તે ટ્યુરિંગ-સંગણનીય નથી,
તો અસરકારક રીતે પણ સંગણનીય હોઈ શકે નહીં.

નિર્ણય સમસ્યા ઉકેલી શકાતી નથી તે સાબિત કરવાની
આપણી રીત અટકવાની સમસ્યાનું તેમાં ન્યૂનીકરણ કરવાની
છે. તેનો અર્થ આ છે: ટ્યુરિંગ મશીનનો સૂચક~$e$
અને ઇનપુટ~$w$ માટે, મશીન અટકે તો $1$ અને
અન્યથા~$0$ મૂલ્ય આપતું વિધેય~$h(e,w)$ ટ્યુરિંગ-સંગણનીય
નથી એ આપણે સાબિત કર્યું છે. આપણે બતાવીશું કે
પ્રથમ-ક્રમનાં વાક્યોનું પ્રામાણ્ય નક્કી કરતું ટ્યુરિંગ
મશીન હોત, તો~$h$ ગણતું ટ્યુરિંગ મશીન પણ હોત.
પરંતુ~$h$ને ટ્યુરિંગ મશીન વડે ગણી શકાતું નથી;
તેથી પ્રામાણ્ય નક્કી કરતું મશીન પણ હોઈ શકે નહીં.

આ રીતનું પહેલું પગલું એ બતાવવાનું છે કે દરેક ઇનપુટ~$w$
અને ટ્યુરિંગ મશીન~$M$ માટે આપણે અસરકારક રીતે
વાક્ય $!T(M, w)$ આપી શકીએ, જે~$M$ની
સૂચનાઓનો ગણ અને ઇનપુટ~$w$ નિરૂપે, અને
વાક્ય~$!E(M, w)$ આપી શકીએ, જે ``$M$ છેવટે
અટકે છે'' એવું વ્યક્ત કરે, એવી રીતે કે:
\begin{quote}
  $\Entails !T(M, w) \lif !E(M,w)$ ત્યારે અને માત્ર
  ત્યારે જ, જ્યારે $M$ ઇનપુટ~$w$ પર અટકે છે.
\end{quote}
આ વાક્યો $!T(M, w)$ અને~$!E(M, w)$નું વર્ણન અને
$!T(M, w) \lif !E(M, w)$ ત્યારે અને માત્ર ત્યારે જ
પ્રામાણ્ય ધરાવે છે જ્યારે $M$ ઇનપુટ~$w$ પર અટકે
છે તેની તપાસ, આપણી સાબિતીનો મુખ્ય ભાગ હશે.
% END OLP-0269
\endgroup


\begingroup
\makeatletter\def\ol@sectioncs{section}\makeatother

% BEGIN OLP-0270
\olfileid{tur}{und}{rep}
\olsection{ટ્યુરિંગ મશીનોનું નિરૂપણ}
\phantomsection\label{guunit:OLP-0270}


\begin{explain}
ટ્યુરિંગ મશીનો અને તેમનું વર્તન પ્રથમ-ક્રમના તર્કના
વાક્ય વડે નિરૂપવા માટે યોગ્ય ભાષા વ્યાખ્યાયિત
કરવી પડશે. આ ભાષાના બે ભાગ છે: મશીનનાં સંરૂપણોનું
વર્ણન કરવા માટેનાં વિધેયો, અને સંગણનનાં પગલાં
(``ક્ષણો'') તથા ટેપ પરનાં સ્થાનોને અંક આપવા માટેની
અભિવ્યક્તિઓ.

બે જાતનાં વિધેયો લઈએ છીએ, અને બંને દ્વિસ્થાનિક
છે: દરેક અવસ્થા~$q$ માટે વિધેય~$\Obj Q_q$,
અને દરેક ટેપ-ચિહ્ન~$\sigma$ માટે
વિધેય~$\Obj S_\sigma$. પહેલાંના વિધાનો
આપણને~$M$ની અવસ્થા અને તેના ટેપ-હેડનું સ્થાન
વર્ણવવા દે છે; બીજાં વિધાનો ટેપની સામગ્રીનું વર્ણન
કરવા દે છે.

ટેપ-હેડનાં સ્થાનો અને થઈ ચૂકેલાં પગલાંની સંખ્યા
વ્યક્ત કરવા સંખ્યાઓ વ્યક્ત કરવાની રીત જોઈએ. તે માટે
અચળ~$\Obj 0$ અને $1$-સ્થાનિક અનુગામી
વિધેય~$\prime$ વાપરીએ છીએ. રૂઢિ મુજબ તેને તેના
આર્ગ્યુમેન્ટની \emph{પછી} લખીએ છીએ (અને કૌંસ
છોડી દઈએ છીએ).

દરેક સંખ્યા $n$ માટે પ્રમાણભૂત પદ~$\num{n}$ છે,
એટલે~$n$નું \emph{આંકિક પદ}, જે~$\Lang L_M$માં
તેનું નિરૂપણ કરે છે. $\num{0}$ એ $\Obj 0$,
$\num{1}$ એ $\Obj 0'$, $\num{2}$ એ $\Obj 0''$,
વગેરે છે. વધુ ઔપચારિક રીતે:
\begin{align*}
\num{0} & = \Obj 0 \\
\num{n+1} &= \num{n}'
\end{align*}
પદ $\num{0}$, એટલે $\Obj 0$, ટેપના સૌથી ડાબા
સ્થાનને તેમજ પ્રથમ પગલા પહેલાંના સમયને (પ્રારંભિક
સંરૂપણને) નામ આપે છે. પદ $\num{1}$, એટલે
$\Obj 0'$, સૌથી ડાબા ખાનાની જમણી બાજુના ખાનાને
અને પ્રથમ પગલા પછીના સમયને નામ આપે છે; આગળ
પણ એમ જ છે.

ટેપનાં સ્થાનોનો ક્રમ (અર્થ ``ની ડાબે'') તથા
સંગણનનાં પગલાંનો ક્રમ (અર્થ ``પહેલાં'') બંને વ્યક્ત
કરવા વિધેય~$<$ પણ લઈએ છીએ.

ભાષા નક્કી થયા પછી~$!T(M, w)$નાં ``સ્વયંસિદ્ધ વાક્યો''
યાદીબદ્ધ કરીએ: એટલે એવા વાક્યો, જે સાથે
મળીને ઇનપુટ~$w$ પર ચાલતાં~$M$નું વર્તન વર્ણવે.
$\Obj 0$, $\prime$ અને $<$ પર શરતો મૂકતાં વાક્યો,
પ્રારંભિક સંરૂપણ વર્ણવતાં વાક્યો અને ચોક્કસ સૂચના
ચલાવ્યા પછી~$M$નું સંરૂપણ કેવું હોય તે વર્ણવતાં
વાક્યો હશે.
\end{explain}

\begin{defn}
  \ollabel{defn:tm-descr}
ટ્યુરિંગ મશીન $M = \tuple{Q, \Sigma, q_0, \delta}$
આપેલું હોય, તો ભાષા~$\Lang L_M$માં નીચેનું આવે છે:
\begin{enumerate}
\item દરેક અવસ્થા~$q \in Q$ માટે દ્વિસ્થાનિક
  વિધેય $\Obj Q_q(x, y)$. સહજ રીતે,
  $\Obj Q_q(\num{m}, \num{n})$ વ્યક્ત કરે છે કે
  ``$n$ પગલાં પછી~$M$,~$m$મું ખાનું વાંચતાં
  અવસ્થા~$q$માં છે.''
\item દરેક ચિહ્ન~$\sigma\in \Sigma$ માટે દ્વિસ્થાનિક
  વિધેય $\Obj S_\sigma(x, y)$. સહજ રીતે,
  $\Obj S_\sigma(\num{m}, \num{n})$ વ્યક્ત કરે છે કે
  ``$n$ પગલાં પછી~$m$મા ખાનામાં ચિહ્ન~$\sigma$ છે.''
\item એક અચળ $\Obj 0$
\item એક એકસ્થાનિક ફલન $\prime$
\item એક દ્વિસ્થાનિક વિધેય $<$
\end{enumerate}
\end{defn}

ટ્યુરિંગ મશીન~$M$, ઇનપુટ
$w = \sigma_{i_1}\dots\sigma_{i_k}$ પર કેવી રીતે
કામ કરે તે વર્ણવતાં વાક્યો નીચે મુજબ છે:
\begin{enumerate}
\item સંખ્યાઓ અને~$<$નું વર્ણન કરતાં સ્વયંસિદ્ધ વાક્યો:
\begin{enumerate}
\item દરેક સંખ્યા તેના અનુગામી કરતાં નાની છે એવું
કહેતું વાક્ય:
\[
\lforall[x][x < x']
\]
\item $<$ સંક્રામક છે તેની ખાતરી આપતું વાક્ય:
\[
\lforall[x][\lforall[y][\lforall[z][
      ((x < y \land y < z) \lif x < z)]]]
\]
\end{enumerate}
\item ઇનપુટ-સંરૂપણનું વર્ણન કરતાં સ્વયંસિદ્ધ વાક્યો:
\begin{enumerate}
\item મશીન શરૂ થાય તે પહેલાં, એટલે $0$~પગલાં પછી,
  $M$ પ્રારંભિક અવસ્થા~$q_0$માં છે અને ખાનું~$1$
  વાંચે છે:
\[
\Obj Q_{q_0}(\num{1}, \num{0})
\]
\item પ્રથમ $k+1$ ખાનાંમાં ચિહ્નો $\TMendtape$,
  $\sigma_{i_1}$, \dots, $\sigma_{i_k}$ છે:
\[
\Obj S_\TMendtape(\num{0}, \num{0}) \land
\Obj S_{\sigma_{i_1}}(\num{1}, \num{0}) \land
\dots \land
\Obj S_{\sigma_{i_k}}(\num{k}, \num{0})
\]
\item બાકીની ટેપ ખાલી છે:
\[
\lforall[x][(\num{k} < x \lif \Obj S_\TMblank(x, \num{0}))]
\]
\end{enumerate}
\item એક સંરૂપણમાંથી આગળના સંરૂપણમાં થતું સંક્રમણ
  વર્ણવતાં સ્વયંસિદ્ધ વાક્યો:

આગળ માટે, $!A(x, y)$ એ નીચેના સ્વરૂપનાં બધાં
વાક્યોનું સંયોજન માનો:
\[
\lforall[z][
  (((z < x \lor x < z) \land \Obj S_\sigma(z, y))
  \lif \Obj S_\sigma(z, y'))]
\]
અહીં $\sigma \in \Sigma$. $!A(\num{m},\num{n})$
વડે આપણે કહીએ છીએ કે ``$m$મા ખાના સિવાય,
$n+1$ પગલાં પછીની ટેપ $n$ પગલાં પછી જેવી જ છે.''
\begin{enumerate}
\item \ollabel{rep-right} દરેક સૂચના
  $\delta(q_i, \sigma) =
  \tuple{q_j, \sigma', \TMright}$ માટે વાક્ય:
\begin{align*}
& \lforall[x][\lforall[y][(
   (\Obj Q_{q_i}(x, y) \land \Obj S_{\sigma}(x, y)) \lif {}]] \\
&\qquad   (\Obj Q_{q_j}(x', y') \land \Obj S_{\sigma'}(x, y') \land
!A(x, y)))
\end{align*}
એનો અર્થ એ છે કે~$y$ પગલાં પછી મશીન
અવસ્થા~$q_i$માં, ચિહ્ન~$\sigma$ ધરાવતું ખાનું~$x$
વાંચે છે, તો $y+1$ પગલાં પછી તે ખાનું~$x+1$ વાંચે છે,
અવસ્થા~$q_j$માં છે, ખાનું~$x$ હવે~$\sigma'$ ધરાવે છે,
અને~$x$ સિવાયનું દરેક ખાનું~$y$ પગલાં પછી જે ચિહ્ન
ધરાવતું હતું તે જ ધરાવે છે.

\item \ollabel{rep-left} દરેક સૂચના
  $\delta(q_i, \sigma) =
  \tuple{q_j, \sigma', \TMleft}$ માટે વાક્ય:
\begin{align*}
& \lforall[x][\lforall[y][
    ((\Obj Q_{q_i}(x', y) \land \Obj S_{\sigma}(x', y)) \lif {}]]\\
& \qquad   (\Obj Q_{q_j}(x, y') \land \Obj S_{\sigma'}(x', y') \land
!A(x', y))) \land {}\\
& \lforall[y][((\Obj Q_{q_i}(\num{0}, y) \land \Obj S_{\sigma}(\num{0},
    y)) \lif {}]\\
& \qquad (\Obj Q_{q_j}(\num{0}, y') \land \Obj S_{\sigma'}(\num{0},
  y') \land !A(\num{0}, y)))
\end{align*}
\sourcecorrection{OLUN-002}{સ્થિર અંગ્રેજી મૂળની
ડાબે-ચાલ સૂચનાની પહેલી શાખામાં લખાતું ખાનું
$x'$ છે, પરંતુ તેનું $!A(x,y)$ એ ખાનાને
યથાવત્ રહેવાની શરતમાંથી બાકાત રાખતું નથી.
અહીં લખાતા ખાનાને બાકાત રાખવા $!A(x',y)$ કર્યું છે.}
આ કેવી રીતે કામ કરે છે તે વિચારો: હવે આપણે
``ખાનું~$x$ વાંચતું હોય તો \dots''થી શરૂ કરતા નથી,
પણ ``ખાનું $x+1$ વાંચતું હોય તો \dots''થી શરૂ કરીએ
છીએ. ડાબે જવાથી આગળના પગલે મશીન ખાનું~$x$
વાંચે છે. પરંતુ લખાણ ખાનું~$x+1$ પર થાય છે.
બાદબાકી કે પૂર્વગામી વિધેય ન હોવાથી આ રીતે કરીએ
છીએ.

નોંધો કે $x+1$ સ્વરૂપની સંખ્યાઓ $1$, $2$, \dots છે;
એટલે મશીન ખાનું~$0$ વાંચતું હોય અને ડાબે જવાનું હોય
તે કિસ્સો અહીં આવતો નથી (અને તે ડાબે જઈ પણ ન શકે,
એટલે એ જ સ્થાને રહે છે). બીજું સંયોજન એ ખાસ કિસ્સો
આવરે છે: $y$ પગલાં પછી મશીન અવસ્થા~$q_i$માં
ખાનું~$0$ વાંચે અને ખાનું~$0$માં ચિહ્ન~$\sigma$
હોય, તો $y+1$ પગલાં પછી પણ તે ખાનું~$0$ વાંચે છે,
હવે અવસ્થા~$q_j$માં છે, ખાનું~$0$માં ચિહ્ન~$\sigma'$
છે, અને ખાનું~$0$ સિવાયનાં ખાનાંમાં~$y$ પગલાં
પછી હતાં તે જ ચિહ્નો રહે છે.
\item \ollabel{rep-stay} દરેક સૂચના
  $\delta(q_i, \sigma) =
  \tuple{q_j, \sigma', \TMstay}$ માટે વાક્ય:
\begin{align*}
& \lforall[x][\lforall[y][(
   (\Obj Q_{q_i}(x, y) \land \Obj S_{\sigma}(x, y)) \lif {}]] \\
&\qquad   (\Obj Q_{q_j}(x, y') \land \Obj S_{\sigma'}(x, y') \land
!A(x, y)))
\end{align*}
\end{enumerate}
\end{enumerate}
ટ્યુરિંગ મશીન~$M$ અને ઇનપુટ~$w$ માટે ઉપરનાં
બધાં વાક્યોનું સંયોજન~$!T(M, w)$ માનો.

~$M$ છેવટે અટકે છે એમ વ્યક્ત કરવા એવું વાક્ય
જોઈએ જે કહે કે ``કેટલાંક પગલાં પછી સંક્રમણ વિધેય
અવ્યાખ્યાયિત હશે.'' $X$ એ એવાં બધા
$\tuple{q, \sigma}$ જોડોનો ગણ માનો કે જેમના માટે
~$\delta(q, \sigma)$ અવ્યાખ્યાયિત છે. હવે~$!E(M, w)$
એ નીચેનું વાક્ય માનો:
\[
\lexists[x][\lexists[y][(\bigvee_{\tuple{q, \sigma} \in
      X}(\Obj Q_q(x, y) \land \Obj S_\sigma(x, y)))]]
\]

ખાસ અટકવાની અવસ્થા~$h$ ધરાવતું ટ્યુરિંગ મશીન
વાપરીએ, તો વાત વધુ સરળ બને છે: તે વખતે
વાક્ય~$!E(M, w)$
\[
\lexists[x][\lexists[y][\Obj Q_h(x, y)]]
\]
મશીન છેવટે અટકે છે એમ વ્યક્ત કરે છે.

\begin{prop}
\ollabel{prop:mlessk}
જો $m < k$, તો $!T(M, w) \Entails \num{m} < \num{k}$.
\end{prop}

\begin{proof}
કસરત.
\end{proof}

\begin{prob}
\olref[tur][und][rep]{prop:mlessk} સાબિત કરો.
(સૂચન: $k-m$ પર આગમન વાપરો.)
\end{prob}
% END OLP-0270
\endgroup


\begingroup
\makeatletter\def\ol@sectioncs{section}\makeatother

% BEGIN OLP-0271
\olfileid{tur}{und}{ver}
\olsection{નિરૂપણની ચકાસણી}
\phantomsection\label{guunit:OLP-0271}


\begin{explain}
આપણું નિરૂપણ કામ કરે છે તેની ચકાસણી માટે બે વાત
સાબિત કરવી પડશે. પ્રથમ, $M$ ઇનપુટ~$w$ પર અટકે,
તો $!T(M, w) \lif !E(M, w)$ પ્રામાણ્ય ધરાવે છે તે
બતાવવું પડશે. પછી ઊલટી દિશા: જો
$!T(M, w) \lif !E(M, w)$ પ્રામાણ્ય ધરાવે, તો
ઇનપુટ~$w$ પર ચલાવતાં $M$ ખરેખર છેવટે અટકે છે.

આ બે પરિણામો સાબિત કરવાની રીતો એકદમ જુદી છે.
પહેલા પરિણામ માટે પ્રથમ-ક્રમના તર્કનું વાક્ય
(એટલે $!T(M, w) \lif !E(M, w)$) પ્રામાણ્ય ધરાવે છે
તે બતાવવું છે. તેનો સૌથી સરળ રસ્તો નિષ્પત્તિ
આપવાનો છે. પરંતુ આપણી સાબિતી બધાં~$M$ અને~$w$
માટે કામ કરવી જોઈએ; એટલે એક જ વાક્યની
નિષ્પત્તિ આપવાની નથી, અસંખ્ય વાક્યો માટે
આપવાની છે. તેથી $M$ અને~$w$ ગમે તે હોય અને
ઇનપુટ~$w$ પર અટકતાં~$M$ને ગમે તેટલાં પગલાં લાગતાં
હોય, $!T(M, w) \lif !E(M, w)$ની નિષ્પત્તિ
મળશે એવું આગમનથી સાબિત કરી શકીએ.

સ્વાભાવિક રીતે આગમનનો આધાર~$M$ અટકવાના સંરૂપણ
સુધી પહોંચે તે પહેલાંના પગલાંની સંખ્યા હશે. આ
સાબિતીમાં આપણે સ્થાપીશું કે ઇનપુટ~$w$ પરના~$M$ના
સંગણનના દરેક પગલા~$n$ માટે
$!T(M, w) \Entails !C(M, w, n)$, જ્યાં
$!C(M, w, n)$ એ ઇનપુટ~$w$ પર ચાલતા~$M$નું
$n$ પગલાં પછીનું સંરૂપણ યોગ્ય રીતે વર્ણવે છે.
હવે માનો કે $M$ ઇનપુટ~$w$
પર~$n$ પગલાં પછી અટકે છે. ત્યારે
$!C(M, w, n)$ અટકવાનું સંરૂપણ વર્ણવશે. જ્યારે
$!C(M, w, n)$ અટકવાનું સંરૂપણ વર્ણવે, ત્યારે
$!C(M, w, n) \Entails !E(M, w)$ પણ બતાવીશું.
  એટલે $M$ ઇનપુટ~$w$ પર અટકે, તો કોઈ~$n$ માટે
$n$ પગલાં પછી~$M$ અટકવાના સંરૂપણમાં હશે. તેથી
$!T(M, w) \Entails !C(M, w, n)$, જ્યાં
$!C(M, w, n)$ એવું સંરૂપણ વર્ણવે છે; અને એ કિસ્સામાં
$!C(M, w, n) \Entails !E(M, w)$ હોવાથી
$!T(M, w) \Entails !E(M, w)$ મળે છે, એટલે
$\Entails !T(M, w) \lif !E(M, w)$.
\sourcecorrection{OLUN-006}{સ્થિર અંગ્રેજી મૂળની આ
છેલ્લી ફલિતતા-પંક્તિમાં વ્યાખ્યાયિત $!T(M,w)$ને
બદલે $T(M,w)$ છપાયું છે; અહીં વ્યાખ્યાયિત
વાક્યનું ચિહ્ન જાળવ્યું છે.}

ઊલટી દિશાની રીત તદ્દન જુદી છે. અહીં આપણે
$\Entails !T(M, w) \lif !E(M, w)$ માનીને
$M$ ઇનપુટ~$w$ પર અટકે છે તે સાબિત કરવાનું છે.
ધારણામાંથી $!T(M, w) \Entails !E(M, w)$ મળે છે:
એટલે જ્યાં $!T(M, w)$ સત્ય છે એવી દરેક
સંરચનામાં $!E(M, w)$ પણ સત્ય છે. તેથી
જેમાં $!T(M, w)$ સત્ય છે એવી
સંરચના~$\Struct{M}$ વર્ણવીશું. તેનો પ્રદેશ
$\Nat$ હશે; બધાં $\Obj Q_q$ અને $\Obj S_\sigma$
નું અર્થઘટન ઇનપુટ~$w$ પર ચાલતા~$M$નાં
સંરૂપણો પરથી નક્કી થશે. દાખલા તરીકે,
$\Sat{M}{\Obj Q_q(\num{m}, \num{n})}$ ત્યારે અને માત્ર
ત્યારે જ્યારે~$M$ ઇનપુટ~$w$ પર~$n$ પગલાં ચાલ્યા
પછી અવસ્થા~$q$માં હોય અને ખાનું~$m$ વાંચતું હોય.
\sourcecorrection{OLUN-007}{સ્થિર અંગ્રેજી મૂળના આ
દાખલામાં ચાલતા મશીનને $T$ લખ્યું છે; સંદર્ભમાં તે
$M$ છે, જ્યારે $!T(M,w)$ મશીનનું વર્ણન કરતું વાક્ય છે.}
હવે ધારણા મુજબ $!T(M, w) \Entails !E(M, w)$,
અને રચના મુજબ $\Sat{M}{!T(M, w)}$;
તેથી $\Sat{M}{!E(M, w)}$. પરંતુ
$\Sat{M}{!E(M, w)}$ ત્યારે અને માત્ર ત્યારે જ્યારે
કોઈ $n \in \Domain{M} = \Nat$ માટે ઇનપુટ~$w$
પર ચાલતું~$M$,~$n$ પગલાં પછી અટકવાના
સંરૂપણમાં હોય.
\end{explain}

\begin{defn}
$!C(M, w, n)$ એ નીચેનું વાક્ય માનો:
\[
\Obj Q_q(\num{m}, \num{n}) \land \Obj S_{\sigma_0}(\num{0}, \num{n})
\land \dots \land \Obj S_{\sigma_k}(\num{k}, \num{n}) \land
\lforall[x][(\num{k} < x \lif \Obj S_\TMblank(x, \num{n}))]
\]
અહીં~$q$, સમય~$n$ પર~$M$ની અવસ્થા છે;
સમય~$n$ પર~$M$ ખાનું~$m$ વાંચે છે;
$0 \le i \le k$ માટે ખાનું~$i$ સમય~$n$ પર
ચિહ્ન~$\sigma_i$ ધરાવે છે;
અને~$k$ એ સમય~$0$ પરની ટેપનું સૌથી જમણું ખાલી
ન હોય એવું ખાનું તથા~$n$ પગલાં પછી હેડે મુલાકાત
લીધેલું સૌથી જમણું ખાનું—આ બેમાંથી વધુ જમણાં
ખાનાનો સૂચક છે.
\end{defn}

\begin{lem}\ollabel{lem:halt-config-implies-halt}
ઇનપુટ~$w$ પર ચાલતું~$M$,~$n$ પગલાં પછી
અટકવાના સંરૂપણમાં હોય, તો
$!C(M, w, n) \Entails !E(M, w)$.
\end{lem}

\begin{proof}
માનો કે ઇનપુટ~$w$ પર~$M$,~$n$ પગલાં પછી અટકે છે.
કોઈ અવસ્થા~$q$, ખાનું~$m$ અને ચિહ્ન~$\sigma$
એવાં છે કે:
\begin{enumerate}
\item $n$ પગલાં પછી~$M$ અવસ્થા~$q$માં ખાનું~$m$
  વાંચે છે, અને તે ખાનામાં~$\sigma$ છે.
\item સંક્રમણ વિધેય~$\delta(q, \sigma)$ અવ્યાખ્યાયિત છે.
\end{enumerate}
$!C(M, w, n)$ આ સંરૂપણનું વર્ણન છે અને તેમાં
$\Obj Q_{q}(\num{m}, \num{n})$ તથા
$\Obj S_{\sigma}(\num{m}, \num{n})$ ખંડો હશે.
આ બંને ખંડો મળીને~$!E(M, w)$ ફલિત કરે છે:
\[
\lexists[x][\lexists[y][(\bigvee_{\tuple{q, \sigma} \in
      X}(\Obj Q_q(x, y) \land \Obj S_\sigma(x, y)))]]
\]
કારણ કે $\tuple{q,\sigma} \in X$ હોવાથી
$\Obj Q_q(\num{m}, \num{n}) \land
\Obj S_{\sigma}(\num{m}, \num{n}) \Entails
\bigvee_{\tuple{q, \sigma} \in X}
(\Obj Q_q(\num{m}, \num{n}) \land
\Obj S_{\sigma}(\num{m}, \num{n}))$.
\sourcecorrection{OLUN-008}{સ્થિર અંગ્રેજી મૂળની
આ ફલિતતા-પંક્તિમાં પહેલાના પસંદ કરેલા $q,\sigma$
બદલે અચાનક $q',\sigma'$ આવે છે અને એક
$S$-વિધેય આગળનો વસ્તુભાષાનો અક્ષરપ્રકાર આપતો આદેશ છૂટી ગયો છે. અહીં
અગાઉના જ $q,\sigma$ અને પૂરો વિધેય-સંકેત રાખ્યો છે.}
\end{proof}

\begin{explain}
એટલે $M$ ઇનપુટ~$w$ પર અટકે, તો કોઈ~$n$ માટે
$!C(M, w, n) \Entails !E(M,w)$. હવે આપણે
બતાવીશું કે કોઈ પણ સમય~$n$ માટે
$!T(M, w) \Entails !C(M, w, n)$.
\end{explain}

\begin{lem}
\ollabel{lem:config}
દરેક~$n$ માટે, $M$ જો~$n$ પગલાં પૂરાં થાય તે
પહેલાં અટક્યું ન હોય, તો
$!T(M, w) \Entails !C(M, w, n)$.
\sourcecorrection{OLUN-009}{સ્થિર અંગ્રેજી મૂળનું
``$n$ પગલાં પછી અટક્યું ન હોય'' વિધાન~$n$મા
પગલે જ અટકવાનું સંરૂપણ બાકાત કરે છે, છતાં નીચેનો
પ્રામાણ્યનો પુરાવો એ જ સંરૂપણ માટે આ ઉપપ્રમેય
વાપરે છે. અહીં~$n$ પગલાં પૂરાં થતાં પહેલાં ન
અટકવાની જરૂરી શરત સ્પષ્ટ કરી છે.}
\end{lem}

\begin{proof}
આગમનનો આધાર: $n = 0$ હોય, તો
$!C(M, w, 0)$ના સંયોજકો~$!T(M,w)$નાં
પ્રારંભિક સંરૂપણ અને સંખ્યાનાં સ્વયંસિદ્ધ વાક્યોમાંથી
ફલિત થાય છે. ખાલી ઇનપુટમાં શરૂઆતથી વાંચાતું
ખાનું ખાલી છે, અને તેની જમણી બાજુનાં બધાં
ખાનાં પણ ખાલી છે, તે ખાલી-ટેપની શરત અને
અનુગામીના સ્વયંસિદ્ધ વાક્યથી મળે છે.
\sourcecorrection{OLUN-024}{સ્થિર અંગ્રેજી મૂળ
આધારકિસ્સાના બધા સંયોજકો~$!T(M,w)$માં
સીધા જ છે એમ કહે છે. ખાલી ઇનપુટ હોય ત્યારે
માથું શરૂઆતમાં પહેલું ખાનું વાંચે છે, એટલે
તે ખાનું ખાલી હોવાનું અને તેની આગળની
ખાલી-ટેપ શરતનું અનુમાન પણ કરવું પડે છે.}

આગમનની ધારણા: $M$ જો~$n$મા પગલા પહેલાં
અટક્યું ન હોય, તો $!T(M,w) \Entails !C(M, w, n)$.
હવે બતાવવું છે કે, $!C(M, w, n)$ અટકવાનું સંરૂપણ
વર્ણવતું ન હોય તો,
$!T(M, w) \Entails !C(M, w, n+1)$.

માનો કે $n \ge 0$ અને ઇનપુટ~$w$ પર શરૂ થયેલું
$M$,~$n$ પગલાં પછી અવસ્થા~$q$માં ખાનું~$m$
વાંચે છે. $M$,~$n$ પગલાં પછી અટકતું ન હોવાથી
~$M$ના કાર્યક્રમમાં નીચેના ત્રણ સ્વરૂપોમાંની કોઈ એક
સૂચના હોવી જોઈએ:
\sourcecorrection{OLUN-010}{સ્થિર અંગ્રેજી મૂળ આ
આગમન-પગલામાં $n>0$ ધારે છે, જેથી આધાર પછીનું
$n=0$થી $n+1$ પગલું છૂટી જાય છે. એ જ વિચાર
$n=0$ માટે પણ લાગુ પડે છે, તેથી અહીં $n\ge0$ છે.}

\begin{enumerate}
\item \ollabel{right} $\delta(q, \sigma) = \tuple{q', \sigma', \TMright}$

\item \ollabel{left} $\delta(q, \sigma) = \tuple{q', \sigma', \TMleft}$

\item \ollabel{stay} $\delta(q, \sigma) = \tuple{q', \sigma', \TMstay}$
\end{enumerate}

આ ત્રણેય કિસ્સા અનુક્રમે જોઈએ.

\begin{enumerate}
\item માનો કે સૂચના~\olref{right}ના સ્વરૂપની છે.
  \olref[rep]{defn:tm-descr}\olref[rep]{rep-right}
  મુજબ આનો અર્થ છે કે
\begin{align*}
& \lforall[x][\lforall[y][((\Obj Q_{q}(x,
  y) \land \Obj S_\sigma(x, y)) \lif {}]]\\
& \qquad (\Obj Q_{q'}(x',
  y') \land \Obj S_{\sigma'}(x, y') \land !A(x, y)))
\intertext{એ $!T(M,w)$નો એક સંયોજક છે. તેથી
  સર્વપરિમાણકના ઉદાહરણથી (અહીં~$x$ માટે
  $\num{m}$ અને~$y$ માટે $\num{n}$ મૂકીને)
  નીચેનું વાક્ય ફલિત થાય છે:}
& (\Obj Q_{q}(\num{m}, \num{n}) \land \Obj S_{\sigma}(\num{m},
\num{n})) \lif {}\\
& \qquad (\Obj Q_{q'}(\num{m}', \num{n}') \land
\Obj S_{\sigma'}(\num{m}, \num{n}') \land !A(\num{m}, \num{n})).
\intertext{આગમનની ધારણા મુજબ $!T(M, w) \Entails !C(M, w, n)$,
  એટલે કે:}
& \Obj Q_q(\num{m}, \num{n}) \land \Obj S_{\sigma_0}(\num{0}, \num{n})
\land \dots \land \Obj S_{\sigma_k}(\num{k}, \num{n}) \land \\
& \qquad\lforall[x][(\num{k} < x \lif \Obj S_\TMblank(x, \num{n}))]\\
\intertext{$n$ પગલાં પછી ખાનું~$m$ ચિહ્ન~$\sigma$
  ધરાવે છે, તેથી અનુરૂપ સંયોજક
  $\Obj S_\sigma(\num{m}, \num{n})$ છે; એટલે
  તેમાંથી આ ફલિત થાય છે:}
& \Obj Q_{q}(\num{m}, \num{n}) \land \Obj S_{\sigma}(\num{m},
   \num{n})
\intertext{હવે આપણને મળે છે:}
& \Obj Q_{q'}(\num{m}', \num{n}') \land \Obj S_{\sigma'}(\num{m},
  \num{n}') \land {}\\
& \qquad\bigwedge_{\substack{0\le i\le k\\i\ne m}}
  \Obj S_{\sigma_i}(\num{i}, \num{n}') \land {}\\
& \qquad \lforall[x][(\num{k} < x \lif \Obj S_\TMblank(x, \num{n}'))]
\end{align*}
\sourcecorrection{OLUN-022}{સ્થિર અંગ્રેજી મૂળના
જમણી-ચાલના પરિણામમાં બદલાયેલા ખાના~$m$નું
નવું ચિહ્ન પહેલેથી અલગ આપેલું છે, છતાં પછીની
યાદી દેખીતી રીતે $0$થી~$k$ સુધીના બધા
જૂનાં ચિહ્નો ફરી જોડે છે. મૂળનું જ ગદ્ય~$m$નો
બાકાત રાખે છે; અહીં એ બાકાત સ્પષ્ટ લખ્યો છે.}
આ રીતે: પહેલી પંક્તિ અગાઉના શરતી વાક્યના
અનુગામીમાંથી સીધી, મોડસ પોનેન્સથી મળે છે.
વચ્ચેની પંક્તિનો દરેક સંયોજક—જેમાં
$S_{\sigma_m}(\num{m},\num{n}')$નો સમાવેશ નથી—
$!C(M, w, n)$ના અનુરૂપ સંયોજક તથા
$!A(\num{m}, \num{n})$માંથી ફલિત થાય છે.

$m < k$ હોય, તો $!T(M,w) \Proves \num{m} < \num {k}$
(\olref[rep]{prop:mlessk}); અને~$<$ની સંક્રામકતાથી
$\lforall[x][(\num{k} < x \lif \num{m} < x)]$
મળે છે. $m = k$ હોય, તો એ જ વાક્ય
$\lforall[x][(\num{k} < x \lif \num{m} < x)]$
માત્ર તર્કથી મળે છે. તેથી છેલ્લી પંક્તિ
$!C(M, w, n)$ના અનુરૂપ સંયોજક,
$\lforall[x][(\num{k} < x \lif \num{m} < x)]$,
અને $!A(\num{m}, \num{n})$માંથી ફલિત થાય છે.
$m<k$ હોય, તો આ જ $!C(M, w, n+1)$ છે.

હવે માનો કે $m=k$. ત્યારે $n+1$ પગલાં પછી
ટેપ-હેડે ખાનું~$k+1$ પણ મુલાકાત લીધું હશે;
તે હવે મુલાકાત લેવાયેલું સૌથી જમણું ખાનું છે.
એટલે $!C(M, w, n+1)$માં નવો સંયોજક
$\Obj S_\TMblank(\num{k}',\num{n}')$ છે, અને
છેલ્લો સંયોજક
$\lforall[x][(\num{k}' < x \lif \Obj S_\TMblank(x, \num{n}'))]$
છે. આ બે વાક્યો પણ ફલિત થાય છે તે
ચકાસવું પડશે.

આપણી પાસે પહેલેથી
$\lforall[x][(\num{k} < x \lif \Obj S_\TMblank(x,
  \num{n}'))]$ છે. ખાસ કરીને તેમાંથી
$\num{k} < \num{k}' \lif
\Obj S_\TMblank(\num{k}', \num{n}')$ મળે છે.
સ્વયંસિદ્ધ વાક્ય $\lforall[x][x < x']$માંથી
$\num{k} < \num{k}'$ મળે છે. મોડસ પોનેન્સથી
$\Obj S_\TMblank(\num{k}',\num{n}')$ ફલિત થાય છે.

વળી, $!T(M,w) \Proves \num{k} < \num{k}'$
હોવાથી~$<$ની સંક્રામકતાના સ્વયંસિદ્ધ વાક્યથી
$\lforall[x][(\num{k}' < x \lif \Obj
  S_\TMblank(x, \num{n}'))]$ મળે છે. (આ
ચકાસણી કસરત તરીકે છોડી છે.)

\item માનો કે સૂચના~\olref{left}ના સ્વરૂપની છે.
  તો \olref[rep]{defn:tm-descr}\olref[rep]{rep-left}
  મુજબ:
\sourcecorrection{OLUN-011}{સ્થિર અંગ્રેજી મૂળના
આ ડાબે-ચાલ કિસ્સામાં લખાતું ખાનું~$x'$ અને
ઉદાહરણમાં~$l+1$ છે, પરંતુ યથાવત્-ખાનાંની
શરત અનુક્રમે $!A(x,y)$ અને $!A(\num{l},\num{n})$
લખાઈ છે. ઉપરાંત શૂન્ય-ખાનાની શાખામાં અહીં
પસંદ કરેલા $q,q'$ને બદલે જૂનાં $q_i,q_j$ છે.
નીચે લખાતું ખાનું બાકાત રાખ્યું છે અને
અવસ્થા-ચિહ્નો એકરૂપ કર્યા છે.}
\begin{align*}
& \lforall[x][\lforall[y][((\Obj Q_{q}(x', y) \land \Obj
    S_{\sigma}(x', y)) \lif {}]]\\
& \qquad (\Obj Q_{q'}(x, y') \land \Obj
  S_{\sigma'}(x', y') \land !A(x', y))) \land {}\\
& \lforall[y][((\Obj Q_{q}(\num{0}, y) \land \Obj S_{\sigma}(\num{0},
    y)) \lif {}]\\
& \qquad (\Obj Q_{q'}(\num{0}, y') \land \Obj S_{\sigma'}(\num{0}, y')
  \land !A(\num{0}, y)))
\intertext{એ $!T(M,w)$નો એક સંયોજક છે. $m>0$
  હોય, તો $l = m - 1$ લો (એટલે $m = l+1$).
  ઉપરના વાક્યના પહેલા સંયોજકમાંથી
  નીચેનું ફલિત થાય છે:}
& (\Obj Q_{q}(\num{l}', \num{n}) \land \Obj S_{\sigma}(\num{l}', \num{n}))
\lif {} \\
& \qquad (\Obj Q_{q'}(\num{l}, \num{n}') \land \Obj S_{\sigma'}(\num{l}',
\num{n}') \land !A(\num{l}', \num{n}))
\intertext{અન્યથા $l = m = 0$ લો અને બીજા
  સંયોજકમાંથી ફલિત થતું નીચેનું વાક્ય
  જુઓ:}
& ((\Obj Q_{q}(\num{0}, \num{n}) \land \Obj S_{\sigma}(\num{0}, \num{n})) \lif {}\\
& \qquad   (\Obj Q_{q'}(\num{0}, \num{n}') \land \Obj S_{\sigma'}(\num{0}, \num{n}') \land
!A(\num{0}, \num{n})))
\intertext{બંનેમાંથી જે લાગુ પડે તેમાંથી
નીચેનું ફલિત થાય છે:}
&  \Obj Q_{q'}(\num{l}, \num{n}') \land \Obj S_{\sigma'}(\num{m},
  \num{n}') \land {}\\
&\qquad  \bigwedge_{\substack{0\le i\le k\\i\ne m}}
  \Obj S_{\sigma_i}(\num{i}, \num{n}') \land {} \\
&\qquad  \lforall[x][(\num{k} < x
  \lif \Obj S_\TMblank(x, \num{n}'))]
\end{align*}
\sourcecorrection{OLUN-023}{સ્થિર અંગ્રેજી મૂળના
ડાબી-ચાલના પરિણામમાં પણ લખાયેલું ખાનું~$m$
અલગ નવા ચિહ્ન સાથે આવે છે; પછીની જૂનાં
ચિહ્નોની યાદીમાં તેને ફરી સામેલ કરવાથી
સંરૂપણ ખોટું થાય. અહીં તે ખાનું સ્પષ્ટ બાકાત છે.}
અગાઉની જેમ. (પહેલા કિસ્સામાં
$\num{l}' \ident \num{l+1} \ident \num{m}$
અને બીજા કિસ્સામાં $\num{l} \ident \num{0}$
એ નોંધો.) પરંતુ આ તો $!C(M, w, n+1)$ જ છે.

\item \olref{stay} કિસ્સો કસરત તરીકે છોડી છે.
\end{enumerate}
આપણે બતાવ્યું કે આવા દરેક~$n$ માટે
$!T(M, w) \Entails !C(M, w, n)$.
\end{proof}

\begin{prob}
\olref[tur][und][ver]{lem:config}ની સાબિતીમાં
કિસ્સો~\olref[tur][und][ver]{stay} પૂર્ણ કરો.
\end{prob}

\begin{prob}
$\Obj S_{\sigma_i}(\num{i}, \num{n})$ અને
$!A(m, n)$માંથી
$\Obj S_{\sigma_i}(\num{i}, \num{n}')$ની
નિષ્પત્તિ આપો (ધારો કે $i \neq m$,
એટલે $i < m$ અથવા $m < i$).
\end{prob}

\begin{prob}
$\lforall[x][(\num{k} < x \lif \Obj
  S_\TMblank(x, \num{n}'))]$,
$\lforall[x][x < x']$ અને
$\lforall[x][\lforall[y][\lforall[z][ ((x < y \land y < z) \lif x <
      z)]]]$માંથી
$\lforall[x][(\num{k}' < x \lif \Obj
  S_\TMblank(x, \num{n}'))]$ની
નિષ્પત્તિ આપો.
\end{prob}

\begin{lem}
\ollabel{lem:valid-if-halt}
$M$ ઇનપુટ~$w$ પર અટકે, તો
$!T(M, w) \lif !E(M, w)$ પ્રામાણ્ય ધરાવે છે.
\end{lem}

\begin{proof}
\olref{lem:config}થી જાણીએ છીએ કે પ્રથમ
અટકવાના સમય સુધીનાં દરેક~$n$ માટે
સમય~$n$ પરનાં~$M$ના સંરૂપણનું
વર્ણન~$!C(M, w, n)$,
$!T(M, w)$માંથી ફલિત થાય છે.
\sourcecorrection{OLUN-012}{સ્થિર અંગ્રેજી મૂળ
અહીં ``કોઈ પણ સમય~$n$'' કહે છે, છતાં મશીન
અટક્યા પછીનું સંરૂપણ તેની વ્યાખ્યા મુજબ આપેલું
નથી. OLUN-009ના સુધારેલા ઉપપ્રમેયને પ્રથમ
અટકવાના સમય સુધી જ વાપર્યો છે.}
માનો કે $M$,~$k$ પગલાં પછી અટકે છે. તે સમયે
તે કોઈ~$m \in \Nat$ માટે ખાનું~$m$ વાંચતું હશે.
તેથી $!C(M, w, k)$ એ~$M$ના અટકવાના
સંરૂપણનું વર્ણન છે: તેમાં
$\Obj Q_q(\num{m}, \num{k})$ તથા
$\Obj S_\sigma(\num{m}, \num{k})$ બંને સંયોજકો છે,
અને $\delta(q,\sigma)$ અવ્યાખ્યાયિત છે. તેથી
\olref{lem:halt-config-implies-halt} મુજબ
$!C(M, w, k) \Entails !E(M, w)$. વળી
$!T(M, w) \Entails !C(M, w, k)$ હોવાથી
$!T(M, w) \Entails !E(M, w)$; એટલે
$!T(M, w) \lif !E(M, w)$ પ્રામાણ્ય ધરાવે છે.
\end{proof}

\begin{explain}
આપણો દાવો ચકાસવાનું પૂરું કરવા ઊલટી દિશા પણ
સ્થાપવી પડશે: $!T(M, w) \lif !E(M, w)$
પ્રામાણ્ય ધરાવે, તો ઇનપુટ~$w$ પર શરૂ કરતાં
$M$ ખરેખર અટકે છે.
\end{explain}

\begin{lem}
\ollabel{lem:halt-if-valid}
$\Entails !T(M, w) \lif !E(M, w)$ હોય, તો
$M$ ઇનપુટ~$w$ પર અટકે છે.
\end{lem}

\begin{proof}
પ્રદેશ~$\Nat$ ધરાવતું $\Lang L_M$-સંરચના
$\Struct M$ લો. તેમાં $\Obj 0$નું અર્થઘટન~$0$,
$\prime$નું અનુગામી વિધેય અને~$<$નું નાની-કરતાં
સંબંધ તરીકે કરીએ છીએ. વિધાનો $\Obj Q_q$ અને
$\Obj S_\sigma$નું અર્થઘટન આ પ્રમાણે છે:
\begin{align*}
  \Assign{\Obj Q_q}{M} & =
\Setabs{\tuple{m, n}}{\begin{array}{ll}\text{ઇનપુટ~$w$ પર શરૂ કરતાં,~$n$ પગલાં પછી,}\\ \text{$M$ અવસ્થા~$q$માં
  ખાનું~$m$ વાંચે છે}\end{array}} \\
  \Assign{\Obj S_\sigma}{M} & = \Setabs{\tuple{m, n}}{\begin{array}{ll}
\text{ઇનપુટ~$w$ પર શરૂ કરતાં,~$n$ પગલાં પછી,}\\ \text{$M$ના ખાના~$m$માં
  ચિહ્ન~$\sigma$ છે}\end{array}}
\end{align*}
બીજા શબ્દોમાં, ઇનપુટ~$w$ પર શરૂ થયેલું~$M$
હકીકતમાં પગલે પગલે શું કરે છે તેનું વર્ણન કરે
એવું સંરચના~$\Struct{M}$ રચીએ છીએ. સ્પષ્ટપણે
$\Sat{M}{!T(M, w)}$. જો
$\Entails !T(M, w) \lif !E(M, w)$, તો
$\Sat{M}{!E(M, w)}$ પણ; એટલે:
\[
\Sat{M}{\lexists[x][\lexists[y][(\bigvee_{\tuple{q, \sigma} \in
      X}(\Obj Q_q(x, y) \land \Obj S_\sigma(x, y)))]]}.
\]
$\Domain{M} = \Nat$ હોવાથી કોઈ
$m$, $n \in \Nat$ તથા એવાં~$q$ અને~$\sigma$
છે કે $\Sat{M}{\Obj Q_q(\num{m}, \num{n}) \land
\Obj S_\sigma(\num{m}, \num{n})}$ અને
$\delta(q, \sigma)$ અવ્યાખ્યાયિત છે.
$\Struct M$ની વ્યાખ્યા મુજબ તેનો અર્થ એ છે કે
ઇનપુટ~$w$ પર શરૂ થયેલું~$M$,~$n$ પગલાં પછી
અવસ્થા~$q$માં છે અને ચિહ્ન~$\sigma$ વાંચે છે;
સંક્રમણ વિધેય અવ્યાખ્યાયિત છે, એટલે~$M$ અટક્યું છે.
\end{proof}
% END OLP-0271
\endgroup


\begingroup
\makeatletter\def\ol@sectioncs{section}\makeatother

% BEGIN OLP-0272
\olfileid{tur}{und}{uns}
\olsection{નિર્ણય સમસ્યા ઉકેલી શકાતી નથી}
\phantomsection\label{guunit:OLP-0272}


\begin{thm}
\ollabel{thm:decision-prob}
નિર્ણય સમસ્યા ઉકેલી શકાતી નથી: એવું કોઈ ટ્યુરિંગ
મશીન~$D$ નથી કે જેની ટેપ પર પ્રથમ-ક્રમના તર્કનું
વાક્ય~$!B$ ઇનપુટ તરીકે હોય ત્યારે~$D$ છેવટે
અટકે અને~$!B$ પ્રામાણ્ય ધરાવે ત્યારે અને માત્ર ત્યારે
જ~$1$, અન્યથા~$0$ આઉટપુટ આપે.
\end{thm}

\begin{proof}
માનો કે નિર્ણય સમસ્યા ઉકેલી શકાતી હોય, એટલે એવું
ટ્યુરિંગ મશીન~$D$ હોય. તો અટકવાની સમસ્યા આ પ્રમાણે
ઉકેલી શકીએ. ટ્યુરિંગ મશીન~$M_e$નો સૂચક~$e$ અને
ઇનપુટ~$w$ આપતાં, તેને અનુરૂપ વાક્ય
$!T(M_e, w) \lif !E(M_e, w)$ ગણતું અને ટેપના સૌથી
ડાબા ખાનાને વાંચતું અટકતું ટ્યુરિંગ મશીન~$E$ બનાવીએ.
પછી $E \concat D$ મશીન ઇનપુટ~$e$ અને~$w$ માટે
પહેલાં $!T(M_e, w) \lif !E(M_e, w)$ ગણે અને ત્યાર પછી
તે ઇનપુટ પર નિર્ણય સમસ્યાનું મશીન~$D$ ચલાવે.
$!T(M_e, w) \lif !E(M_e, w)$ પ્રામાણ્ય ધરાવે ત્યારે
અને માત્ર ત્યારે $D$ આઉટપુટ~$1$ સાથે અટકે છે;
અન્યથા આઉટપુટ~$0$ આપે છે.
\olref[ver]{lem:halt-if-valid} અને
\olref[ver]{lem:valid-if-halt} મુજબ
$!T(M_e, w) \lif !E(M_e, w)$ પ્રામાણ્ય
ધરાવે ત્યારે અને માત્ર ત્યારે~$M_e$ ઇનપુટ~$w$ પર
અટકે છે. એટલે ઇનપુટ~$e$ અને~$w$ પર $E\concat D$
એવું મશીન બને કે~$M_e$ ઇનપુટ~$w$ પર અટકે ત્યારે
અને માત્ર ત્યારે આઉટપુટ~$1$ સાથે, અન્યથા આઉટપુટ~$0$
સાથે અટકે. બીજા શબ્દોમાં, $E \concat D$ અટકવાની
સમસ્યા ઉકેલી નાખે. પરંતુ
\olref[hal]{thm:halting-problem}થી આપણે જાણીએ છીએ
કે આવું કોઈ ટ્યુરિંગ મશીન હોઈ શકતું નથી.
\end{proof}

\begin{cor}\ollabel{cor:undecidable-sat}%
પ્રથમ-ક્રમના તર્કનું મનસ્વી વાક્ય સંતોષ્ય છે કે
નહીં તે નક્કી કરી શકાતું નથી.
\end{cor}

\begin{proof}
  માનો કે સંતોષ્યતા ટ્યુરિંગ મશીન~$S$ વડે નક્કી કરી
  શકાતી હોય. તો નિર્ણય સમસ્યા આ રીતે ઉકેલી શકીએ:
  વાક્ય~$!B$ ઇનપુટ તરીકે આપેલું હોય, તો
  ટેપ પરનું~$!B$ એક ખાનું જમણી તરફ ખસેડો.
  \sourcecorrection{OLUN-005}{સ્થિર અંગ્રેજી મૂળ આ
  વાક્યમાં ઇનપુટને $B$ કહે છે; આસપાસની સાબિતીમાં
  તે જ વાક્ય $!B$ છે, તેથી અહીં ચિહ્ન એકરૂપ કર્યું છે.}
  ખાનું~$1$ પર પાછા જઈ~$\lnot$ ચિહ્ન લખો.

  હવે ટ્યુરિંગ મશીન~$S$ ચલાવો. ટેપ પર આઉટપુટ~$1$
  (જો $\lnot !B$ સંતોષ્ય હોય) અથવા~$0$ (જો
  $\lnot !B$ અસંતોષ્ય હોય) સાથે તે છેવટે અટકે છે.
  ખાનું~$1$ પર~$\TMstroke$ હોય, તો તેને ભૂંસી નાખો;
  ખાનું~$1$ ખાલી હોય, તો ત્યાં~$\TMstroke$ લખો,
  અને પછી અટકો.

  આ ટ્યુરિંગ મશીન હંમેશાં અટકે છે અને~$\lnot !B$
  અસંતોષ્ય હોય ત્યારે અને માત્ર ત્યારે આઉટપુટ~$1$,
  અન્યથા~$0$ આપે છે. $!B$ પ્રામાણ્ય ધરાવે ત્યારે
  અને માત્ર ત્યારે~$\lnot !B$ અસંતોષ્ય છે. એટલે
  મશીન~$!B$ પ્રામાણ્ય ધરાવે ત્યારે અને માત્ર ત્યારે
  આઉટપુટ~$1$, અન્યથા~$0$ આપે છે; આમ તે નિર્ણય
  સમસ્યા ઉકેલી નાખે.
\end{proof}

\begin{explain}
એટલે ``$!B$ પ્રથમ-ક્રમના તર્કનું પ્રમાણભૂત વાક્ય
છે?'' એ પ્રશ્નનો હંમેશાં સાચો ``હા'' અથવા ``ના''
જવાબ આપતું ટ્યુરિંગ મશીન નથી. પરંતુ સાચો જવાબ
``હા'' હોય ત્યારે તે આપતું, અને જવાબ ``ના'' હોય તો
અટકતું જ ન હોય એવું ટ્યુરિંગ મશીન \emph{છે}.
આ વાત પ્રથમ-ક્રમના તર્કની યથાર્થતા અને પૂર્ણતાના
પ્રમેયોમાંથી તથા નિષ્પત્તિઓની અસરકારક પરિગણના
કરી શકાય છે એ હકીકતમાંથી ફલિત થાય છે.
\end{explain}

\begin{thm}
  \ollabel{thm:valid-ce}%
  પ્રથમ-ક્રમનાં વાક્યોનું પ્રામાણ્ય અર્ધનિર્ણેય
  છે: એવું ટ્યુરિંગ મશીન~$E$ છે કે જેની ટેપ પર
  પ્રથમ-ક્રમનું વાક્ય~$!B$ ઇનપુટ તરીકે હોય,
  તો~$!B$ પ્રામાણ્ય ધરાવે ત્યારે અને માત્ર ત્યારે
  $E$ છેવટે આઉટપુટ~$1$ સાથે અટકે છે; અન્યથા તે
  અટકતું નથી.
\end{thm}

\begin{proof}
  પ્રથમ-ક્રમના તર્કની બધી શક્ય નિષ્પત્તિઓ એક
  પછી એક અસરકારક વિધિ વડે ઉત્પન્ન કરી શકાય છે.
  મશીન~$E$ એમ કરે છે અને~$\Proves !B$ બતાવતી
  નિષ્પત્તિ મળે ત્યારે આઉટપુટ~$1$ સાથે અટકે છે.
  યથાર્થતાના પ્રમેયથી, $E$ આઉટપુટ~$1$ સાથે અટકે,
  તો~$\Entails !B$ છે. પૂર્ણતાના પ્રમેયથી,
  $\Entails !B$ હોય તો~$\Proves !B$ બતાવતી
  નિષ્પત્તિ છે. $E$ બધી શક્ય નિષ્પત્તિઓ
  પદ્ધતિસર ઉત્પન્ન કરતું હોવાથી તે
  $\Proves !B$ બતાવતી વ્યૂત્પત્તિ છેવટે
  શોધશે અને આઉટપુટ~$1$
  સાથે અટકશે.
\end{proof}
% END OLP-0272
\endgroup


\begingroup
\makeatletter\def\ol@sectioncs{section}\makeatother

% BEGIN OLP-0273
\olfileid{tur}{und}{tra}
\olsection{ટ્રાખ્ટેનબ્રોટનું પ્રમેય}
\phantomsection\label{guunit:OLP-0273}


\begin{explain}
  \olref[rep]{sec}માં આપણે ટ્યુરિંગ મશીન~$M$
  અને ઇનપુટ શબ્દમાળા~$w$ માટે વાક્યો
  $!T(M,w)$ તથા~$!E(M,w)$ વ્યાખ્યાયિત કર્યા.
  પછી \olref[ver]{lem:valid-if-halt} અને
  \olref[ver]{lem:halt-if-valid}માં બતાવ્યું કે
  $!T(M,w) \lif !E(M,w)$ ત્યારે અને માત્ર ત્યારે
  પ્રામાણ્ય ધરાવે છે જ્યારે ઇનપુટ~$w$ પર શરૂ થયેલું
  $M$ છેવટે અટકે છે. અટકવાની સમસ્યા અનિર્ણેય
  હોવાથી, આમાંથી પ્રથમ-ક્રમના તર્કનાં વાક્યોનું
  પ્રામાણ્ય અને સંતોષ્યતા અનિર્ણેય છે એ ફલિત થાય છે
  (\cref{tur:und:uns:thm:decision-prob,tur:und:uns:cor:undecidable-sat}).

  પરંતુ વાક્યનું પ્રામાણ્ય અને સંતોષ્યતા કોઈ પણ
  સંરચનાઓ માટે વ્યાખ્યાયિત છે—સાન્ત હોય કે
  અનંત. તમને એવું લાગે કે કોઈ વાક્ય
  સાન્ત સંરચનામાં સંતોષ્ય છે કે નહીં (અથવા
  બધાં સાન્ત સંરચનાઓમાં પ્રામાણ્ય ધરાવે છે કે
  નહીં) તે નક્કી કરવું સરળ હશે. નિર્ણય સમસ્યા
  ઉકેલી ન શકાય તેની સાબિતી અનુરૂપ ફેરવીને
  બતાવી શકીએ કે એવું નથી.

  પહેલાં \olref[ver]{lem:halt-if-valid}ની
  સાબિતી ફરી જુઓ: ત્યાં આપણે~$!T(M,w)$નું
  નિદર્શ~$\Struct M$ બનાવ્યું હતું, જે
  ઇનપુટ~$w$ પર શરૂ થયેલું મશીન~$M$ શું કરે
  છે તેનું ચોક્કસ વર્ણન કરે છે. એ નિદર્શનો
  પ્રદેશ~$\Nat$ હોવાથી તે અનંત હતો. પરંતુ
  $M$ ઇનપુટ~$w$ પર ખરેખર અટકતું હોય, તો
  એ જ રીતે સાન્ત નિદર્શ~$\Struct M'$ બનાવી શકીએ.
  માનો કે ઇનપુટ~$w$ પર શરૂ થયેલું~$M$,
  $k$ પગલાં પછી અટકે છે. પ્રદેશ
  $\Domain{M'}$ તરીકે $\{0, \dots, n\}$ લો,
  જ્યાં~$n$ એ~$k+1$ અને~$w$ની લંબાઈમાંથી
  મોટું છે, અને મૂકો:
  \sourcecorrection{OLUN-013}{સ્થિર અંગ્રેજી મૂળ
  $n$ માટે $k$ અને ઇનપુટની લંબાઈમાંથી મોટું
  મૂલ્ય લે છે. જો $n=k$ હોય, તો અટકવાના
  સમયનો પૂર્વગામી સંક્રમણ અથવા છેલ્લે લખાતું
  ખાનું સ્વ-લૂપના છેડે આવી શકે છે. $n>k$
  રાખવાથી વપરાતી સમય-સંખ્યા અને અટકવા પહેલાં
  લખાતાં ખાનાં એ છેડાથી નીચે રહે છે.}
  \[
    \Assign{\prime}{M'}(x) =
    \begin{cases}
      x + 1 &\text{જો $x < n$ હોય}\\
      n &\text{અન્યથા,}
    \end{cases}
  \]
  અને $\tuple{x,y} \in \Assign{<}{M'}$ ત્યારે
  અને માત્ર ત્યારે જ્યારે $x < y$ અથવા
  $x = y = n$. બાકીની રીતે $\Struct{M'}$ની
  વ્યાખ્યા~$\Struct{M}$ જેવી જ છે.
  $\Struct{M'}$ની વ્યાખ્યા મુજબ,
  \olref[ver]{lem:halt-if-valid}ની સાબિતીની
  જેમ, $\Sat{M'}{!T(M,w)}$. અને આપણે~$M$
  ઇનપુટ~$w$ પર અટકે છે એમ માન્યું હોવાથી
  $\Sat{M'}{!E(M,w)}$. તેથી
  $\Struct{M'}$ એ~$!T(M,w) \land !E(M,w)$નું
  સાન્ત નિદર્શ છે (અહીં~$\lif$ને બદલે~$\land$
  મૂક્યું છે).

  આ સાબિતીનો અડધો ભાગ થયો: $M$ ઇનપુટ~$w$
  પર અટકે, તો $!T(M,w) \land !E(M,w)$ને
  સાન્ત નિદર્શ છે. દુર્ભાગ્યે ઊલટી દિશા સાચી
  નથી: કેટલાંક ટ્યુરિંગ મશીનો ઇનપુટ~$w$ પર
  અટકતાં નથી, છતાં $!T(M,w) \land !E(M,w)$ને
  સાન્ત નિદર્શ મળે છે. દાખલા તરીકે, એક જ
  અવસ્થા~$q_0$ અને સૂચના
  $\delta(q_0,\TMblank) = \tuple{q_0,\TMblank,\TMstay}$
  ધરાવતું મશીન~$M$ લો. ખાલી ઇનપુટ
  $w = \emptyseq$ પર તે કદી અટકતું નથી:
  તે અનંત ચક્રમાં ચાલે છે, પરંતુ ટેપ બદલતું કે
  હેડ ખસેડતું નથી. બધાં સંરૂપણો સરખાં જ છે
  (એક જ અવસ્થા, હેડનું સ્થાન અને ટેપની સામગ્રી).
  આપણે $!T(M,\emptyseq) \land !E(M,\emptyseq)$
  સંતોષે એવું સાન્ત સંરચના~$\Struct{M''}$
  વ્યાખ્યાયિત કરી શકીએ (કસરત). તેમ છતાં,
  આવાં સંરચનાઓ દૂર થાય તે રીતે
  $!T(M,w)$માં યોગ્ય ફેરફાર કરી શકીએ.

  \begin{prob}
    એક જ અવસ્થા~$q_0$ અને એક જ સૂચના
    $\delta(q_0,\TMblank) = \tuple{q_0,\TMblank,\TMstay}$
    ધરાવતું ટ્યુરિંગ મશીન~$M$ લો.
    $\Domain{M''} = \{0, 1, 2\}$,
    $\Assign{\prime}{M''}(0) =
    \Assign{\prime}{M''}(1) = 1$ અને
    $\Assign{\prime}{M''}(2) = 2$, તથા
    $\Assign{<}{M''} = \{\tuple{0,1}, \tuple{1,1}, \tuple{2,2}\}$
    લો.
    \sourcecorrection{OLUN-015}{સ્થિર અંગ્રેજી
    મૂળ ``એક જ અવસ્થા'' કહેતાં છતાં સૂચનામાં
    અપરિચિત $q$ આપે છે; અનુગામીની બીજી
    કિંમતમાં $M''$ને બદલે $M'$ છે. અહીં
    બંને જગ્યાએ આ જ મશીન અને નિદર્શનાં
    યોગ્ય ચિહ્નો રાખ્યાં છે.}
    $\Assign{\Obj Q_{q_0}}{M''}$,
    $\Assign{\Obj S_{\TMblank}}{M''}$ અને
    $\Assign{\Obj S_{\TMendtape}}{M''}$
    એ રીતે વ્યાખ્યાયિત કરો કે
    $!T(M,\emptyseq)$ તથા $!E(M,\emptyseq)$
    સાચાં થાય, અને કેમ થાય છે તે સમજાવો.
    સંકેત: $\delta(q_0, \TMendtape)$
    અવ્યાખ્યાયિત છે એ નોંધો. ખાતરી કરો કે
    \begin{align*}
    & \Obj Q_{q_0}(\num{1}, \num{n}) \land \Obj S_{\TMendtape}(\num{0}, \num{n})
      \land
      \lforall[x][(\num{0} < x \lif \Obj S_\TMblank(x, \num{n}))] \text{\quad બધાં $n \in \Nat$ માટે}\\
    & \lexists[y][(\Obj Q_{q_0}(\num{0}, y) \land \Obj S_{\TMendtape}(\num{0}, y))]
    \end{align*}
    બંને~$\Struct{M''}$માં સત્ય છે.
  \end{prob}
\end{explain}

ઇનપુટ $w = \sigma_{i_1}\dots\sigma_{i_k}$ પર
ટ્યુરિંગ મશીન~$M$નું કામ વર્ણવતાં વાક્યો
વિચારો:
\begin{enumerate}
\item સંખ્યાઓ અને~$<$ વર્ણવતાં સ્વયંસિદ્ધ વાક્યો
  (\olref[rep]{sec}માં $!T(M,w)$ની વ્યાખ્યા જેવા જ).
\item પ્રારંભિક સંરૂપણ વર્ણવતાં સ્વયંસિદ્ધ વાક્યો:
  $!T(M,w)$ની વ્યાખ્યા જેવા જ.
\item એક સંરૂપણમાંથી બીજા સંરૂપણમાં થતું સંક્રમણ
  વર્ણવતાં સ્વયંસિદ્ધ વાક્યો:

આગળ માટે $!A(x, y)$ અગાઉ જેવું જ માનો અને
નીચેનું વ્યાખ્યાયિત કરો:
\[
  !B(y) \ident \lforall[x][(x < y \lif \eq/[x][y])].
\]
\begin{enumerate}
\item \ollabel{rep-right} દરેક સૂચના
  $\delta(q_i, \sigma) =
  \tuple{q_j, \sigma', \TMright}$ માટે વાક્ય:
\begin{align*}
& \lforall[x][\lforall[y][(
   (\Obj Q_{q_i}(x, y) \land \Obj S_{\sigma}(x, y)) \lif {}]] \\
&\qquad   (\Obj Q_{q_j}(x', y') \land \Obj S_{\sigma'}(x, y') \land
!A(x, y) \land !B(y')))
\end{align*}
\item \ollabel{rep-left} દરેક સૂચના
  $\delta(q_i, \sigma) =
  \tuple{q_j, \sigma', \TMleft}$ માટે વાક્ય:
\begin{align*}
& \lforall[x][\lforall[y][
    ((\Obj Q_{q_i}(x', y) \land \Obj S_{\sigma}(x', y)) \lif {}]]\\
& \qquad   (\Obj Q_{q_j}(x, y') \land \Obj S_{\sigma'}(x', y') \land
!A(x', y) \land !B(y'))) \land {}\\
& \lforall[y][((\Obj Q_{q_i}(\num{0}, y) \land \Obj S_{\sigma}(\num{0},
    y)) \lif {}]\\
& \qquad (\Obj Q_{q_j}(\num{0}, y') \land \Obj S_{\sigma'}(\num{0},
  y') \land !A(\num{0}, y) \land !B(y')))
\end{align*}
\sourcecorrection{OLUN-016}{સ્થિર અંગ્રેજી મૂળની
સામાન્ય ડાબે-ચાલ શાખામાં લખાતું ખાનું~$x'$
છતાં $!A(x,y)$ છે અને આગળના સમય માટે
$!B(y')$ નથી. અહીં લખાતું ખાનું
યથાવત્-શરતમાંથી બાકાત છે અને આ શાખામાં પણ
નવો સમય અગાઉના સમયોથી જુદો રહે છે.}
\item \ollabel{rep-stay} દરેક સૂચના
  $\delta(q_i, \sigma) =
  \tuple{q_j, \sigma', \TMstay}$ માટે વાક્ય:
\begin{align*}
& \lforall[x][\lforall[y][(
   (\Obj Q_{q_i}(x, y) \land \Obj S_{\sigma}(x, y)) \lif {}]] \\
&\qquad (\Obj Q_{q_j}(x, y') \land \Obj S_{\sigma'}(x, y') \land
!A(x, y) \land !B(y')))
\end{align*}
\end{enumerate}
જેમ જોઈ શકાય છે,~$M$નાં સંક્રમણો વર્ણવતાં
વાક્યો, અંતે $!B(y')$ ઉમેર્યું છે તે સિવાય,
$!T(M,w)$નાં અનુરૂપ વાક્ય જેવા જ છે.
$!B(y')$ ખાતરી કરે છે કે ``આગળના'' સંરૂપણનો
સૂચક~$y'$ અગાઉની બધી સંખ્યાઓ
$\Obj 0$, $\Obj 0'$, \dots{}થી જુદો છે.
\end{enumerate}
ટ્યુરિંગ મશીન~$M$ અને ઇનપુટ~$w$ માટે ઉપરનાં
બધાં વાક્યોનું સંયોજન~$!T'(M, w)$ માનો.

\begin{lem}\ollabel{lem:halts-sat}
  ઇનપુટ~$w$ પર શરૂ થયેલું~$M$ અટકે, તો
  $!T'(M,w) \land !E(M,w)$ને સાન્ત નિદર્શ છે.
\end{lem}

\begin{proof}
  $\Struct{M'}$ને
  \olref[ver]{lem:halt-if-valid}ની સાબિતી જેવું
  જ લો, પરંતુ:
  \begin{align*}
    \Domain{M'} & = \{0, \dots, n\},\\
    \Assign{\prime}{M'}(x) & =
    \begin{cases}
      x + 1 &\text{જો $x < n$ હોય}\\
      n &\text{અન્યથા,}
    \end{cases} \\
    \tuple{x,y} \in \Assign{<}{M'} &\text{ત્યારે અને માત્ર ત્યારે જ્યારે $x < y$ અથવા $x = y = n$,}
  \end{align*}
  જ્યાં $n = \max(k+1,\len{w})$ અને~$k$ એ
  એવી સૌથી નાની પગલાંની સંખ્યા છે કે
  ઇનપુટ~$w$ પર શરૂ થયેલું~$M$, $k$ પગલાં
  પછી અટક્યું હોય.
  \sourcecorrection{OLUN-014}{સ્થિર અંગ્રેજી
  મૂળ અહીં ફરી $n=\max(k,\len{w})$ લે છે.
  OLUN-013 મુજબ અટકવાના સમય અને છેલ્લા
  લખાતા ખાનાને છેડાની સ્વ-લૂપથી દૂર રાખવા
  અહીં $k+1$ લીધું છે; એટલે $!B(\num{k})$
  પણ અંતિમ સંક્રમણથી ફલિત થઈ શકે છે.}
  $\Sat{M'}{!T'(M,w) \land !E(M,w)}$
  તેની ચકાસણી કસરત તરીકે છોડી છે.
  \sourcecorrection{OLUN-017}{સ્થિર અંગ્રેજી
  મૂળના આ સંતોષ-દાવામાં $E(M,w)$ આગળનું
  મેટાસંકેત~$!$ છૂટ્યું છે; અહીં વ્યાખ્યાયિત
  વાક્ય~$!E(M,w)$ જ રાખ્યું છે.}
\end{proof}

\begin{prob}
  $\Sat{M'}{!T'(M,w) \land !E(M,w)}$
  સાબિત કરીને \olref[tur][und][tra]{lem:halts-sat}ની
  સાબિતી પૂર્ણ કરો.
  \sourcecorrection{OLUN-018}{સ્થિર અંગ્રેજી
  મૂળની આ કસરતમાં $!T'(M,w)$ને બદલે
  $!T(M,w)$ છે અને $E(M,w)$ આગળનું $!$
  છૂટે છે; ઉપપ્રમેયમાં દર્શાવેલું જ સંયોજન
  અહીં ચકાસવાનું છે.}
\end{prob}

\begin{lem}\ollabel{lem:sat-halts}
  $!T'(M,w) \land !E(M,w)$ને સાન્ત નિદર્શ
  હોય, તો ઇનપુટ~$w$ પર શરૂ થયેલું~$M$ અટકે છે.
\end{lem}

\begin{proof}
  પ્રતિનિષેધ સાબિત કરીએ. માનો કે ઇનપુટ~$w$
  પર શરૂ થયેલું~$M$ અટકતું નથી.
  $!T'(M,w) \land !E(M,w)$ને કોઈ નિદર્શ
  જ ન હોય, તો વાત પૂરી. તેથી માનો કે
  $\Struct{M}$ એ~$!T'(M,w) \land !E(M, w)$નું
  નિદર્શ છે. હવે બતાવવું છે કે તે સાન્ત હોઈ
  શકતું નથી.
  \sourcecorrection{OLUN-019}{સ્થિર અંગ્રેજી
  મૂળ અહીં જૂના~$!T(M,w)$નું નિદર્શ માને છે,
  પણ નવી~$!B$ શરત પરથી અનંતપણું બતાવવા
  $!T'(M,w)$નું નિદર્શ જરૂરી છે.}

  \olref[ver]{lem:config}ની જેમ સાબિત કરી
  શકીએ કે ઇનપુટ~$w$ પર શરૂ થયેલું~$M$ જો
  $n$ પગલાં પહેલાં અટક્યું ન હોય અને
  $n \ge 1$ હોય, તો $!T'(M,w) \Entails
  !C(M, w, n) \land !B(\num{n})$.
  $M$ ઇનપુટ~$w$ પર અટકતું નથી, તેથી
  દરેક ધન~$n \in \Nat$ માટે
  $!T'(M,w) \Entails !C(M, w, n) \land
  !B(\num{n})$. પ્રારંભિક સંરૂપણનું
  $!C(M,w,0)$ ઇનપુટનાં સ્વયંસિદ્ધ વાક્યોમાંથી
  અલગથી મળે છે.
  \sourcecorrection{OLUN-020}{સ્થિર અંગ્રેજી
  મૂળ $!B(\num{n})$ને $n=0$ માટે પણ
  ફલિત કહે છે; પ્રારંભિક સ્વયંસિદ્ધ વાક્યો
  એ દાવો આપતાં નથી. પ્રથમ સંક્રમણથી
  આગળના બધા $n\ge1$ માટે~$!B$ મળે છે,
  અને જુદા આંકિક પદો બતાવવા એટલું પૂરતું છે.}
  \olref[rep]{prop:mlessk}થી, દરેક~$k < n$
  માટે $!T'(M,w) \Entails \num{k} < \num{n}$.
  વળી $!B(\num{n}) \Entails
  \num{k} < \num{n} \lif
  \eq/[\num{k}][\num{n}]$. તેથી બધા
  $k < n$ માટે
  $\Sat{M}{\eq/[\num{k}][\num{n}]}$:
  એટલે અસંખ્ય પદો~$\num{k}$નાં મૂલ્યો
  $\Struct{M}$માં જુદાં જ હોવાં જોઈએ.
  પરંતુ તેના માટે $\Domain{M}$ અનંત હોવો
  જરૂરી છે. તેથી~$\Struct{M}$,
  $!T'(M,w) \land !E(M, w)$નું સાન્ત નિદર્શ
  હોઈ શકતું નથી.
\end{proof}

\begin{prob}
  ઇનપુટ~$w$ પર શરૂ થયેલું~$M$ જો~$n$
  પગલાં પહેલાં અટક્યું ન હોય અને~$n \ge 1$,
  તો $!T'(M,w) \Entails !B(\num{n})$
  સાબિત કરીને
  \olref[tur][und][tra]{lem:sat-halts}ની સાબિતી
  પૂર્ણ કરો.
  \sourcecorrection{OLUN-021}{સ્થિર અંગ્રેજી
  મૂળ આ કસરતમાં $n=0$ પણ આવરી લે છે.
  OLUN-020 મુજબ તે પ્રારંભિક કિસ્સામાં
  આ દાવો જરૂરી નથી કે ફલિત પણ થતો નથી.}
\end{prob}

\begin{thm}[ટ્રાખ્ટેનબ્રોટનું પ્રમેય]
  \ollabel{thm:trakhtenbrodt}
  પ્રથમ-ક્રમના તર્કનું મનસ્વી વાક્ય સાન્ત
  નિદર્શ ધરાવે છે કે નહીં (એટલે સાન્ત
  નિદર્શમાં સંતોષ્ય છે કે નહીં) તે અનિર્ણેય છે.
\end{thm}

\begin{proof}
  માનો કે સાન્ત નિદર્શમાં સંતોષ્યતા નક્કી કરતું
  ટ્યુરિંગ મશીન~$F$ હોય. તો કોઈ પણ ટ્યુરિંગ
  મશીન~$M$ અને ઇનપુટ~$w$ માટે વાક્ય
  $!T'(M,w) \land !E(M,w)$ ગણીને~$F$ વડે
  તેને સાન્ત નિદર્શ છે કે નહીં તે નક્કી કરી
  શકીએ. \cref{tur:und:tra:lem:halts-sat,tur:und:tra:lem:sat-halts}
  મુજબ એવું નિદર્શ ત્યારે અને માત્ર ત્યારે છે
  જ્યારે ઇનપુટ~$w$ પર શરૂ થયેલું~$M$ અટકે.
  એટલે~$F$ વડે અટકવાની સમસ્યા ઉકેલી શકાય;
  પરંતુ તે ઉકેલી શકાતી નથી એ જાણીએ છીએ.
\end{proof}

\begin{cor}\ollabel{cor:fproof-incomp}%
  \emph{સાન્ત} નિદર્શોમાં પ્રામાણ્ય માટે
  યથાર્થ અને પૂર્ણ એવી કોઈ નિષ્પત્તિ
  પદ્ધતિ હોઈ શકતી નથી: એટલે એવી નિષ્પત્તિ પદ્ધતિ,
  જેમાં દરેક સાન્ત સંરચના~$\Struct{M}$
  માટે $\Sat{M}{!B}$ હોય ત્યારે અને માત્ર
  ત્યારે $\Proves !B$ હોય.
\end{cor}

\begin{proof}
  કસરત.
\end{proof}

\begin{prob}
  \olref[tur][und][tra]{cor:fproof-incomp}
  સાબિત કરો. $!B$ દરેક સાન્ત સંરચનામાં
  સંતોષાય ત્યારે અને માત્ર ત્યારે~$\lnot !B$
  સાન્ત નિદર્શમાં સંતોષ્ય નથી એ નોંધો.
  \olref[tur][und][uns]{thm:valid-ce}ના અર્થમાં
  સાન્ત નિદર્શમાં સંતોષ્યતા અર્ધનિર્ણેય કેમ
  છે તે સમજાવો. આથી સાન્ત નિદર્શોમાં
  પ્રામાણ્ય માટેની નિષ્પત્તિ પદ્ધતિ હોત,
  તો સાન્ત નિદર્શમાં સંતોષ્યતા નિર્ણેય બનત
  એવું કેમ થાય તે દલીલ કરો.
\end{prob}
% END OLP-0273
\endgroup


\OLEndChapterHook
% END OLP-0264
\endgroup


\OLEndPartHook
% END OLP-0252
